\hypertarget{symbology-reference}{%
\section{Symbology Reference}\label{symbology-reference}}

\hypertarget{point-of-sale}{%
\subsection{Point of Sale}\label{point-of-sale}}

\hypertarget{ean-13}{%
\subsubsection{EAN-13}\label{ean-13}}

\textbf{EAN-13} is an extension of the \protect\hyperlink{upc-a}{UPC-A}
barcode symbology that usually carries a GTIN-13. It was designed by the
International Article Numbering Association in 1976 for identification
of retail goods at point of sale outside of the US.

Also known as: EAN, UCC-13, European Article Number, International
Article Number, JAN, JAN-13, IAN, WPC, SAAN, UCCET, ABAC, BCCI, ICA,
MANA, KANC, ANA, ANC.

Variants:

\begin{itemize}
\tightlist
\item
  EAN-13+2 is an extension of EAN-13 that includes a
  \protect\hyperlink{ean-2}{two-digit add-on}.
\item
  EAN-13+5 is an extension of EAN-13 that includes a
  \protect\hyperlink{ean-5}{five-digit add-on}.
\item
  EAN-99 is a special form of EAN-13 starting with \texttt{99} that is
  used as an in-store coupon.
\item
  \protect\hyperlink{ean-8}{EAN-8} is a barcode symbology derived from
  EAN-13 that is designed for small packaging. It uses a distinct
  numbering system based on GTIN-8.
\item
  \protect\hyperlink{isbn}{ISBN} is a variant of EAN-13 used to identify
  books.
\item
  \protect\hyperlink{ismn}{ISMN} is a variant of EAN-13 used to identify
  printed music.
\item
  \protect\hyperlink{issn}{ISSN} is a variant of EAN-13 used to identify
  periodicals.
\item
  \protect\hyperlink{gs1-composite-symbols}{EAN-13 Composite} is a
  variant of EAN-13 that should be used when a CC-A or CC-B GS1
  Composite 2D component is required.
\end{itemize}

Standards: ISO/IEC 15420, BS EN 797, GS1 General Specifications.

\hypertarget{data-and-options}{%
\paragraph{Data and Options}\label{data-and-options}}

\begin{itemize}
\tightlist
\item
  The data field for a EAN-13 may contain twelve or thirteen digits,
  optionally followed by a space then two or five digits if an
  \protect\hyperlink{ean-2}{EAN-2} or \protect\hyperlink{ean-5}{EAN-5}
  add-on is required.
\item
  If twelve digits of primary data are supplied then the check digit is
  calculated automatically. Otherwise the provided check digit will be
  verified.
\item
  The \textbf{includetext} option should normally be supplied.
\item
  The \textbf{guardwhitespace} option enables the display of whitespace
  guard marks.
\end{itemize}

\hypertarget{examples}{%
\paragraph{Examples}\label{examples}}

Identical symbols, input provided with and without a check digit:

\begin{verbatim}
Data:    9771473968012
Options: includetext guardwhitespace
Encoder: ean13
\end{verbatim}

\begin{verbatim}
Data:    977147396801
Options: includetext guardwhitespace
Encoder: ean13
\end{verbatim}

\includegraphics{images/ean13-1.eps}

A symbol that includes a \protect\hyperlink{ean-5}{five-digit add-on}:

\begin{verbatim}
Data:    9771473968012 54499
Options: includetext guardwhitespace
Encoder: ean13
\end{verbatim}

\includegraphics{images/ean13-2.eps}

\hypertarget{ean-8}{%
\subsubsection{EAN-8}\label{ean-8}}

\textbf{EAN-8} is derived from the \protect\hyperlink{ean-13}{EAN-13}
barcode symbology and is designed for small packaging. It usually
carries a GTIN-8.

Also known as: UCC-8, JAN-8.

Variants:

\begin{itemize}
\tightlist
\item
  EAN-8+2 is an extension of EAN-8 that includes a
  \protect\hyperlink{ean-2}{two-digit add-on}.
\item
  EAN-8+5 is an extension of EAN-8 that includes a
  \protect\hyperlink{ean-5}{five-digit add-on}.
\item
  EAN-Velocity is a special form of EAN-8 starting with \texttt{0} that
  is used for in-store coupons.
\item
  \protect\hyperlink{ean-13}{EAN-13} is a longer variant of EAN-8 which
  has a distinct number system based on GTIN-13.
\item
  \protect\hyperlink{gs1-composite-symbols}{EAN-8 Composite} is a
  variant of EAN-8 that should be used when a CC-A or CC-B GS1 composite
  2D component is required.
\end{itemize}

Standards: ISO/IEC 15420, BS EN 797, GS1 General Specifications.

\hypertarget{data-and-options-1}{%
\paragraph{Data and Options}\label{data-and-options-1}}

\begin{itemize}
\tightlist
\item
  The data field takes either seven or eight digits, optionally followed
  by a space then two or five digits if an
  \protect\hyperlink{ean-2}{EAN-2} or \protect\hyperlink{ean-5}{EAN-5}
  add-on is required.
\item
  If seven digits of primary data are supplied then the check digit is
  calculated automatically. Otherwise the provided check digit will be
  verified.
\item
  The \textbf{includetext} option should normally be supplied.
\item
  The \textbf{guardwhitespace} option enables the display of whitespace
  guard marks.
\end{itemize}

\hypertarget{examples-1}{%
\paragraph{Examples}\label{examples-1}}

Identical symbols, input provided with and without a check digit:

\begin{verbatim}
Data:    01335583
Options: includetext
Encoder: ean8
\end{verbatim}

\begin{verbatim}
Data:    0133558
Options: includetext
Encoder: ean8
\end{verbatim}

\includegraphics{images/ean8-1.eps}

Truncated with whitespace guards:

\begin{verbatim}
Data:    01335583
Options: includetext height=0.5 guardwhitespace
Encoder: ean8
\end{verbatim}

\includegraphics{images/ean8-2.eps}

\hypertarget{upc-a}{%
\subsubsection{UPC-A}\label{upc-a}}

The \textbf{UPC-A} barcode symbology is used for identification of
retail goods at point of sale inside of the US. It usually carries a
GTIN-12.

Also known as: UPC, UCC-12, Universal Product Code.

Variants:

\begin{itemize}
\tightlist
\item
  UPC-A+2 is an extension of UPC-A that includes a
  \protect\hyperlink{ean-2}{two-digit add-on}.
\item
  UPC-A+5 is an extension of UPC-A that includes a
  \protect\hyperlink{ean-5}{five-digit add-on}.
\item
  \protect\hyperlink{upc-e}{UPC-E} is a barcode symbology derived from
  UPC-A that is designed for small packaging.
\item
  \protect\hyperlink{gs1-composite-symbols}{UPC-A Composite} is a
  variant of UPC-A that should be used when a CC-A or CC-B
  \protect\hyperlink{gs1-composite-symbols}{GS1 composite} 2D component
  is required.
\item
  A UPC-A symbol can be converted to an
  \protect\hyperlink{ean-13}{EAN-13} symbol by prefixing the GTIN-12
  with 0 to make the equivalent GTIN-13.
\end{itemize}

Standards: ISO/IEC 15420, BS EN 797, GS1 General Specifications.

\hypertarget{data-and-options-2}{%
\paragraph{Data and Options}\label{data-and-options-2}}

\begin{itemize}
\tightlist
\item
  The data field for a UPC-A may contain eleven or twelve digits,
  optionally followed by a space then two or five digits if an
  \protect\hyperlink{ean-2}{EAN-2} or \protect\hyperlink{ean-5}{EAN-5}
  add-on is required.
\item
  Alternatively, the data field may contain seven or eight digits of a
  \protect\hyperlink{upc-e}{UPC-E} to produce the equivalent UPC-A
  symbol.
\item
  If eleven digits of primary data are supplied then the check digit is
  calculated automatically. Otherwise the provided check digit will be
  verified.
\item
  The \textbf{includetext} option should normally be supplied.
\end{itemize}

\hypertarget{examples-2}{%
\paragraph{Examples}\label{examples-2}}

Identical symbols, input provided with and without a check digit:

\begin{verbatim}
Data:    788581014974
Options: includetext
Encoder: upca
\end{verbatim}

\begin{verbatim}
Data:    78858101497
Options: includetext
Encoder: upca
\end{verbatim}

\includegraphics{images/upca-1.eps}

A symbol that includes a \protect\hyperlink{ean-5}{five-digit add-on}:

\begin{verbatim}
Data:    788581014974 54499
Options: includetext guardwhitespace
Encoder: upca
\end{verbatim}

\includegraphics{images/upca-2.eps}

\hypertarget{upc-e}{%
\subsubsection{UPC-E}\label{upc-e}}

\textbf{UPC-E} is a compacted form of the
\protect\hyperlink{upc-a}{UPC-A} barcode symbology that usually carries
a GTIN-12 with a number system of \emph{0} or \emph{1} that has been
zero compressed.

Variants:

\begin{itemize}
\tightlist
\item
  UPC-E0 is a UPC-E with a number system of \emph{0}.
\item
  UPC-E1 is a UPC-E with a number system of \emph{1}.
\item
  UPC-E+2 is an extension of UPC-E that includes a
  \protect\hyperlink{ean-2}{two-digit add-on}.
\item
  UPC-E+5 is an extension of UPC-E that includes a
  \protect\hyperlink{ean-5}{five-digit add-on}.
\item
  \protect\hyperlink{upc-a}{UPC-A} is the full size form of UPC-E.
\item
  \protect\hyperlink{gs1-composite-symbols}{UPC-E Composite} is a
  variant of UPC-E that should be used when a CC-A or CC-B GS1 Composite
  2D component is required.
\end{itemize}

Standards: ISO/IEC 15420, BS EN 797, GS1 General Specifications.

\hypertarget{data-and-options-3}{%
\paragraph{Data and Options}\label{data-and-options-3}}

\begin{itemize}
\tightlist
\item
  The data field takes either seven or eight digits, optionally followed
  by a space then two or five digits if an
  \protect\hyperlink{ean-2}{EAN-2} or \protect\hyperlink{ean-5}{EAN-5}
  add-on is required.
\item
  Alternatively, the data field may contain eleven or twelve digits of a
  \protect\hyperlink{upc-a}{UPC-A} to produce the equivalent UPC-E
  symbol, provided that the input can be zero suppressed.
\item
  If seven digits of primary data are supplied then the check digit is
  calculated automatically. Otherwise the provided check digit will be
  verified.
\item
  The \textbf{includetext} option should normally be supplied.
\end{itemize}

\hypertarget{examples-3}{%
\paragraph{Examples}\label{examples-3}}

Identical symbols, input provided with and without a check digit:

\begin{verbatim}
Data:    01234565
Options: includetext
Encoder: upce
\end{verbatim}

\begin{verbatim}
Data:    0123456
Options: includetext
Encoder: upce
\end{verbatim}

\includegraphics{images/upce-1.eps}

A truncated symbol:

\begin{verbatim}
Data:    01234565
Options: includetext height=0.5
Encoder: upce
\end{verbatim}

\includegraphics{images/upce-2.eps}

\hypertarget{isbn}{%
\subsubsection{ISBN}\label{isbn}}

An \textbf{ISBN} barcode is a variant of
\protect\hyperlink{ean-13}{EAN-13} that is used to identify books.

Also known as: ISBN-13, International Standard Book Number, Bookland
EAN-13.

Variants:

\begin{itemize}
\tightlist
\item
  ISBN-10 is a legacy format that was depreciated for public use after
  1st January 2007.
\end{itemize}

Standards: ISO 2108, ISO/IEC 15420, BS EN 797, GS1 General
Specifications.

\hypertarget{data-and-options-4}{%
\paragraph{Data and Options}\label{data-and-options-4}}

\begin{itemize}
\tightlist
\item
  The data should contain twelve or thirteen digits separated
  appropriately by dash characters \emph{-}.
\item
  The data can also be provided in legacy ISBN-10 format as nine or ten
  digits separated appropriately by dash characters \emph{-}. This will
  be automatically upgraded to the ISBN-13 format.
\item
  If the last digit of the primary data is not given then the ISBN check
  digit is calculated automatically, otherwise it will be verified.
\item
  The primary data can optionally be followed by a space then two or
  five digits if an \protect\hyperlink{ean-2}{EAN-2} or
  \protect\hyperlink{ean-5}{EAN-5} add-on is required.
\item
  The \textbf{includetext} option should normally be supplied.
\item
  The \textbf{guardwhitespace} option enables the display of whitespace
  guard marks.
\item
  The following options are also relevant to this barcode symbology:

  \begin{itemize}
  \tightlist
  \item
    \textbf{isbntextfont}: PostScript font name for text above symbol
  \item
    \textbf{isbntextsize}: Font size for the text above symbol, in
    points
  \item
    \textbf{isbntextxoffset}: Horizontal position of ISBN text, in
    points
  \item
    \textbf{isbntextyoffset}: Vertical position of ISBN text, in points
  \end{itemize}
\item
  \emph{Deprecated.} The \textbf{legacy} option prevents ISBN-10 input
  from being upgraded to ISBN-13 and will result in a symbol that is
  obsolete and should not be used at point of sale.
\end{itemize}

\hypertarget{example-isbn}{%
\paragraph{Example ISBN}\label{example-isbn}}

Identical symbols, input provided with and without an ISBN check digit:

\begin{verbatim}
Data:    978-1-873671-00-9
Options: includetext
Encoder: isbn
\end{verbatim}

\begin{verbatim}
Data:    978-1-873671-00
Options: includetext
Encoder: isbn
\end{verbatim}

\includegraphics{images/isbn-1.eps}

An ISBN with a \protect\hyperlink{ean-5}{five-digit add-on}:

\begin{verbatim}
Data:    978-1-873671-00-9 54499
Options: includetext guardwhitespace
Encoder: isbn
\end{verbatim}

\includegraphics{images/isbn-2.eps}

The following ISBN-10 input will be automatically upgraded to a valid
ISBN-13 symbol:

\begin{verbatim}
Data:    1-86074-271-2
Options: includetext
Encoder: isbn
\end{verbatim}

\begin{verbatim}
Data:    1-86074-271
Options: includetext
Encoder: isbn
\end{verbatim}

\includegraphics{images/isbn-3.eps}

\hypertarget{example-isbn-10}{%
\paragraph{Example ISBN-10}\label{example-isbn-10}}

Note that ISBN-10 is legacy format not for use at P.O.S.

The following will generate an obsolete ISBN-10 symbol:

\begin{verbatim}
Data:    1-86074-271-8
Options: legacy includetext guardwhitespace
Encoder: isbn
\end{verbatim}

\begin{verbatim}
Data:    1-86074-271
Options: legacy includetext guardwhitespace
Encoder: isbn
\end{verbatim}

\includegraphics{images/isbn-4.eps}

\hypertarget{ismn}{%
\subsubsection{ISMN}\label{ismn}}

An \textbf{ISMN} barcode is a variant of
\protect\hyperlink{ean-13}{EAN-13} with a prefix \emph{979} that is used
to identify printed music.

Also known as: International Standard Music Number, ISMN-13.

Variants:

\begin{itemize}
\tightlist
\item
  ISMN-10 is a legacy format that was depreciated for public use.
\end{itemize}

Standards: ISO 10957, ISO/IEC 15420, BS EN 797, GS1 General
Specifications.

\hypertarget{data-and-options-5}{%
\paragraph{Data and Options}\label{data-and-options-5}}

\begin{itemize}
\tightlist
\item
  The data should contain twelve or thirteen digits separated
  appropriately by dash characters \emph{-}.
\item
  The data can also be provided in legacy ISMN-10 format start \emph{M-}
  then eight or nine digits separated appropriately by dash characters
  \emph{-}. This will be automatically upgraded to the ISMN-13 format.
\item
  If the last digit of the primary data is not given then the ISMN check
  digit is calculated automatically, otherwise it will be verified.
\item
  The primary data can optionally be followed by a space then two or
  five digits if an \protect\hyperlink{ean-2}{EAN-2} or
  \protect\hyperlink{ean-5}{EAN-5} add-on is required.
\item
  The \textbf{includetext} option should normally be supplied.
\item
  The \textbf{guardwhitespace} option enables the display of whitespace
  guard marks.
\item
  The following options are also relevant to this barcode symbology:

  \begin{itemize}
  \tightlist
  \item
    \textbf{ismntextfont}: PostScript font name for text above symbol
  \item
    \textbf{ismntextsize}: Font size for the text above symbol, in
    points
  \item
    \textbf{ismntextxoffset}: Horizontal position of ISMN text, in
    points
  \item
    \textbf{ismntextyoffset}: Vertical position of ISMN text, in points
  \end{itemize}
\item
  \emph{Deprecated.} The \textbf{legacy} option prevents ISMN-10 input
  from being upgraded to ISMN-13 and will result in a symbol that is
  obsolete and should not be used at point of sale.
\end{itemize}

\hypertarget{example-ismn}{%
\paragraph{Example ISMN}\label{example-ismn}}

Identical symbols, input provided with and without an ISMN check digit:

\begin{verbatim}
Data:    979-0-2600-0043-8
Options: includetext
Encoder: ismn
\end{verbatim}

\begin{verbatim}
Data:    979-0-2600-0043
Options: includetext
Encoder: ismn
\end{verbatim}

\includegraphics{images/ismn-1.eps}

The following ISMN-10 input will be automatically upgraded to a valid
ISMN-13 symbol:

\begin{verbatim}
Data:    M-345-24680-5
Options: includetext
Encoder: ismn
\end{verbatim}

\begin{verbatim}
Data:    M-345-24680
Options: includetext
Encoder: ismn
\end{verbatim}

\includegraphics{images/ismn-2.eps}

\hypertarget{example-ismn-10}{%
\paragraph{Example ISMN-10}\label{example-ismn-10}}

Note that ISMN-10 is a legacy format not for use at P.O.S.

The following will generate an obsolete ISMN-10 symbol:

\begin{verbatim}
Data:    M-345-24680-5
Options: legacy includetext guardwhitespace
Encoder: ismn
\end{verbatim}

\begin{verbatim}
Data:    M-345-24680
Options: legacy includetext guardwhitespace
Encoder: ismn
\end{verbatim}

\includegraphics{images/ismn-3.eps}

\hypertarget{issn}{%
\subsubsection{ISSN}\label{issn}}

An \textbf{ISSN} barcode is an \protect\hyperlink{ean-13}{EAN-13} with
prefix \emph{977} used to identify periodicals.

Also known as: International Standard Serial Number.

Standards: ISO 3297, ISO/IEC 15420, BS EN 797, GS1 General
Specifications.

\hypertarget{data-and-options-6}{%
\paragraph{Data and Options}\label{data-and-options-6}}

\begin{itemize}
\tightlist
\item
  The data should contain the seven or eight digits ISSN separated by a
  dash characters \emph{-}, followed by a two-digit sequence variant,
  optionally followed by two or five digits if a
  \protect\hyperlink{ean-2}{two-digit add-on} or
  \protect\hyperlink{ean-5}{five-digit add-on} is required.
\item
  If the last digit of the ISSN data is not given then the ISSN check
  digit is calculated automatically, otherwise it will be verified.
\item
  The \textbf{includetext} option should normally be supplied.
\item
  The \textbf{guardwhitespace} option enables the display of whitespace
  guard marks.
\item
  The following options are also relevant to this barcode symbology:

  \begin{itemize}
  \tightlist
  \item
    \textbf{issntextfont}: PostScript font name for text above symbol
  \item
    \textbf{issntextsize}: Font size for the text above symbol, in
    points
  \item
    \textbf{issntextxoffset}: Horizontal position of ISSN text, in
    points
  \item
    \textbf{issntextyoffset}: Vertical position of ISSN text, in points
  \end{itemize}
\end{itemize}

A sequence variant is a two-digit number that usually starts at zero and
is incremented whenever the recommended retail price is amended, where
applicable.

\hypertarget{example}{%
\paragraph{Example}\label{example}}

Identical symbols, input provided with and without an ISSN check digit
and having sequence number \emph{00}:

\begin{verbatim}
Data:    0317-8471 00
Options: includetext guardwhitespace
Encoder: issn
\end{verbatim}

\begin{verbatim}
Data:    0317-847 00
Options: includetext guardwhitespace
Encoder: issn
\end{verbatim}

\includegraphics{images/issn-1.eps}

An ISSN with sequence number \emph{03} and a
\protect\hyperlink{ean-2}{two-digit add-on} representing issue number
\emph{17}:

\begin{verbatim}
Data:    0317-8471 03 17
Options: includetext guardwhitespace
Encoder: issn
\end{verbatim}

\begin{verbatim}
Data:    0317-847 03 17
Options: includetext guardwhitespace
Encoder: issn
\end{verbatim}

\includegraphics{images/issn-2.eps}

\hypertarget{two-dimensional}{%
\subsection{Two-Dimensional}\label{two-dimensional}}

\hypertarget{aztec-code}{%
\subsubsection{Aztec Code}\label{aztec-code}}

\textbf{Aztec Code} is a 2D matrix-style barcode symbology. It can
encode full 256-character extended-ASCII.

Variants:

\begin{itemize}
\tightlist
\item
  \protect\hyperlink{aztec-runes}{Aztec Runes} are a set of small
  barcode symbols that are used for special applications.
\end{itemize}

Standards: ISO/IEC 24778, ANSI/AIM BC13 - ISS Aztec Code.

\hypertarget{data-and-options-7}{%
\paragraph{Data and Options}\label{data-and-options-7}}

\begin{itemize}
\tightlist
\item
  The data field can contain any extended ASCII data. The default
  interpretation of data by readers is in accordance with ISO/IEC
  8859-1. When supported by the receiver characters from other code
  pages can be encoded using Extended Channel Interpretation (ECI).
\item
  When the \textbf{parse} option is specified, any instances of
  \texttt{\^{}NNN} in the data field are replaced with their equivalent
  ASCII value, useful for specifying unprintable characters.
\item
  When the \textbf{parsefnc} option is specified, non-data function
  characters can be specified by escape sequences:

  \begin{itemize}
  \tightlist
  \item
    \texttt{\^{}FNC1}: FNC1
  \item
    \texttt{\^{}ECI000000} to \texttt{\^{}ECI999999}: ECI indicators
  \end{itemize}
\item
  The \textbf{eclevel} option is used to specify the percentage of error
  correction to be applied when expanding the data, by default
  \texttt{23}.
\item
  The \textbf{ecaddchars} option is used to specify how many additional
  error correction characters to apply the data once expanded by the
  eclevel percentage, by default \texttt{3}.
\item
  The \textbf{layers} option is used to specify a particular number of
  layers in which to encode the data, between \texttt{1} and
  \texttt{32}. By default the encoder will create a symbol with be
  minimal number of layers to encode the given data.
\item
  The \textbf{format} option is used to select between
  \texttt{format=full} and \texttt{format=compact} symbol types. By
  default the encoder will choose the most appropriate format to create
  a symbol of minimal size.
\item
  The \textbf{readerinit} option denotes that the symbol is used for
  programming the barcode reader.
\item
  The \textbf{raw} option denotes that the data field is providing the
  input as a pre-encoded bitstream suitable for direct low-level
  encoding.
\item
  \emph{Deprecated: Use \protect\hyperlink{aztec-runes}{Aztec Runes}
  instead.} The \textbf{format} option can also be used to create Aztec
  Code "runes", using \texttt{format=rune}. In this case the rune symbol
  number should be given in the data field.
\end{itemize}

\hypertarget{examples-4}{%
\paragraph{Examples}\label{examples-4}}

\begin{verbatim}
Data:    This is Aztec Code
Options: 
Encoder: azteccode
\end{verbatim}

\includegraphics{images/aztec-1.eps}

\begin{verbatim}
Data:    This is ^065ztec Code
Options: parse eclevel=50 ecaddchars=0
Encoder: azteccode
\end{verbatim}

\includegraphics{images/aztec-6.eps}

\begin{verbatim}
Data:    ABC123
Options: layers=3 format=full
Encoder: azteccode
\end{verbatim}

\includegraphics{images/aztec-2.eps}

\begin{verbatim}
Data:    ABC123
Options: format=compact
Encoder: azteccode
\end{verbatim}

\includegraphics{images/aztec-5.eps}

\begin{verbatim}
Data:    00100111001000000101001101111000010100111100101000000110
Options: raw
Encoder: azteccode
\end{verbatim}

\includegraphics{images/aztec-4.eps}

\hypertarget{aztec-runes}{%
\subsubsection{Aztec Runes}\label{aztec-runes}}

\textbf{Aztec Runes} are a set of small barcode symbols that are used
for special applications.

Variants:

\protect\hyperlink{aztec-code}{Aztec Code} is a 2D matrix-style barcode
symbology that can encode full 256 character extended-ASCII.

Standards: ISO/IEC 24778, ANSI/AIM BC13 - ISS Aztec Code.

\hypertarget{data-and-options-8}{%
\paragraph{Data and Options}\label{data-and-options-8}}

\begin{itemize}
\tightlist
\item
  The data field contains the rune number \texttt{0} to \texttt{255}.
\end{itemize}

\hypertarget{examples-5}{%
\paragraph{Examples}\label{examples-5}}

\begin{verbatim}
Data:    25
Options: 
Encoder: aztecrune
\end{verbatim}

\includegraphics{images/aztec-3.eps}

\hypertarget{data-matrix}{%
\subsubsection{Data Matrix}\label{data-matrix}}

The \textbf{Data Matrix} symbology is 2D matrix-style barcode that can
encode full 256 character extended-ASCII.

Also known as: Data Matrix ECC 200, DataMatrix.

Variants:

\begin{itemize}
\tightlist
\item
  \protect\hyperlink{data-matrix-rectangular}{Data Matrix Rectangular}
  is a variant of Data Matrix that uses rectangular symbols.
\item
  \protect\hyperlink{data-matrix-rectangular-extension}{Data Matrix
  Rectangular Extension} is a variant of Data Matrix Rectangular that
  provides a greater range of symbol dimensions.
\item
  \protect\hyperlink{gs1-datamatrix}{GS1 DataMatrix} is a variant of
  Data Matrix that should be used when encoding data that is in
  \protect\hyperlink{gs1-application-identifier-standard-format}{GS1
  Application Identifier standard format}.
\item
  \protect\hyperlink{gs1-digital-link-data-matrix}{GS1 Digital Link Data
  Matrix} is a variant of Data Matrix that should be used when encoding
  a GS1 Digital Link URI.
\item
  \protect\hyperlink{hibc-symbols}{HIBC Data Matrix} is a variant of
  Data Matrix that should be used when encoding HIBC formatted data.
\item
  \protect\hyperlink{mailmark}{Royal Mail Mailmark} is a variant of Data
  Matrix that is used by the Royal Mail postal service on letters and
  other small mailpieces.
\end{itemize}

Standards: ISO/IEC 16022, ANSI/AIM BC11 - ISS Data Matrix.

\hypertarget{data-and-options-9}{%
\paragraph{Data and Options}\label{data-and-options-9}}

\begin{itemize}
\tightlist
\item
  The data field can contain any extended ASCII data. The default
  interpretation of data by readers is in accordance with ISO/IEC
  8859-1. When supported by the receiver characters from other code
  pages can be encoded using Extended Channel Interpretation (ECI).
\item
  When the \textbf{parse} option is specified, any instances of
  \texttt{\^{}NNN} in the data field are replaced with their equivalent
  ASCII value. This is useful for specifying unprintable characters.
\item
  When the \textbf{parsefnc} option is specified, non-data function
  characters can be specified by escape sequences:

  \begin{itemize}
  \tightlist
  \item
    \texttt{\^{}FNC1}: FNC1. \emph{Recommendation: Use
    \protect\hyperlink{gs1-datamatrix}{GS1 DataMatrix} for encoding GS1
    AI syntax data.}
  \item
    \texttt{\^{}PROG}: PROG - Reader programming
  \item
    \texttt{\^{}ECI000000} to \texttt{\^{}ECI999999}: ECI indicators
  \end{itemize}
\item
  Whenever the data field contains suitable input, the encoder will
  compress ISO/IEC 15434 structured data (internally using the 05 Macro
  and 06 Macro codewords.)
\item
  When the \textbf{dotty} option is specified the modules are rendered
  as dots rather than squares. The dot radius can be adjusted using the
  \textbf{inkspread} option.
\item
  The \textbf{columns} and \textbf{rows} options are used to specify the
  size of the symbol.
\item
  The \textbf{version} option can also be used to specify the symbol
  size, as \texttt{version=RxC}. Valid options are \emph{10x10},
  \emph{12x12}, \emph{14x14}, \emph{16x16}, \emph{18x18}, \emph{20x20},
  \emph{22x22}, \emph{24x24}, \emph{26x26}, \emph{32x32}, \emph{36x36},
  \emph{40x40}, \emph{44x44}, \emph{48x48}, \emph{52x52}, \emph{64x64},
  \emph{72x72}, \emph{80x80}, \emph{88x88}, \emph{96x96},
  \emph{104x104}, \emph{120x120}, \emph{132x132}, \emph{144x144}
\item
  If \textbf{columns}, \textbf{rows} and \textbf{version} are
  unspecified the encoder will default to creating a symbol of the
  specified \textbf{format} that is the minimum size to represent the
  given data.
\item
  The \textbf{raw} option denotes that the data field is providing the
  input as pre-encoded codewords in \texttt{\^{}NNN} format, suitable
  for direct low-level encoding.
\item
  \emph{Deprecated}: The \textbf{format} option is used to specify the
  shape of the symbol, either:

  \begin{itemize}
  \tightlist
  \item
    \texttt{square} (default)
  \item
    \texttt{rectangle}: Use
    \protect\hyperlink{data-matrix-rectangular}{Data Matrix Rectangular}
    for rectangular symbols.
  \end{itemize}
\item
  \emph{Deprecated: Use
  \protect\hyperlink{data-matrix-rectangular-extension}{Data Matrix
  Rectangular Extension}.} The \textbf{dmre} option enables Data Matrix
  Rectangular Extension symbols based on the ISO/IEC standard, which
  increases the number of rectangular symbol sizes available.
\item
  \emph{Deprecated: Use
  \protect\hyperlink{data-matrix-rectangular-extension}{Data Matrix
  Rectangular Extension} with dimdmre option.} The \textbf{dindmre}
  option enables Data Matrix Rectangular Extension symbols based on the
  obsolete DIN standard.
\item
  \emph{Deprecated: For internal use.} The \textbf{mailmark} option
  signifies the requirement to use Mailmark-specific encoding.
\end{itemize}

\hypertarget{examples-6}{%
\paragraph{Examples}\label{examples-6}}

Identical symbols, automatically sized, the latter showing use of the
parse option:

\begin{verbatim}
Data:    This is Data Matrix
Options: 
Encoder: datamatrix
\end{verbatim}

\begin{verbatim}
Data:    This is ^068ata Matrix
Options: parse
Encoder: datamatrix
\end{verbatim}

\includegraphics{images/datamatrix-1.eps}

Identical symbols with a fixed size:

\begin{verbatim}
Data:    Fixed size
Options: version=48x48
Encoder: datamatrix
\end{verbatim}

\begin{verbatim}
Data:    Fixed size
Options: rows=48 columns=48
Encoder: datamatrix
\end{verbatim}

\includegraphics{images/datamatrix-2.eps}

A rectangular symbol with a fixed size:

\begin{verbatim}
Data:    Rectangular
Options: format=rectangle version=16x48
Encoder: datamatrix
\end{verbatim}

\includegraphics{images/datamatrix-3.eps}

An advanced use containing "hidden data" in the padding codewords as
used by some non-standard, "high-security" applications. The technique
works by filling the symbol using raw codewords formed from the standard
data immediately followed by the non-standard padding data beginning
with \texttt{\^{}129}:

\begin{verbatim}
Data:    ^066^067^068^142^052^129^161^056^206^101^251^147
Options: version=16x16 raw
Encoder: datamatrix
\end{verbatim}

\hypertarget{data-matrix-rectangular}{%
\subsubsection{Data Matrix Rectangular}\label{data-matrix-rectangular}}

The \textbf{Data Matrix Rectangular} symbology is a rectangular variant
of \protect\hyperlink{data-matrix}{Data Matrix}.

Variants:

\begin{itemize}
\tightlist
\item
  \protect\hyperlink{data-matrix}{Data Matrix} is the standard square
  variant of the symbology.
\item
  \protect\hyperlink{data-matrix-rectangular-extension}{Data Matrix
  Rectangular Extension} is a variant of Data Matrix Rectangular that
  provides a greater range of symbol dimensions.
\end{itemize}

Standards: ISO/IEC 16022, ANSI/AIM BC11 - ISS Data Matrix.

\hypertarget{data-and-options-10}{%
\paragraph{Data and Options}\label{data-and-options-10}}

\begin{itemize}
\tightlist
\item
  The data field can contain any extended ASCII data. The default
  interpretation of data by readers is in accordance with ISO/IEC
  8859-1. When supported by the receiver characters from other code
  pages can be encoded using Extended Channel Interpretation (ECI).
\item
  When the \textbf{parse} option is specified, any instances of
  \texttt{\^{}NNN} in the data field are replaced with their equivalent
  ASCII value. This is useful for specifying unprintable characters.
\item
  When the \textbf{parsefnc} option is specified, non-data function
  characters can be specified by escape sequences:

  \begin{itemize}
  \tightlist
  \item
    \texttt{\^{}FNC1}: FNC1. \emph{Recommendation: Use
    \protect\hyperlink{gs1-datamatrix}{GS1 DataMatrix} for encoding GS1
    AI syntax data.}
  \item
    \texttt{\^{}PROG}: PROG - Reader programming
  \item
    \texttt{\^{}ECI000000} to \texttt{\^{}ECI999999}: ECI indicators
  \end{itemize}
\item
  Whenever the data field contains suitable input, the encoder will
  compress ISO/IEC 15434 structured data (internally using the 05 Macro
  and 06 Macro codewords.)
\item
  When the \textbf{dotty} option is specified the modules are rendered
  as dots rather than squares. The dot radius can be adjusted using the
  \textbf{inkspread} option.
\item
  \emph{Deprecated: Use
  \protect\hyperlink{data-matrix-rectangular-extension}{Data Matrix
  Rectangular Extension}.} The \textbf{dmre} option enables Data Matrix
  Rectangular Extension symbols based on the ISO/IEC standard, which
  increases the number of rectangular symbol sizes available.
\item
  \emph{Deprecated: Standard is obsolete. Use
  \protect\hyperlink{data-matrix-rectangular-extension}{Data Matrix
  Rectangular Extension} with the dindmre option.} The \textbf{dindmre}
  option enables Data Matrix Rectangular Extension symbols based on the
  obsolete DIN standard.
\item
  The \textbf{columns} and \textbf{rows} options are used to specify the
  size of the symbol.
\item
  The \textbf{version} option can also be used to specify the symbol
  size, as \texttt{version=RxC}. Valid options are \emph{8x18},
  \emph{8x32}, \emph{12x26}, \emph{12x36}, \emph{16x36}, \emph{16x48}.
\item
  If \textbf{columns}, \textbf{rows} and \textbf{version} are
  unspecified the encoder will default to creating a symbol that is the
  minimum size to represent the given data.
\item
  The \textbf{raw} option denotes that the data field is providing the
  input as pre-encoded codewords in \texttt{\^{}NNN} format, suitable
  for direct low-level encoding.
\end{itemize}

\hypertarget{examples-7}{%
\paragraph{Examples}\label{examples-7}}

A rectangular symbol with a fixed size:

\begin{verbatim}
Data:    Rectangular
Options: version=16x48
Encoder: datamatrixrectangular
\end{verbatim}

\includegraphics{images/datamatrix-3.eps}

\hypertarget{data-matrix-rectangular-extension}{%
\subsubsection{Data Matrix Rectangular
Extension}\label{data-matrix-rectangular-extension}}

The \textbf{Data Matrix Rectangular Extension} symbology is an
enhancement to \protect\hyperlink{data-matrix-rectangular}{Data Matrix
Rectangular} that provides a greater range of symbol dimensions.

Also known as: DMRE.

Variants:

\begin{itemize}
\tightlist
\item
  \protect\hyperlink{data-matrix}{Data Matrix} is the standard square
  variant of the symbology.
\item
  \protect\hyperlink{data-matrix-rectangular}{Data Matrix Rectangular}
  is the original rectangular variant of Data Matrix that provides a
  limited range of symbol dimensions.
\end{itemize}

Standards: ISO/IEC 21471.

\hypertarget{data-and-options-11}{%
\paragraph{Data and Options}\label{data-and-options-11}}

\begin{itemize}
\tightlist
\item
  The data field can contain any extended ASCII data. The default
  interpretation of data by readers is in accordance with ISO/IEC
  8859-1. When supported by the receiver characters from other code
  pages can be encoded using Extended Channel Interpretation (ECI).
\item
  When the \textbf{parse} option is specified, any instances of
  \texttt{\^{}NNN} in the data field are replaced with their equivalent
  ASCII value. This is useful for specifying unprintable characters.
\item
  When the \textbf{parsefnc} option is specified, non-data function
  characters can be specified by escape sequences:

  \begin{itemize}
  \tightlist
  \item
    \texttt{\^{}FNC1}: FNC1
  \item
    \texttt{\^{}PROG}: PROG - Reader programming
  \item
    \texttt{\^{}ECI000000} to \texttt{\^{}ECI999999}: ECI indicators
  \end{itemize}
\item
  Whenever the data field contains suitable input, the encoder will
  compress ISO/IEC 15434 structured data (internally using the 05 Macro
  and 06 Macro codewords.)
\item
  When the \textbf{dotty} option is specified the modules are rendered
  as dots rather than squares. The dot radius can be adjusted using the
  \textbf{inkspread} option.
\item
  The \textbf{columns} and \textbf{rows} options are used to specify the
  size of the symbol.
\item
  The \textbf{version} option can also be used to specify the symbol
  size, as \texttt{version=RxC}. Valid options are \emph{8x18},
  \emph{8x32}, \emph{8x48}, \emph{8x64}, \emph{8x80}, \emph{8x96},
  \emph{8x120}, \emph{8x144}, \emph{12x26}, \emph{12x36}, \emph{12x64},
  \emph{12x88}, \emph{16x36}, \emph{16x48}, \emph{16x64}, \emph{20x36},
  \emph{20x44}, \emph{20x64}, \emph{22x48}, \emph{24x48}, \emph{24x64},
  \emph{26x40}, \emph{26x48}, \emph{26x64}.
\item
  If \textbf{columns}, \textbf{rows} and \textbf{version} are
  unspecified the encoder will default to creating a symbol that is the
  minimum size to represent the given data.
\item
  The \textbf{raw} option denotes that the data field is providing the
  input as pre-encoded codewords in \texttt{\^{}NNN} format, suitable
  for direct low-level encoding.
\item
  \emph{Deprecated: Standard is obsolete.} The \textbf{dindmre} option
  select Data Matrix Rectangular Extension symbols whose dimensions are
  based on the obsolete DIN standard, rather than the current ISO
  standard.

  \begin{itemize}
  \tightlist
  \item
    With \texttt{dindmre} the valid version \textbf{options} are
    \emph{8x18}, \emph{8x32}, \emph{8x48}, \emph{8x64}, \emph{12x26},
    \emph{12x36}, \emph{12x64}, \emph{16x36}, \emph{16x48},
    \emph{16x64}, \emph{24x32}, \emph{24x36}, \emph{24x48},
    \emph{24x64}, \emph{26x32}, \emph{26x40}, \emph{26x48},
    \emph{26x64}.
  \end{itemize}
\end{itemize}

\hypertarget{examples-8}{%
\paragraph{Examples}\label{examples-8}}

A DMRE symbol with a fixed size:

\begin{verbatim}
Data:    1234
Options: version=8x80
Encoder: datamatrixrectangularextension
\end{verbatim}

\includegraphics{images/datamatrix-4.eps}

\hypertarget{han-xin-code}{%
\subsubsection{Han Xin Code}\label{han-xin-code}}

The \textbf{Han Xin Code} symbology is a 2D matrix-style barcode
symbology that can encode full 256 character extended-ASCII.

Also known as: Chinese Sensible.

Standards: GB/T 21049-2007.

\hypertarget{data-and-options-12}{%
\paragraph{Data and Options}\label{data-and-options-12}}

\begin{itemize}
\tightlist
\item
  The data field can contain any extended ASCII data. The default
  interpretation of data by readers is in accordance with ISO/IEC
  8859-1.
\item
  When the \textbf{parse} option is specified, any instances of
  \texttt{\^{}NNN} in the data field are replaced with their equivalent
  ASCII value, useful for specifying unprintable characters.
\item
  The \textbf{eclevel} option is used to specify the error correction
  level:

  \begin{itemize}
  \tightlist
  \item
    \texttt{eclevel=L1} - Lowest
  \item
    \texttt{eclevel=L2}
  \item
    \texttt{eclevel=L3}
  \item
    \texttt{eclevel=L4} - Highest
  \end{itemize}
\item
  The \textbf{version} option is used to specify the size of the symbol,
  1 to 84.
\item
  If unspecified the encoder will select the version of the symbol that
  is the minimum size to represent the given data at the selected error
  correction level.
\end{itemize}

\hypertarget{examples-9}{%
\paragraph{Examples}\label{examples-9}}

\begin{verbatim}
Data:    Han Xin Code
Options: version=10 eclevel=L4
Encoder: hanxin
\end{verbatim}

\includegraphics{images/hanxin-1.eps}

\hypertarget{micropdf417}{%
\subsubsection{MicroPDF417}\label{micropdf417}}

The \textbf{MicroPDF417} barcode symbology is 2D stacked-linear barcode
based on \protect\hyperlink{pdf417}{PDF417} that can encode full 256
character extended-ASCII.

Variants:

\begin{itemize}
\tightlist
\item
  \protect\hyperlink{pdf417}{PDF417} is a larger variant of the
  MicroPDF417 barcode.
\item
  \protect\hyperlink{hibc-symbols}{HIBC MicroPDF417} is a variant of
  MicroPDF417 that should be used when encoding HIBC formatted data.
\end{itemize}

Standards: ISO/IEC 24728, AIM ISS - MicroPDF417.

\hypertarget{data-and-options-13}{%
\paragraph{Data and Options}\label{data-and-options-13}}

\begin{itemize}
\item
  The data field can contain any extended ASCII data. The default
  interpretation of data by readers is in accordance with ISO/IEC
  8859-1. When supported by the receiver characters from other code
  pages can be encoded using Extended Channel Interpretation (ECI).
\item
  When the \textbf{parse} option is specified, any instances of
  \texttt{\^{}NNN} in the data field are replaced with their equivalent
  ASCII value, useful for specifying unprintable characters.
\item
  When the \textbf{parsefnc} option is specified, non-data function
  characters can be specified by escape sequences:

  \begin{itemize}
  \tightlist
  \item
    \texttt{\^{}ECI000000} to \texttt{\^{}ECI811799}: ECI indicators
  \end{itemize}
\item
  The \textbf{columns} and \textbf{rows} options are used to specify the
  size of the symbol. Valid values are:

  \begin{itemize}
  \tightlist
  \item
    \emph{1x11}, \emph{1x14}, \emph{1x17}, \emph{1x20}, \emph{1x24},
    \emph{1x28}, \emph{2x8}, \emph{2x11}, \emph{2x14}, \emph{2x17},
    \emph{2x20}, \emph{2x23}, \emph{2x26}, \emph{3x6}, \emph{3x8},
    \emph{3x10}, \emph{3x12}, \emph{3x15}, \emph{3x20}, \emph{3x26},
    \emph{3x32}, \emph{3x38}, \emph{3x44}, \emph{4x4}, \emph{4x6},
    \emph{4x8}, \emph{4x10}, \emph{4x12}, \emph{4x15}, \emph{4x20},
    \emph{4x26}, \emph{4x32}, \emph{4x38}, \emph{4x44}
  \end{itemize}
\item
  If the \textbf{columns} and \textbf{rows} are unspecified the encoder
  will default to creating a symbol that is the minimum size to
  represent the given data.
\item
  The \textbf{rowmult} option is used to specify how tall each bar is,
  with respect to the minimum module width. The default is 3.
\item
  The \textbf{raw} option denotes that the data field is providing the
  input as a pre-encoded codewords in \texttt{\^{}NNN} format, suitable
  for direct low-level encoding.
\item
  \emph{Deprecated: Internal use.} The \textbf{cca} option identifies
  this symbol as a \emph{CC-A} 2D component of a
  \protect\hyperlink{gs1-composite-symbols}{GS1 Composite} symbol.

  \begin{itemize}
  \tightlist
  \item
    Special size rules apply when the \textbf{cca} option is given, in
    which case the \textbf{columns} and \textbf{rows} options that are
    used to specify the size of the symbol must be one of:

    \begin{itemize}
    \tightlist
    \item
      \emph{2x5}, \emph{2x6}, \emph{2x7}, \emph{2x8}, \emph{2x9},
      \emph{2x10}, \emph{2x12}, \emph{3x4}, \emph{3x5}, \emph{3x6},
      \emph{3x7}, \emph{3x8}, \emph{4x3}, \emph{4x4}, \emph{4x5},
      \emph{4x6}, \emph{4x7}
    \end{itemize}
  \end{itemize}
\item
  \emph{Deprecated: Internal use.} The \textbf{ccb} option identifies
  this symbol as a \emph{CC-B} 2D component of a
  \protect\hyperlink{gs1-composite-symbols}{GS1 Composite} symbol.
\end{itemize}

\hypertarget{examples-10}{%
\paragraph{Examples}\label{examples-10}}

\begin{verbatim}
Data:    MicroPDF417
Options: 
Encoder: micropdf417
\end{verbatim}

\includegraphics{images/micropdf417-1.eps}

\begin{verbatim}
Data:    MicroP^068F417
Options: parse rows=15 columns=4
Encoder: micropdf417
\end{verbatim}

\includegraphics{images/micropdf417-2.eps}

\hypertarget{pdf417}{%
\subsubsection{PDF417}\label{pdf417}}

The \textbf{PDF417} barcode symbology is 2D stacked-linear barcode that
can encode full 256 character extended-ASCII.

Variants:

\begin{itemize}
\tightlist
\item
  \protect\hyperlink{compact-pdf417}{Compact PDF417} is a shortened form
  of the PDF417 barcode that is used in applications where the space for
  the symbol is restricted.
\item
  \protect\hyperlink{micropdf417}{MicroPDF417} is a smaller variant of
  the PDF417 barcode.
\item
  \protect\hyperlink{hibc-symbols}{HIBC PDF417} is a variant of PDF417
  that should be used when encoding HIBC formatted data.
\end{itemize}

Standards: ISO/IEC 15438, DD ENV 12925, AIM USS - PDF417.

\hypertarget{data-and-options-14}{%
\paragraph{Data and Options}\label{data-and-options-14}}

\begin{itemize}
\tightlist
\item
  The data field can contain any extended ASCII data. The default
  interpretation of data by readers is in accordance with ISO/IEC
  8859-1. When supported by the receiver characters from other code
  pages can be encoded using Extended Channel Interpretation (ECI).
\item
  When the \textbf{parse} option is specified, any instances of
  \texttt{\^{}NNN} in the data field are replaced with their equivalent
  ASCII value, useful for specifying unprintable characters.
\item
  When the \textbf{parsefnc} option is specified, non-data function
  characters can be specified by escape sequences:

  \begin{itemize}
  \tightlist
  \item
    \texttt{\^{}ECI000000} to \texttt{\^{}ECI811799}: ECI indicators.
  \end{itemize}
\item
  The \textbf{eclevel} option is used to specify the error correction
  level, from 1 to 5. The default is to choose a standard level of error
  correction that is determined by the encoded data length.
\item
  The \textbf{columns} option specifies the number of columns (or groups
  of bars) in the output symbol, from 1 to 30.
\item
  The \textbf{rows} option specifies the minimum number of rows in the
  symbol, from 3 to 90.
\item
  If \textbf{rows} is unspecified the encoder will select a number that
  creates a symbol that is the minimum size to represent the given data.
\item
  The \textbf{rowmult} option is used to specify how tall each bar is,
  with respect to the minimum module width. The default is 3.
\item
  The \textbf{raw} option denotes that the data field is providing the
  input as a pre-encoded codewords in \texttt{\^{}NNN} format, suitable
  for direct low-level encoding.
\item
  \emph{Deprecated: Use \protect\hyperlink{compact-pdf417}{Compact
  PDF417} instead. The \textbf{compact} option is used to create a
  compact/truncated PDF417 symbol that has fewer bars per row that a
  standard symbol and hence is more narrow.}
\item
  \emph{Deprecated: Internal use.} The \textbf{ccc} option identifies
  this symbol as a \emph{CC-C} 2D component of a
  \protect\hyperlink{gs1-composite-symbols}{GS1 Composite} symbol.
\end{itemize}

\hypertarget{examples-11}{%
\paragraph{Examples}\label{examples-11}}

\begin{verbatim}
Data:    PDF417
Options: 
Encoder: pdf417
\end{verbatim}

\includegraphics{images/pdf417-1.eps}

\begin{verbatim}
Data:    P^068F417
Options: parse columns=2 rows=15
Encoder: pdf417
\end{verbatim}

\includegraphics{images/pdf417-2.eps}

\begin{verbatim}
Data:    Strong error correction
Options: columns=2 eclevel=5
Encoder: pdf417
\end{verbatim}

\includegraphics{images/pdf417-4.eps}

\begin{verbatim}
Data:    ^453^178^121^239
Options: raw columns=2
Encoder: pdf417
\end{verbatim}

\includegraphics{images/pdf417-5.eps}

\hypertarget{compact-pdf417}{%
\subsubsection{Compact PDF417}\label{compact-pdf417}}

\textbf{Compact PDF417} is a shortened form of the
\protect\hyperlink{pdf417}{PDF417} barcode that is used in applications
where the space for the symbol is restricted.

Also known as: Truncated PDF417

Variants:

\begin{itemize}
\tightlist
\item
  \protect\hyperlink{pdf417}{PDF417} is the larger, more popular
  variant.
\item
  \protect\hyperlink{micropdf417}{MicroPDF417} is a smaller variant of
  the PDF417 barcode.
\item
  \protect\hyperlink{hibc-symbols}{HIBC PDF417} is a variant of PDF417
  that should be used when encoding HIBC formatted data.
\end{itemize}

Standards: ISO/IEC 15438, DD ENV 12925, AIM USS - PDF417.

\hypertarget{data-and-options-15}{%
\paragraph{Data and Options}\label{data-and-options-15}}

\begin{itemize}
\tightlist
\item
  The data field can contain any extended ASCII data. The default
  interpretation of data by readers is in accordance with ISO/IEC
  8859-1. When supported by the receiver characters from other code
  pages can be encoded using Extended Channel Interpretation (ECI).
\item
  When the \textbf{parse} option is specified, any instances of
  \texttt{\^{}NNN} in the data field are replaced with their equivalent
  ASCII value, useful for specifying unprintable characters.
\item
  When the \textbf{parsefnc} option is specified, non-data function
  characters can be specified by escape sequences:

  \begin{itemize}
  \tightlist
  \item
    \texttt{\^{}ECI000000} to \texttt{\^{}ECI811799}: ECI indicators
  \end{itemize}
\item
  The \textbf{eclevel} option is used to specify the error correction
  level, from 1 to 5. The default is to choose a standard level of error
  correction that is determined by the encoded data length.
\item
  The \textbf{columns} option specifies the number of columns (or groups
  of bars) in the output symbol, from 1 to 30.
\item
  The \textbf{rows} option specifies the minimum number of rows in the
  symbol, from 3 to 90.
\item
  If \textbf{rows} is unspecified the encoder will select a number that
  creates a symbol that is the minimum size to represent the given data.
\item
  The \textbf{rowmult} option is used to specify how tall each bar is,
  with respect to the minimum module width. The default is 3.
\item
  The \textbf{raw} option denotes that the data field is providing the
  input as a pre-encoded codewords in \texttt{\^{}NNN} format, suitable
  for direct low-level encoding.
\end{itemize}

\hypertarget{examples-12}{%
\paragraph{Examples}\label{examples-12}}

\begin{verbatim}
Data:    A truncated PDF417
Options: columns=4
Encoder: pdf417compact
\end{verbatim}

\includegraphics{images/pdf417-3.eps}

\hypertarget{qr-code}{%
\subsubsection{QR Code}\label{qr-code}}

The \textbf{QR Code} symbology is a 2D matrix-style barcode symbology
that can encode full 256 character extended-ASCII.

Also known as: Quick Response Code.

Variants:

\begin{itemize}
\tightlist
\item
  \protect\hyperlink{micro-qr-code}{Micro QR Code} is a small QR Code
  that is used in applications that require a small symbol space.
\item
  \protect\hyperlink{gs1-qr-code}{GS1 QR Code} is a variant of Data
  Matrix that should be used when encoding data that is in
  \protect\hyperlink{gs1-application-identifier-standard-format}{GS1
  Application Identifier standard format}.
\item
  \protect\hyperlink{gs1-digital-link-qr-code}{GS1 Digital Link QR Code}
  is a variant of Data Matrix that should be used when encoding a GS1
  Digital Link URI.
\item
  \protect\hyperlink{hibc-symbols}{HIBC QR Code} is a variant of QR Code
  that should be used when encoding HIBC formatted data.
\item
  \protect\hyperlink{swiss-qr-code}{Swiss QR Code} is a variant of the
  QR Code for use with the QR-bill application supporting Swiss
  payments.
\end{itemize}

Standards: ISO/IEC 18004, JIS X 0510, ITS - QR Code, AIM ISS - QR Code.

\hypertarget{data-and-options-16}{%
\paragraph{Data and Options}\label{data-and-options-16}}

\begin{itemize}
\tightlist
\item
  The data field can contain any extended ASCII data. The default
  interpretation of data by readers is in accordance with ISO/IEC
  8859-1. When supported by the receiver characters from other code
  pages can be encoded using Extended Channel Interpretation (ECI).
\item
  When the \textbf{parse} option is specified, any instances of
  \texttt{\^{}NNN} in the data field are replaced with their equivalent
  ASCII value, useful for specifying unprintable characters.
\item
  When the \textbf{parsefnc} option is specified, non-data function
  characters can be specified by escape sequences:

  \begin{itemize}
  \tightlist
  \item
    \texttt{\^{}FNC1}: FNC1. \emph{Recommendation: Use
    \protect\hyperlink{gs1-qr-code}{GS1 QR Code} for encoding GS1 AI
    syntax data.}
  \item
    \texttt{\^{}ECI000000} to \texttt{\^{}ECI999999}: ECI indicators
  \end{itemize}
\item
  The \textbf{eclevel} option is used to specify the error correction
  level:

  \begin{itemize}
  \tightlist
  \item
    \texttt{eclevel=L} - Low
  \item
    \texttt{eclevel=M} - Medium (default)
  \item
    \texttt{eclevel=Q} - Quality
  \item
    \texttt{eclevel=H} - High
  \end{itemize}
\item
  The \textbf{version} option is used to specify the size of the symbol,
  1 to 40.
\item
  If unspecified the encoder will select the version of the symbol that
  is the minimum size to represent the given data at the selected error
  correction level.
\item
  \emph{Deprecated}: The \textbf{format} option is used to specify the
  symbol type, either:

  \begin{itemize}
  \tightlist
  \item
    \texttt{full} (default)
  \item
    \texttt{micro}: Use \protect\hyperlink{micro-qr-code}{Micro QR Code}
    for Micro QR Code symbols.
  \item
    \texttt{rmqr}: The standard is still in draft.
  \end{itemize}
\end{itemize}

\hypertarget{examples-13}{%
\paragraph{Examples}\label{examples-13}}

\begin{verbatim}
Data:    QR Code
Options: 
Encoder: qrcode
\end{verbatim}

\begin{verbatim}
Data:    QR ^067ode
Options: parse
Encoder: qrcode
\end{verbatim}

\includegraphics{images/qrcode-1.eps}

\begin{verbatim}
Data:    QR CODE 1234
Options: version=10 eclevel=Q
Encoder: qrcode
\end{verbatim}

\includegraphics{images/qrcode-2.eps}

\hypertarget{micro-qr-code}{%
\subsubsection{Micro QR Code}\label{micro-qr-code}}

The \textbf{Micro QR Code} symbology is a smaller variant of
\protect\hyperlink{qr-code}{QR Code} that is used in applications that
require a small symbol space.

Also known as: Micro Quick Response Code.

Variants:

\begin{itemize}
\tightlist
\item
  \protect\hyperlink{qr-code}{QR Code} is the more popular, larger
  variant.
\end{itemize}

Standards: ISO/IEC 18004, JIS X 0510, ITS - QR Code, AIM ISS - QR Code.

\hypertarget{data-and-options-17}{%
\paragraph{Data and Options}\label{data-and-options-17}}

\begin{itemize}
\tightlist
\item
  The data field can contain any extended ASCII data. The default
  interpretation of data by readers is in accordance with ISO/IEC
  8859-1.
\item
  An appropriate size will be selected to work around the following
  restrictions:

  \begin{itemize}
  \tightlist
  \item
    An M1 symbol is only compatible with numeric data.
  \item
    An M2 symbol is only compatible with alphanumeric data.
  \end{itemize}
\item
  When the \textbf{parse} option is specified, any instances of
  \texttt{\^{}NNN} in the data field are replaced with their equivalent
  ASCII value, useful for specifying unprintable characters.
\item
  The \textbf{version} option is used to specify the size of the symbol,
  either \texttt{version=M1}, \texttt{version=M2}, \texttt{version=M3}
  or \texttt{version=M4}.
\item
  The \textbf{eclevel} option is used to specify the error correction
  level:

  \begin{itemize}
  \tightlist
  \item
    \texttt{eclevel=L} - Low (default)
  \item
    \texttt{eclevel=M} - Medium; Not compatible with M1 symbols
  \item
    \texttt{eclevel=Q} - Quality; Only compatible with M4 symbols
  \end{itemize}
\item
  If unspecified the encoder will select the version of the symbol that
  is the minimum size to represent the given data at the selected error
  correction level.
\end{itemize}

\hypertarget{examples-14}{%
\paragraph{Examples}\label{examples-14}}

\begin{verbatim}
Data:    01234567
Options: 
Encoder: microqrcode
\end{verbatim}

\includegraphics{images/qrcode-5.eps}

\hypertarget{swiss-qr-code}{%
\subsubsection{Swiss QR Code}\label{swiss-qr-code}}

\textbf{Swiss QR Code} is a variant of the
\protect\hyperlink{qr-code}{QR Code} barcode symbology for use with the
QR-bill application supporting Swiss payments.

Standards: Swiss Implementation Guidelines QR-bill, ISO/IEC 18004, ITS -
QR Code

\hypertarget{data-and-options-18}{%
\paragraph{Data and Options}\label{data-and-options-18}}

\begin{itemize}
\tightlist
\item
  The data field input is provided in the format described in the Swiss
  Implementation Guidelines QR-bill document.
\item
  The \textbf{version} option is used to specify the size of the symbol,
  5 to 25. This is not typically specified.
\item
  If the \textbf{version} is unspecified the encoder will default to
  creating a symbol that is the minimum size to represent the given data
  at the selected error correction level.
\end{itemize}

Note: Swiss QR Code specifies that the image shall be 46mm x 46mm
(excluding quiet zone), irrespective of QR Code version, and is
overlayed with a fixed 7mm x 7mm Swiss Cross pattern. The Swiss Cross
pitch bears no relation to the module width of the QR symbol, so we must
assume that the application provides sufficient resolution throughout
the print process that defects produced by issues such as pixel grazing
are insignificant. Therefore, unlike most other encoders, we do not pick
a fixed module width to assist with grid fitting but rather ensure that
the image is scaled to a fixed size, per the specification.

\hypertarget{example-1}{%
\paragraph{Example}\label{example-1}}

\begin{verbatim}
Data:
         SPC
         0200
         1
         CH5800791123000889012
         S
         Robert Schneider AG
         Rue du Lac
         1268
         2501
         Biel
         CH
         
         
         
         
         
         
         
         199.95
         CHF
         K
         Pia-Maria Rutschmann-Schnyder
         Grosse Marktgasse 28
         9400 Rorschach
         
         
         CH
         SCOR
         RF18539007547034
         
         EPD
Options:
Encoder: swissqrcode
\end{verbatim}

\includegraphics{images/swissqrcode-1.eps}

\hypertarget{one-dimensional}{%
\subsection{One-Dimensional}\label{one-dimensional}}

\hypertarget{code-128}{%
\subsubsection{Code 128}\label{code-128}}

\textbf{Code 128} is an arbitrarily long, high-density barcode symbology
that can be used to encode full 256 character extended-ASCII.

Also known as: USD-6, USS-128, Code 128A, Code 128B, Code 128C.

Variants:

\begin{itemize}
\tightlist
\item
  \protect\hyperlink{gs1-128}{GS1-128} is a variant of Code 128 that
  should be used when encoding data that is in
  \protect\hyperlink{gs1-application-identifier-standard-format}{GS1
  Application Identifier standard format}.
\item
  \protect\hyperlink{hibc-symbols}{HIBC Code 128} is a variant of Code
  128 that should be used when encoding HIBC formatted data.
\end{itemize}

Standards: ISO/IEC 15417, ANSI/AIM BC4 - ISS Code 128, BS EN 799.

\hypertarget{data-and-options-19}{%
\paragraph{Data and Options}\label{data-and-options-19}}

\begin{itemize}
\tightlist
\item
  The data field input can consist of any extended ASCII data. The
  default interpretation of data by readers is in accordance with
  ISO/IEC 8859-1.
\item
  The mandatory check digit is calculated automatically.
\item
  When the \textbf{parse} option is specified, any instances of
  \texttt{\^{}NNN} in the data field are replaced with their equivalent
  ASCII or extended-ASCII value, useful for specifying unprintable
  characters, e.g. \texttt{\^{}029} for \emph{GS}, \texttt{\^{}209} for
  \emph{Ñ}, etc.
\item
  \emph{FNC4} function characters will be inserted automatically to
  allow the encoding of extended ASCII characters.
\item
  When the \textbf{parsefnc} option is specified, non-data function
  characters can be specified by escape sequences:

  \begin{itemize}
  \tightlist
  \item
    \texttt{\^{}FNC1}: FNC1. \emph{Recommendation: Use
    \protect\hyperlink{gs1-128}{GS1-128} for encoding GS1 AI syntax
    data.}
  \item
    \texttt{\^{}FNC2}: FNC2
  \item
    \texttt{\^{}FNC3}: FNC3
  \end{itemize}
\item
  The \textbf{raw} option denotes that the data field is providing the
  input as pre-encoded codewords in \texttt{\^{}NNN} format, suitable
  for direct low-level encoding.
\item
  \emph{Deprecated: For internal use.} When the \textbf{parsefnc} option
  is specified, the special pseudo characters \texttt{\^{}LNKA} and
  \texttt{\^{}LNKC} at the end of the symbol indicate that a GS1-128
  symbol includes a \emph{CC-A/B} or \emph{CC-C} GS1 composite 2D
  component.
\end{itemize}

\hypertarget{example-2}{%
\paragraph{Example}\label{example-2}}

\begin{verbatim}
Data:    Count0123456789!
Options: includetext
Encoder: code128
\end{verbatim}

\includegraphics{images/code128-1.eps}

\hypertarget{code-39}{%
\subsubsection{Code 39}\label{code-39}}

The \textbf{Code 39} barcode symbology is discrete, variable length and
self-checking.

Also known as: Code 3 of 9, LOGMARS, Alpha39, USD-3, USS-39.

Variants:

\begin{itemize}
\tightlist
\item
  \protect\hyperlink{code-39-extended}{Code 39 Extended} is a variant of
  Code 39 that can be used to encode full 128 character ASCII with the
  use of shift character combinations.
\item
  \protect\hyperlink{hibc-symbols}{HIBC Code39} is a variant of Code 39
  that should be used when encoding HIBC formatted data.
\item
  AIM USD-2 is a subset of Code 39 containing the characters A-Z, 0-9,
  \emph{space}, \texttt{-} and \texttt{.}.
\end{itemize}

Standards: ISO/IEC 16388, ANSI/AIM BC1 - USS Code 39, BS EN 800, MIL STD
1189.

\hypertarget{data-and-options-20}{%
\paragraph{Data and Options}\label{data-and-options-20}}

\begin{itemize}
\tightlist
\item
  The data field can hold any of the following:

  \begin{itemize}
  \tightlist
  \item
    Numbers 0-9
  \item
    Capital letters A-Z
  \item
    Symbols \texttt{-.\$/+\%*} and \emph{space}
  \end{itemize}
\item
  The \textbf{includecheck} option calculates the check digit.
\item
  The \textbf{includecheckintext} option makes the calculated checksum
  appear in the human readable text.
\item
  The \textbf{hidestars} option suppresses the asterisks in the human
  readable text.
\end{itemize}

\hypertarget{examples-15}{%
\paragraph{Examples}\label{examples-15}}

\begin{verbatim}
Data:    CODE39
Options: includetext
Encoder: code39
\end{verbatim}

\includegraphics{images/code39-1.eps}

\begin{verbatim}
Data:    CODE39
Options: includecheck includetext includecheckintext
Encoder: code39
\end{verbatim}

\includegraphics{images/code39-2.eps}

\begin{verbatim}
Data:    CODE39
Options: hidestars includecheck includetext
Encoder: code39
\end{verbatim}

\includegraphics{images/code39-3.eps}

\hypertarget{code-39-extended}{%
\subsubsection{Code 39 Extended}\label{code-39-extended}}

The \textbf{Code 39 Extended} barcode symbology is discrete, variable
length and self-checking. It is based on
\protect\hyperlink{code-39}{Code 39} but can encode full 128 character
ASCII by using shift combinations.

Also known as: Code 39 Full ASCII.

Variants:

\begin{itemize}
\tightlist
\item
  \protect\hyperlink{code-39}{Code 39} is a simpler variant of Code 39
  Extended.
\end{itemize}

Standards: ISO/IEC 16388, ANSI/AIM BC1 - USS Code 39, BS EN 800.

\hypertarget{data-and-options-21}{%
\paragraph{Data and Options}\label{data-and-options-21}}

\begin{itemize}
\tightlist
\item
  The data field can consist of any ASCII data.
\item
  When the \textbf{parse} option is specified, any instances of
  \texttt{\^{}NNN} in the data field are replaced with their equivalent
  ASCII value, useful for specifying unprintable characters, e.g.
  \texttt{\^{}029} for \emph{GS}, etc.
\item
  The \textbf{includecheck} option calculates the check digit.
\item
  The \textbf{includecheckintext} causes the calculated check digit to
  appear in the human readable text.
\item
  The \textbf{hidestars} option suppresses the asterisks in the human
  readable text.
\end{itemize}

\hypertarget{examples-16}{%
\paragraph{Examples}\label{examples-16}}

\begin{verbatim}
Data:    Code39 Ext!
Options: includetext includecheck
Encoder: code39ext
\end{verbatim}

\includegraphics{images/code39ext-1.eps}

\begin{verbatim}
Data:    Code39^029Extended
Options: parse includecheck
Encoder: code39ext
\end{verbatim}

\includegraphics{images/code39ext-2.eps}

\hypertarget{code-93}{%
\subsubsection{Code 93}\label{code-93}}

\textbf{Code 93} is a continuous, variable length, self-checking barcode
symbology.

Also known as: USD-7, USS-93.

Variants:

\begin{itemize}
\tightlist
\item
  \protect\hyperlink{code-93-extended}{Code 93 Extended} is a variant of
  Code 93 that can be used to encode full 128 character ASCII with the
  use of special shift character combinations.
\end{itemize}

Standards: ANSI/AIM BC5 - USS Code 93, ITS 93i.

\hypertarget{data-and-options-22}{%
\paragraph{Data and Options}\label{data-and-options-22}}

\begin{itemize}
\tightlist
\item
  The data field can hold any of the following:

  \begin{itemize}
  \tightlist
  \item
    Numbers 0-9
  \item
    Capital letters A-Z
  \item
    Symbols \texttt{-.\$/+\%*} and \emph{space}
  \end{itemize}
\item
  The \textbf{includecheck} option calculates the two check digits.
\item
  \emph{Deprecated: For internal use.} The \textbf{parsefnc} option
  allows the special shift characters to be supplied as escape
  sequences:

  \begin{itemize}
  \tightlist
  \item
    \texttt{\^{}SFT\$}: (\$)
  \item
    \texttt{\^{}SFT\%}: (\%)
  \item
    \texttt{\^{}SFT/}: (/)
  \item
    \texttt{\^{}SFT+}: (+)
  \end{itemize}
\end{itemize}

\hypertarget{examples-17}{%
\paragraph{Examples}\label{examples-17}}

\begin{verbatim}
Data:    CODE93
Options: includecheck
Encoder: code93
\end{verbatim}

\includegraphics{images/code93-1.eps}

Code 93 including a special shift combination \texttt{(/)A} representing
\emph{!}:

\begin{verbatim}
Data:    CODE93^SFT/A
Options: parsefnc includecheck
Encoder: code93
\end{verbatim}

\includegraphics{images/code93-2.eps}

\hypertarget{code-93-extended}{%
\subsubsection{Code 93 Extended}\label{code-93-extended}}

The \textbf{Code 93 Extended} barcode symbology is continuous, variable
length and self-checking. It is based on
\protect\hyperlink{code-93}{Code 93} but can encode full 128 character
ASCII using four additional shift characters: \emph{(\$)} \emph{(\%)}
\emph{(/)} \emph{(+)}

Also known as: Code 93 Full ASCII.

Variants:

\begin{itemize}
\tightlist
\item
  \protect\hyperlink{code-93}{Code 93} is a simpler variant of the Code
  93 Extended barcode symbology.
\end{itemize}

Standards: ANSI/AIM BC5 - USS Code 93, ITS 93i.

\hypertarget{data-and-options-23}{%
\paragraph{Data and Options}\label{data-and-options-23}}

\begin{itemize}
\tightlist
\item
  The data field can consist of any ASCII data.
\item
  When the \textbf{parse} option is specified, any instances of
  \texttt{\^{}NNN} in the data field are replaced with their equivalent
  ASCII value, useful for specifying unprintable characters, e.g.
  \texttt{\^{}029} for \emph{GS}, etc.
\item
  The \textbf{includecheck} option calculates the two check digits.
\end{itemize}

\hypertarget{examples-18}{%
\paragraph{Examples}\label{examples-18}}

\begin{verbatim}
Data:    Code93Ext!
Options: includecheck
Encoder: code93ext
\end{verbatim}

\includegraphics{images/code93ext-1.eps}

\begin{verbatim}
Data:    Code93^029Extended
Options: parse includecheck
Encoder: code93ext
\end{verbatim}

\includegraphics{images/code93ext-2.eps}

\hypertarget{interleaved-2-of-5}{%
\subsubsection{Interleaved 2 of 5}\label{interleaved-2-of-5}}

\textbf{Interleaved 2 of 5} is a high-density numeric barcode symbology.

Also known as: ITF, Code 2 of 5 Interleaved, USD-1, USS-Interleaved 2 of
5.

Variants:

\begin{itemize}
\tightlist
\item
  \protect\hyperlink{itf-14}{ITF-14} is a variant of Interleaved 2 of 5
  that should be used when encoding a fourteen-digit GTIN.
\end{itemize}

Standards: ISO/IEC 16390, ANSI/AIM BC2 - USS Interleaved 2 of 5, BS EN
801.

\hypertarget{data-and-options-24}{%
\paragraph{Data and Options}\label{data-and-options-24}}

\begin{itemize}
\tightlist
\item
  The data can consist of any number of digits.
\item
  If the length of the symbol including the possible check digit would
  be odd then the data is prefixed by \emph{0}.
\item
  The \textbf{includecheck} option calculates the check digit.
\item
  The \textbf{includecheckintext} option makes the calculated checksum
  appear in the human readable text.
\end{itemize}

\hypertarget{examples-19}{%
\paragraph{Examples}\label{examples-19}}

\begin{verbatim}
Data:    0123456789
Options: 
Encoder: interleaved2of5
\end{verbatim}

\includegraphics{images/itf-1.eps}

\begin{verbatim}
Data:    2401234567
Options: includecheck includetext includecheckintext
Encoder: interleaved2of5
\end{verbatim}

\includegraphics{images/itf-2.eps}

\hypertarget{supply-chain}{%
\subsection{Supply Chain}\label{supply-chain}}

\hypertarget{ean-14}{%
\subsubsection{EAN-14}\label{ean-14}}

\textbf{EAN-14} is an implementation of the
\protect\hyperlink{gs1-128}{GS1-128} barcode symbology with \emph{AI
(01)} that is used to encode a lone GTIN-14.

Also known as: UCC-14.

Standards: ISO/IEC 15417, ANSI/AIM BC4-1999 ISS, BS EN 799, GS1 General
Specifications.

\hypertarget{data-and-options-25}{%
\paragraph{Data and Options}\label{data-and-options-25}}

\begin{itemize}
\tightlist
\item
  The data field input must be a solitary \emph{AI (01)} with thirteen
  or fourteen digits of a GTIN, i.e. \texttt{(01)...}.
\item
  Arbitrary spacing may be placed between the digits to format the human
  readable text without interfering with the encoded data.
\item
  If thirteen digits of primary data are supplied then the check digit
  is calculated automatically. Otherwise the provided check digit will
  be verified.
\end{itemize}

\hypertarget{examples-20}{%
\paragraph{Examples}\label{examples-20}}

Identical symbols, input provided with and without a check digit:

\begin{verbatim}
Data:    (01)04601234567893
Options: includetext
Encoder: ean14
\end{verbatim}

\begin{verbatim}
Data:    (01)0460123456789
Options: includetext
Encoder: ean14
\end{verbatim}

\includegraphics{images/ean14-1.eps}

\hypertarget{itf-14}{%
\subsubsection{ITF-14}\label{itf-14}}

\textbf{ITF-14} is an implementation of the
\protect\hyperlink{interleaved-2-of-5}{Interleaved 2 of 5} barcode
symbology that is typically used to encode a GTIN-14, GTIN-13 or
GTIN-12.

Also known as: UPC Shipping Container Symbol, SCS, UPC Case Code.

Standards: ISO/IEC 16390, ANSI/AIM BC2-1995 USS, BS EN 801, GS1 General
Specifications.

\hypertarget{data-and-options-26}{%
\paragraph{Data and Options}\label{data-and-options-26}}

\begin{itemize}
\tightlist
\item
  The data consists of either thirteen or fourteen digits.
\item
  Arbitrary spacing may be placed between the digits to format the human
  readable text without interfering with the encoded data.
\item
  If thirteen digits are supplied then the check digit is calculated
  automatically. Otherwise the provided check digit will be verified.
\end{itemize}

\hypertarget{examples-21}{%
\paragraph{Examples}\label{examples-21}}

Identical symbols, input provided with and without a check digit:

\begin{verbatim}
Data:    04601234567893
Options: includetext
Encoder: itf14
\end{verbatim}

\begin{verbatim}
Data:    0460123456789
Options: includetext
Encoder: itf14
\end{verbatim}

\includegraphics{images/itf-14-1.eps}

\hypertarget{sscc-18}{%
\subsubsection{SSCC-18}\label{sscc-18}}

\textbf{SSCC-18} is an implementation of the
\protect\hyperlink{gs1-128}{GS1-128} barcode symbology with \emph{AI
(00)} that is typically used to encode an eighteen-digit shipping
container serial number.

Also known as: EAN-18, NVE.

Standards: ISO/IEC 15417, ANSI/AIM BC4-1999 ISS, BS EN 799, GS1 General
Specifications.

\hypertarget{data-and-options-27}{%
\paragraph{Data and Options}\label{data-and-options-27}}

\begin{itemize}
\tightlist
\item
  The data field input must be a solitary \emph{AI (00)} with seventeen
  or eighteen digits of a Serial Shipping Container Code, i.e.
  \texttt{(00)...}.
\item
  Arbitrary spacing may be placed between the digits to format the human
  readable text without interfering with the encoded data.
\item
  If seventeen digits of primary data are supplied then the check digit
  is calculated automatically. Otherwise the provided check digit will
  be verified.
\end{itemize}

\hypertarget{example-3}{%
\paragraph{Example}\label{example-3}}

Identical symbols, input provided with and without a check digit:

\begin{verbatim}
Data:    (00)006141411234567890
Options: includetext
Encoder: sscc18
\end{verbatim}

\begin{verbatim}
Data:    (00)00614141123456789
Options: includetext
Encoder: sscc18
\end{verbatim}

\includegraphics{images/sscc18-1.eps}

\hypertarget{gs1-datamatrix}{%
\subsubsection{GS1 DataMatrix}\label{gs1-datamatrix}}

\textbf{GS1 DataMatrix} is an implementation of the
\protect\hyperlink{data-matrix}{Data Matrix} (ECC 200) barcode symbology
with \protect\hyperlink{gs1-application-identifier-standard-format}{GS1
formatted} data.

Standards: ISO/IEC 16022, ANSI/AIM BC11 ISS, GS1 General Specifications.

\hypertarget{data-and-options-28}{%
\paragraph{Data and Options}\label{data-and-options-28}}

\begin{itemize}
\tightlist
\item
  The data field input is provided in
  \protect\hyperlink{gs1-application-identifier-standard-format}{GS1
  Application Identifier standard format} starting with the mandatory
  \emph{(01)} Application Identifier.
\item
  Note that it is not necessary to specify \emph{FNC1} characters since
  the encoder will insert these automatically where necessary.
\item
  \emph{Deprecated.} The \textbf{dontlint} option disables GS1
  Application Identifier syntax validation allowing images to be
  generated for input that does not comply with GS1 standards.
\item
  The \textbf{format} option is used to specify the shape of the symbol,
  either \texttt{square} (default) or \texttt{rectangle}.
\item
  The \textbf{columns} and \textbf{rows} options are used to specify the
  size of the symbol.
\item
  The \textbf{version} option can also be used to specify the symbol
  size, as \texttt{version=RxC}. Valid options are:

  \begin{itemize}
  \tightlist
  \item
    With \texttt{format=square}: \emph{10x10}, \emph{12x12},
    \emph{14x14}, \emph{16x16}, \emph{18x18}, \emph{20x20},
    \emph{22x22}, \emph{24x24}, \emph{26x26}, \emph{32x32},
    \emph{36x36}, \emph{40x40}, \emph{44x44}, \emph{48x48},
    \emph{52x52}, \emph{64x64}, \emph{72x72}, \emph{80x80},
    \emph{88x88}, \emph{96x96}, \emph{104x104}, \emph{120x120},
    \emph{132x132}, \emph{144x144}
  \item
    With \texttt{format=rectangle}: \emph{8x18}, \emph{8x32},
    \emph{12x26}, \emph{12x36}, \emph{16x36}, \emph{16x48}
  \end{itemize}
\item
  If \textbf{columns}, \textbf{rows} and \textbf{version} are
  unspecified the encoder will default to creating a symbol of the
  specified \textbf{format} that is the minimum size to represent the
  given data.
\end{itemize}

\hypertarget{example-4}{%
\paragraph{Example}\label{example-4}}

\begin{verbatim}
Data:    (01)95012345678903(3103)000123
Options: 
Encoder: gs1datamatrix
\end{verbatim}

\includegraphics{images/gs1datamatrix-1.eps}

\hypertarget{gs1-qr-code}{%
\subsubsection{GS1 QR Code}\label{gs1-qr-code}}

\textbf{GS1 QR Code} is an implementation of the
\protect\hyperlink{qr-code}{QR Code} barcode symbology with
\protect\hyperlink{gs1-application-identifier-standard-format}{GS1
formatted data}.

Standards: ISO/IEC 18004, ITS - QR Code, GS1 General Specifications.

\hypertarget{data-and-options-29}{%
\paragraph{Data and Options}\label{data-and-options-29}}

\begin{itemize}
\tightlist
\item
  The data field input is provided in
  \protect\hyperlink{gs1-application-identifier-standard-format}{GS1
  Application Identifier standard format} starting with the mandatory
  \emph{(01)} and \emph{(8200)} Application Identifiers.
\item
  Note that it is not necessary to specify \emph{FNC1} characters since
  the encoder will insert these automatically where necessary.
\item
  \emph{Deprecated.} The \textbf{dontlint} option disables GS1
  Application Identifier syntax validation allowing images to be
  generated for input that does not comply with GS1 standards.
\item
  The \textbf{eclevel} option is used to specify the error correction
  level:

  \begin{itemize}
  \tightlist
  \item
    \texttt{eclevel=L} - Low
  \item
    \texttt{eclevel=M} - Medium (default)
  \item
    \texttt{eclevel=Q} - Quality
  \item
    \texttt{eclevel=H} - High
  \end{itemize}
\item
  The \textbf{version} option is used to specify the size of the symbol,
  1 to 40.
\item
  If the \textbf{version} is unspecified the encoder will default to
  creating a symbol that is the minimum size to represent the given data
  at the selected error correction level.
\end{itemize}

\hypertarget{example-5}{%
\paragraph{Example}\label{example-5}}

\begin{verbatim}
Data:    (01)03453120000011(8200)http://www.abc.net
Options: 
Encoder: gs1qrcode
\end{verbatim}

\includegraphics{images/gs1qrcode-1.eps}

\begin{verbatim}
Data:    (01)03453120000011(8200)http://abc.net(10)XYZ(410)9501101020917
Options: 
Encoder: gs1qrcode
\end{verbatim}

\includegraphics{images/gs1qrcode-2.eps}

\hypertarget{gs1-digital-link-datamatrix}{%
\subsubsection{GS1 Digital Link
DataMatrix}\label{gs1-digital-link-datamatrix}}

\textbf{GS1 Digital Link DataMatrix} is an implementation of the
\protect\hyperlink{data-matrix}{Data Matrix} (ECC 200) barcode symbology
containing a GS1 Digital Link URI.

Standards: ISO/IEC 16022, ANSI/AIM BC11 ISS, GS1 Digital Link Standard:
URI Syntax.

\hypertarget{data-and-options-30}{%
\paragraph{Data and Options}\label{data-and-options-30}}

\begin{itemize}
\tightlist
\item
  The data field input is provided as a GS1 Digital Link URI.
\item
  \emph{Deprecated.} The \textbf{dontlint} option disables GS1
  Application Identifier syntax validation allowing images to be
  generated for input that does not comply with GS1 standards.
\item
  The \textbf{format} option is used to specify the shape of the symbol,
  either \texttt{square} (default) or \texttt{rectangle}.
\item
  The \textbf{columns} and \textbf{rows} options are used to specify the
  size of the symbol.
\item
  The \textbf{version} option can also be used to specify the symbol
  size, as \texttt{version=RxC}. Valid options are:

  \begin{itemize}
  \tightlist
  \item
    With \texttt{format=square}: \emph{10x10}, \emph{12x12},
    \emph{14x14}, \emph{16x16}, \emph{18x18}, \emph{20x20},
    \emph{22x22}, \emph{24x24}, \emph{26x26}, \emph{32x32},
    \emph{36x36}, \emph{40x40}, \emph{44x44}, \emph{48x48},
    \emph{52x52}, \emph{64x64}, \emph{72x72}, \emph{80x80},
    \emph{88x88}, \emph{96x96}, \emph{104x104}, \emph{120x120},
    \emph{132x132}, \emph{144x144}
  \item
    With \texttt{format=rectangle}: \emph{8x18}, \emph{8x32},
    \emph{12x26}, \emph{12x36}, \emph{16x36}, \emph{16x48}
  \end{itemize}
\item
  If \textbf{columns}, \textbf{rows} and \textbf{version} are
  unspecified the encoder will default to creating a symbol of the
  specified \textbf{format} that is the minimum size to represent the
  given data.
\end{itemize}

\hypertarget{example-6}{%
\paragraph{Example}\label{example-6}}

\begin{verbatim}
Data:    https://id.gs1.org/01/12312312312333/22/ABC%2d123?99=XYZ%2f987#test
Options: 
Encoder: gs1dldatamatrix
\end{verbatim}

\hypertarget{gs1-qr-digital-link-code}{%
\subsubsection{GS1 QR Digital Link
Code}\label{gs1-qr-digital-link-code}}

\textbf{GS1 Digital Link QR Code} is an implementation of the
\protect\hyperlink{qr-code}{QR Code} barcode symbology containing a GS1
Digital Link URI.

Standards: ISO/IEC 18004, ITS - QR Code, GS1 Digital Link Standard: URI
Syntax.

\hypertarget{data-and-options-31}{%
\paragraph{Data and Options}\label{data-and-options-31}}

\begin{itemize}
\tightlist
\item
  The data field input is provided as a GS1 Digital Link URI.
\item
  \emph{Deprecated.} The \textbf{dontlint} option disables GS1
  Application Identifier syntax validation allowing images to be
  generated for input that does not comply with GS1 standards.
\item
  The \textbf{eclevel} option is used to specify the error correction
  level:

  \begin{itemize}
  \tightlist
  \item
    \texttt{eclevel=L} - Low
  \item
    \texttt{eclevel=M} - Medium (default)
  \item
    \texttt{eclevel=Q} - Quality
  \item
    \texttt{eclevel=H} - High
  \end{itemize}
\item
  The \textbf{version} option is used to specify the size of the symbol,
  1 to 40.
\item
  If the \textbf{version} is unspecified the encoder will default to
  creating a symbol that is the minimum size to represent the given data
  at the selected error correction level.
\end{itemize}

\hypertarget{example-7}{%
\paragraph{Example}\label{example-7}}

\begin{verbatim}
Data:    https://id.gs1.org/01/12312312312333/22/ABC%2d123?99=XYZ%2f987#test
Options: 
Encoder: gs1dlqrcode
\end{verbatim}

\hypertarget{gs1-128}{%
\subsubsection{GS1-128}\label{gs1-128}}

\textbf{GS1-128} is an implementation of the
\protect\hyperlink{code-128}{Code 128} barcode symbology which carries
\protect\hyperlink{gs1-application-identifier-standard-format}{GS1
formatted} data, including a GTIN-14.

Also known as: UCC/EAN-128, EAN-128, UCC-128.

Variants:

\begin{itemize}
\tightlist
\item
  \protect\hyperlink{gs1-composite-symbols}{GS1-128 Composite} is a
  variant of GS1-128 that should be used when a CC-A, CC-B or CC-C GS1
  composite 2D component is required.
\item
  \protect\hyperlink{ean-14}{EAN-14} is a variant of GS1-128 that should
  be used when encoding a lone fourteen-digit GTIN.
\item
  \protect\hyperlink{sscc-18}{SSCC-18} is a variant of GS1-128 that
  should be used when encoding a lone eighteen-digit SSCC.
\end{itemize}

Standards: ISO/IEC 15417, ANSI/AIM BC4-1999 ISS, BS EN 799, GS1 General
Specifications.

\hypertarget{data-and-options-32}{%
\paragraph{Data and Options}\label{data-and-options-32}}

\begin{itemize}
\tightlist
\item
  The data field input is provided in
  \protect\hyperlink{gs1-application-identifier-standard-format}{GS1
  Application Identifier standard format}.
\item
  Note that it is not necessary to specify \emph{FNC1} characters since
  the encoder will insert these automatically where necessary.
\item
  \emph{Deprecated.} The \textbf{dontlint} option disables GS1
  Application Identifier syntax validation allowing images to be
  generated for input that does not comply with GS1 standards.
\item
  \emph{Deprecated: For internal use.} The \textbf{linkagea} option
  specifies that the symbol includes a \emph{CC-A} or \emph{CC-B} GS1
  composite 2D component.
\item
  \emph{Deprecated: For internal use.} The \textbf{linkagec} option
  specifies that the symbol includes a \emph{CC-C} GS1 composite 2D
  component.
\end{itemize}

\hypertarget{examples-22}{%
\paragraph{Examples}\label{examples-22}}

GTIN \emph{95012345678903}; Weight \emph{0.123kg}:

\begin{verbatim}
Data:    (01)95012345678903(3103)000123
Options: includetext
Encoder: gs1-128
\end{verbatim}

\includegraphics{images/ean128-1.eps}

GTIN \emph{0061414199996}; Expiration date \emph{1st Jan 2010}; Batch
\emph{123ABC}; Serial \emph{1234567890}:

\begin{verbatim}
Data:    (01)0061414199996(17)100101(10)123ABC(21)1234567890
Options: includetext
Encoder: gs1-128
\end{verbatim}

\includegraphics{images/ean128-2.eps}

\hypertarget{gs1-dotcode}{%
\subsubsection{GS1 DotCode}\label{gs1-dotcode}}

\textbf{GS1 DotCode} is an implementation of the
\protect\hyperlink{dotcode}{DotCode} barcode symbology with
\protect\hyperlink{gs1-application-identifier-standard-format}{GS1
formatted} data.

Standards: AIM - ISS DotCode, GS1 General Specifications.

\hypertarget{data-and-options-33}{%
\paragraph{Data and Options}\label{data-and-options-33}}

\begin{itemize}
\tightlist
\item
  The data field input is provided in
  \protect\hyperlink{gs1-application-identifier-standard-format}{GS1
  Application Identifier standard format}.
\item
  Note that it is not necessary to specify \emph{FNC1} characters since
  the encoder will insert these automatically where necessary.
\item
  \emph{Deprecated.} The \textbf{dontlint} option disables GS1
  Application Identifier syntax validation allowing images to be
  generated for input that does not comply with GS1 standards.
\item
  The \textbf{ratio} is used to specify the aspect ratio of the symbol.
  The default is 1.5.
\item
  The \textbf{columns} and \textbf{rows} options are used to specify the
  size of the symbol. When these are not specified a symbol is generated
  that maintains the selected aspect ratio.
\item
  \emph{Deprecated: Use only if performance is a concern.} The
  \textbf{fast} option is used to enable the optional high-speed symbol
  masking algorithm.
\end{itemize}

\hypertarget{examples-23}{%
\paragraph{Examples}\label{examples-23}}

\begin{verbatim}
Data:    (235)5vBZIF%!<B;?oa%(01)01234567890128(8008)19052001
Options: rows=16
Encoder: gs1dotcode
\end{verbatim}

\includegraphics{images/gs1dotcode-1.eps}

\hypertarget{gs1-databar-family}{%
\subsection{GS1 DataBar Family}\label{gs1-databar-family}}

\hypertarget{gs1-databar-omnidirectional}{%
\subsubsection{GS1 DataBar
Omnidirectional}\label{gs1-databar-omnidirectional}}

\textbf{GS1 DataBar Omnidirectional} is a fixed-length, linear barcode
symbology that can be used to encode a GTIN-14 for use at point of sale.

Also known as: RSS-14

Variants:

\begin{itemize}
\tightlist
\item
  \protect\hyperlink{gs1-databar-stacked-omnidirectional}{GS1 DataBar
  Stacked Omnidirectional} is a variant of GS1 DataBar Omnidirectional
  for use where a taller, narrower symbol is required.
\item
  \protect\hyperlink{gs1-composite-symbols}{GS1 DataBar Omnidirectional
  Composite} is a variant of GS1 DataBar Omnidirectional that should be
  used when a CC-A or CC-B GS1 composite 2D component is required.
\end{itemize}

Standards: ISO/IEC 24724, ITS Reduced Space Symbology (RSS), AIM ISS -
Reduced Space Symbology (RSS), GS1 General Specifications.

\hypertarget{data-and-options-34}{%
\paragraph{Data and Options}\label{data-and-options-34}}

\begin{itemize}
\tightlist
\item
  The data field input is must be a solitary \emph{AI (01)} with
  thirteen or fourteen digits of a GTIN, i.e. \texttt{(01)...}.
\item
  If thirteen digits of \emph{AI (01)} data are supplied then the check
  digit is calculated automatically, otherwise the provided check digit
  will be verified.
\item
  \emph{Deprecated: For internal use.} The \textbf{linkage} option
  signifies the presence of a GS1 composite 2D component.
\end{itemize}

\hypertarget{examples-24}{%
\paragraph{Examples}\label{examples-24}}

Identical symbols, input provided with and without a check digit:

\begin{verbatim}
Data:    (01)24012345678905
Options: 
Encoder: databaromni
\end{verbatim}

\begin{verbatim}
Data:    (01)2401234567890
Options: 
Encoder: databaromni
\end{verbatim}

\includegraphics{images/databaromni-1.eps}

\hypertarget{gs1-databar-stacked-omnidirectional}{%
\subsubsection{GS1 DataBar Stacked
Omnidirectional}\label{gs1-databar-stacked-omnidirectional}}

\textbf{GS1 DataBar Stacked Omnidirectional} is a fixed-length, stacked
linear barcode symbology that can be used to encode a GTIN-14 for use a
point of sale.

Also known as: RSS-14 Stacked Omnidirectional.

Variants:

\begin{itemize}
\tightlist
\item
  \protect\hyperlink{gs1-databar-omnidirectional}{GS1 DataBar
  Omnidirectional} is a variant of GS1 DataBar Stacked Omnidirectional
  for use where a shorter, wider symbol is required.
\item
  \protect\hyperlink{gs1-composite-symbols}{GS1 DataBar Stacked
  Omnidirectional Composite} is a variant of GS1 DataBar Stacked
  Omnidirectional that should be used when a CC-A or CC-B
  \protect\hyperlink{gs1-composite-symbols}{GS1 composite} 2D component
  is required.
\end{itemize}

Standards: ISO/IEC 24724, ITS Reduced Space Symbology (RSS), AIM ISS -
Reduced Space Symbology (RSS), GS1 General Specifications.

\hypertarget{data-and-options-35}{%
\paragraph{Data and Options}\label{data-and-options-35}}

\begin{itemize}
\tightlist
\item
  The data field input must be a solitary \emph{AI (01)} with thirteen
  or fourteen digits of a GTIN, i.e. \texttt{(01)...}.
\item
  If thirteen digits of \emph{AI (01)} data are supplied then the check
  digit is calculated automatically, otherwise the provided check digit
  will be verified.
\item
  \emph{Deprecated: For internal use.} The \textbf{linkage} option
  signifies the presence of a GS1 composite 2D component.
\end{itemize}

\hypertarget{examples-25}{%
\paragraph{Examples}\label{examples-25}}

Identical symbols, input provided with and without a check digit:

\begin{verbatim}
Data:    (01)24012345678905
Options: 
Encoder: databarstackedomni
\end{verbatim}

\begin{verbatim}
Data:    (01)2401234567890
Options: 
Encoder: databarstackedomni
\end{verbatim}

\includegraphics{images/databarstackedomni-1.eps}

\hypertarget{gs1-databar-expanded}{%
\subsubsection{GS1 DataBar Expanded}\label{gs1-databar-expanded}}

\textbf{GS1 DataBar Expanded} is a variable-length, linear barcode
symbology that can be used to encode a GTIN-14 alongside a number of
other application identifiers for use at point of sale.

Also known as: RSS Expanded.

Variants:

\begin{itemize}
\tightlist
\item
  \protect\hyperlink{gs1-databar-expanded-stacked}{GS1 DataBar Expanded
  Stacked} is a variant of GS1 DataBar Expanded for use where a taller,
  narrower symbol is required.
\item
  \protect\hyperlink{gs1-composite-symbols}{GS1 DataBar Expanded
  Composite} is a variant of GS1 DataBar Expanded that should be used
  when a CC-A or CC-B GS1 composite 2D component is required.
\item
  \protect\hyperlink{gs1-north-american-coupon-code}{GS1 North American
  Coupon Code} is an applicaiton of GS1 DataBar Expanded for use with a
  paperless coupon system.
\end{itemize}

Standards: ISO/IEC 24724, ITS Reduced Space Symbology (RSS), AIM ISS -
Reduced Space Symbology (RSS), GS1 General Specifications.

\hypertarget{data-and-options-36}{%
\paragraph{Data and Options}\label{data-and-options-36}}

\begin{itemize}
\tightlist
\item
  The data field input is provided in
  \protect\hyperlink{gs1-application-identifier-standard-format}{GS1
  Application Identifier standard format}.
\item
  Note that it is not necessary to specify \emph{FNC1} characters since
  the encoder will insert these automatically where necessary.
\item
  \emph{Deprecated.} The \textbf{dontlint} option disables GS1
  Application Identifier syntax validation allowing images to be
  generated for input that does not comply with GS1 standards.
\item
  If the data contains a number of application identifiers matching any
  of the specifications below then they should be provided in this given
  order for maximum encoding efficiency:

  \begin{itemize}
  \tightlist
  \item
    \texttt{(01)9...(3103)...}
  \item
    \texttt{(01)9...(3202)...}
  \item
    \texttt{(01)9...(3203)...}
  \item
    \texttt{(01)9...(310x/320x)...(11/13/15/17)...}
  \item
    \texttt{(01)9...(310x/320x)...}
  \item
    \texttt{(01)9...(392x)...}
  \item
    \texttt{(01)9...(393x)...}
  \item
    \texttt{(01)...}
  \end{itemize}
\item
  \emph{Deprecated: For internal use.} The \textbf{linkage} option
  signifies the presence of a GS1 composite 2D component.
\end{itemize}

\hypertarget{examples-26}{%
\paragraph{Examples}\label{examples-26}}

\begin{verbatim}
Data:    (01)95012345678903(3103)000123
Options: 
Encoder: databarexpanded
\end{verbatim}

\includegraphics{images/databarexpanded-1.eps}

\hypertarget{gs1-databar-expanded-stacked}{%
\subsubsection{GS1 DataBar Expanded
Stacked}\label{gs1-databar-expanded-stacked}}

\textbf{GS1 DataBar Expanded Stacked} is a variable-length,
stacked-linear barcode symbology that can be used to encode a GTIN-14
alongside a number of other application identifiers for use at point of
sale.

Also known as: RSS Expanded Stacked.

Variants:

\begin{itemize}
\tightlist
\item
  \protect\hyperlink{gs1-databar-expanded}{GS1 DataBar Expanded} is a
  variant of GS1 DataBar Expanded Stacked for use where a shorter, wider
  symbol is required.
\item
  \protect\hyperlink{gs1-composite-symbols}{GS1 DataBar Expanded Stacked
  Composite} is a variant of GS1 DataBar Expanded Stacked that should be
  used when a CC-A or CC-B GS1 composite 2D component is required.
\item
  \protect\hyperlink{gs1-north-american-coupon-code}{GS1 North American
  Coupon Code} is an applicaiton of GS1 DataBar Expanded for use with a
  paperless coupon system.
\end{itemize}

Standards: ISO/IEC 24724, ITS Reduced Space Symbology (RSS), AIM ISS -
Reduced Space Symbology (RSS), GS1 General Specifications.

\hypertarget{data-and-options-37}{%
\paragraph{Data and Options}\label{data-and-options-37}}

\begin{itemize}
\tightlist
\item
  The data field input is provided in
  \protect\hyperlink{gs1-application-identifier-standard-format}{GS1
  Application Identifier standard format}.
\item
  Note that it is not necessary to specify \emph{FNC1} characters since
  the encoder will insert these automatically where necessary.
\item
  \emph{Deprecated.} The \textbf{dontlint} option disables GS1
  Application Identifier syntax validation allowing images to be
  generated for input that does not comply with GS1 standards.
\item
  If the data contains a number of application identifiers matching any
  of the specifications below then they should be provided in this given
  order for maximum encoding efficiency:

  \begin{itemize}
  \tightlist
  \item
    \texttt{(01)9...(3103)...}
  \item
    \texttt{(01)9...(3202)...}
  \item
    \texttt{(01)9...(3203)...}
  \item
    \texttt{(01)9...(310x/320x)...(11/13/15/17)...}
  \item
    \texttt{(01)9...(310x/320x)...}
  \item
    \texttt{(01)9...(392x)...}
  \item
    \texttt{(01)9...(393x)...}
  \item
    \texttt{(01)...}
  \end{itemize}
\item
  The \textbf{segments} option is used to specify the maximum number of
  segments per row which must be an even number. The default is
  \emph{4}.
\item
  \emph{Deprecated: For internal use.} The \textbf{linkage} option
  signifies the presence of a GS1 composite 2D component.
\end{itemize}

\hypertarget{examples-27}{%
\paragraph{Examples}\label{examples-27}}

\begin{verbatim}
Data:    (01)95012345678903(3103)000123
Options: segments=4
Encoder: databarexpandedstacked
\end{verbatim}

\includegraphics{images/databarexpandedstacked-1.eps}

\hypertarget{gs1-databar-truncated}{%
\subsubsection{GS1 DataBar Truncated}\label{gs1-databar-truncated}}

\textbf{GS1 DataBar Truncated} is a fixed-length, linear barcode
symbology that can be used to encode a GTIN-14 for in-house
applications.

Also known as: RSS-14 Truncated.

Variants:

\begin{itemize}
\tightlist
\item
  \protect\hyperlink{gs1-databar-stacked}{GS1 DataBar Stacked} is a
  variant of GS1 DataBar Truncated for use where a taller, narrower
  symbol is required.
\item
  \protect\hyperlink{gs1-composite-symbols}{GS1 DataBar Truncated
  Composite} is a variant of GS1 DataBar Truncated that should be used
  when a CC-A or CC-B GS1 composite 2D component is required.
\end{itemize}

Standards: ISO/IEC 24724, ITS Reduced Space Symbology (RSS), AIM ISS -
Reduced Space Symbology (RSS), GS1 General Specifications.

\hypertarget{data-and-options-38}{%
\paragraph{Data and Options}\label{data-and-options-38}}

\begin{itemize}
\tightlist
\item
  The data field input must be a solitary \emph{AI (01)} with 13 or 14
  digits of a GTIN, i.e. \texttt{(01)...}.
\item
  If thirteen digits of \emph{AI (01)} data are supplied then the check
  digit is calculated automatically, otherwise the provided check digit
  will be verified.
\item
  \emph{Deprecated: For internal use.} The \textbf{linkage} option
  signifies the presence of a GS1 composite 2D component.
\end{itemize}

\hypertarget{examples-28}{%
\paragraph{Examples}\label{examples-28}}

Identical symbols, input provided with and without a check digit:

\begin{verbatim}
Data:    (01)24012345678905
Options: 
Encoder: databartruncated
\end{verbatim}

\begin{verbatim}
Data:    (01)2401234567890
Options: 
Encoder: databartruncated
\end{verbatim}

\includegraphics{images/databartruncated-1.eps}

\hypertarget{gs1-databar-stacked}{%
\subsubsection{GS1 DataBar Stacked}\label{gs1-databar-stacked}}

\textbf{GS1 DataBar Stacked} is a fixed-length, stacked linear barcode
symbology that can be used to encode a GTIN-14 for in-house
applications.

Also known as: RSS-14 Stacked.

Variants:

\begin{itemize}
\tightlist
\item
  \protect\hyperlink{gs1-databar-truncated}{GS1 DataBar Truncated} is a
  variant of GS1 DataBar Stacked for use where a shorter, wider symbol
  is required.
\item
  \protect\hyperlink{gs1-composite-symbols}{GS1 DataBar Stacked
  Composite} is a variant of GS1 DataBar Stacked that should be used
  when a CC-A or CC-B GS1 composite 2D component is required.
\end{itemize}

Standards: ISO/IEC 24724, ITS Reduced Space Symbology (RSS), AIM ISS -
Reduced Space Symbology (RSS), GS1 General Specifications.

\hypertarget{data-and-options-39}{%
\paragraph{Data and Options}\label{data-and-options-39}}

\begin{itemize}
\tightlist
\item
  The data field input must be a solitary \emph{AI (01)} with thirteen
  or fourteen digits of a GTIN, i.e. \texttt{(01)...}.
\item
  If thirteen digits of \emph{AI (01)} data are supplied then the check
  digit is calculated automatically, otherwise the provided check digit
  will be verified.
\item
  \emph{Deprecated: For internal use.} The \textbf{linkage} option
  signifies the presence of a GS1 composite 2D component.
\end{itemize}

\hypertarget{examples-29}{%
\paragraph{Examples}\label{examples-29}}

Identical symbols, input provided with and without a check digit:

\begin{verbatim}
Data:    (01)24012345678905
Options: 
Encoder: databarstacked
\end{verbatim}

\begin{verbatim}
Data:    (01)2401234567890
Options: 
Encoder: databarstacked
\end{verbatim}

\includegraphics{images/databarstacked-1.eps}

\hypertarget{gs1-databar-limited}{%
\subsubsection{GS1 DataBar Limited}\label{gs1-databar-limited}}

\textbf{GS1 DataBar Limited} is fixed-length, linear barcode symbology
that can be used to encode a GTIN-14 beginning with \emph{0} or \emph{1}
for in-house applications.

Also known as: RSS Limited.

Variants:

\begin{itemize}
\tightlist
\item
  \protect\hyperlink{gs1-composite-symbols}{GS1 DataBar Limited
  Composite} is a variant of GS1 DataBar Limited that should be used
  when a CC-A or CC-B GS1 composite 2D component is required.
\end{itemize}

Standards: ISO/IEC 24724, ITS Reduced Space Symbology (RSS), AIM ISS -
Reduced Space Symbology (RSS), GS1 General Specifications.

\hypertarget{data-and-options-40}{%
\paragraph{Data and Options}\label{data-and-options-40}}

\begin{itemize}
\tightlist
\item
  The data field input must be a solitary \emph{AI (01)} with thirteen
  or fourteen digits of a GTIN starting with \emph{0} or \emph{1}, i.e.
  \texttt{(01)0...} or \texttt{(01)1...}.
\item
  If thirteen digits of \emph{AI (01)} data are supplied then the check
  digit is calculated automatically, otherwise the provided check digit
  will be verified.
\item
  \emph{Deprecated: For internal use.} The \textbf{linkage} option
  signifies the presence of a GS1 composite 2D component.
\end{itemize}

\hypertarget{examples-30}{%
\paragraph{Examples}\label{examples-30}}

Identical symbols, input provided with and without a check digit:

\begin{verbatim}
Data:    (01)15012345678907
Options: 
Encoder: databarlimited
\end{verbatim}

\begin{verbatim}
Data:    (01)1501234567890
Options: 
Encoder: databarlimited
\end{verbatim}

\includegraphics{images/databarlimited-1.eps}

\hypertarget{gs1-north-american-coupon-code}{%
\subsubsection{GS1 North American Coupon
Code}\label{gs1-north-american-coupon-code}}

\textbf{GS1 North American Coupon Code} is an implementation of the
\protect\hyperlink{gs1-databar-expanded}{GS1 DataBar Expanded} barcode
symbology with AI (8110) that is used as a paperless coupon system at
point of sale.

Also known as: GS1 DataBar Coupon, U.S. Coupon Code.

Standards: North American Coupon Application Guideline Using GS1 DataBar
Expanded Symbols, ISO/IEC 24724, GS1 General Specifications.

\hypertarget{data-and-options-41}{%
\paragraph{Data and Options}\label{data-and-options-41}}

\begin{itemize}
\tightlist
\item
  The data field input is provided in
  \protect\hyperlink{gs1-application-identifier-standard-format}{GS1
  Application Identifier standard format} and must be a solitary AI
  (8110).
\item
  The \textbf{segments} option is used to specify the maximum number of
  segments per row which must be an even number. The default is
  \emph{4}.
\end{itemize}

\hypertarget{examples-31}{%
\paragraph{Examples}\label{examples-31}}

\begin{verbatim}
Data:    (8110)106141416543213500110000310123196000
Options: includetext segments=8
Encoder: gs1northamericancoupon
\end{verbatim}

\includegraphics{images/gs1northamericancoupon-1.eps}

\hypertarget{postal-symbols}{%
\subsection{Postal Symbols}\label{postal-symbols}}

\hypertarget{australia-post-4-state-customer-code}{%
\subsubsection{Australia Post 4 State Customer
Code}\label{australia-post-4-state-customer-code}}

The \textbf{Australia Post 4 State Customer Code} is a barcode used by
the Australian Postal Service to encode the data on letter mail.

\hypertarget{data-and-options-42}{%
\paragraph{Data and Options}\label{data-and-options-42}}

\begin{itemize}
\tightlist
\item
  The first two characters of the data field are digits used to specify
  the manditory FCC type of the symbols, either \texttt{11},
  \texttt{45}, \texttt{59} or \texttt{67}.
\item
  The next eight characters are digits that specify the manditory DPID.
\item
  The number of remaining characters varies according to the given FCC
  code and these specify the contents of the customer information field
  in one of two alphabets:

  \begin{itemize}
  \tightlist
  \item
    The \textbf{custinfoenc} option should be supplied as
    \texttt{custinfoenc=numeric} if the customer information field is to
    be encoded using the numeric alphabet which can contain the digits
    0-9.
  \item
    Otherwise the customer information field is encoded using the
    default character encoding, \texttt{custinfoenc=character}, which
    permits any of the following characters:

    \begin{itemize}
    \tightlist
    \item
      Upper case letters \texttt{A-Z}
    \item
      Lower case letters \texttt{a-z}
    \item
      Digits \texttt{0-9}
    \item
      Symbols \texttt{space} and \texttt{\#}
    \end{itemize}
  \end{itemize}
\end{itemize}

\hypertarget{examples-32}{%
\paragraph{Examples}\label{examples-32}}

FCC 62 symbol with character customer data:

\begin{verbatim}
Data:    6279438541AaaB 155
Options: custinfoenc=character
Encoder: auspost
\end{verbatim}

\includegraphics{images/auspost-1.eps}

FCC 59 symbol with numeric customer data:

\begin{verbatim}
Data:    593221132401234567
Options: custinfoenc=numeric
Encoder: auspost
\end{verbatim}

\includegraphics{images/auspost-2.eps}

\hypertarget{deutsche-post-identcode}{%
\subsubsection{Deutsche Post Identcode}\label{deutsche-post-identcode}}

\textbf{Deutsche Post Identcode} is an implementation of the
\protect\hyperlink{interleaved-2-of-5}{Interleaved 2 of 5} barcode
symbology that is used by German Post for mail routing.

Also known as: DHL Identcode.

\hypertarget{data-and-options-43}{%
\paragraph{Data and Options}\label{data-and-options-43}}

\begin{itemize}
\tightlist
\item
  The data consists of a consecutive string of eleven or twelve digits
  consisting of:

  \begin{itemize}
  \tightlist
  \item
    Two-digit primary distribution centre identifier
  \item
    Three-digit customer identifier
  \item
    Six-digit mail piece identifier
  \item
    One-digit check digit (may be omitted)
  \end{itemize}
\item
  If eleven digits are supplied then the check digit is calculated
  automatically. Otherwise the provided check digit will be verified.
\end{itemize}

\hypertarget{example-8}{%
\paragraph{Example}\label{example-8}}

Identical symbols, input provided with an without a check digit:

\begin{verbatim}
Data:    563102430313
Options: includetext
Encoder: identcode
\end{verbatim}

\begin{verbatim}
Data:    56310243031
Options: includetext
Encoder: identcode
\end{verbatim}

\includegraphics{images/identcode-1.eps}

\hypertarget{deutsche-post-leitcode}{%
\subsubsection{Deutsche Post Leitcode}\label{deutsche-post-leitcode}}

The \textbf{Deutsche Post Leitcode} barcode symbology is an
implementation of the \protect\hyperlink{interleaved-2-of-5}{Interleaved
2 of 5} barcode that is used by German Post for mail routing.

Also known as: DHL Leitcode.

\hypertarget{data-and-options-44}{%
\paragraph{Data and Options}\label{data-and-options-44}}

\begin{itemize}
\tightlist
\item
  The data consists of a consecutive string of thirteen or fourteen
  digits consisting of:

  \begin{itemize}
  \tightlist
  \item
    Five-digit postal code
  \item
    Three-digit street identifier
  \item
    Three-digit house number
  \item
    Two-digit product code
  \item
    One-digit check digit (may be omitted)
  \end{itemize}
\item
  If thirteen digits are supplied then the check digit is calculated
  automatically. Otherwise the provided check digit will be verified.
\end{itemize}

\hypertarget{examples-33}{%
\paragraph{Examples}\label{examples-33}}

Identical symbols, input provided with and without a check digit:

\begin{verbatim}
Data:    21348075016401
Options: includetext
Encoder: leitcode
\end{verbatim}

\begin{verbatim}
Data:    2134807501640
Options: includetext
Encoder: leitcode
\end{verbatim}

\includegraphics{images/leitcode-1.eps}

\hypertarget{japan-post-4-state-barcode}{%
\subsubsection{Japan Post 4 State
Barcode}\label{japan-post-4-state-barcode}}

The \textbf{Japan Post 4 state barcode} symbology is used by the Japan
Post service to encode the delivery point identifier on letter mail.

\hypertarget{data-and-options-45}{%
\paragraph{Data and Options}\label{data-and-options-45}}

\begin{itemize}
\tightlist
\item
  The data may contain any of the following characters:

  \begin{itemize}
  \tightlist
  \item
    Capital letters \texttt{A-Z}
  \item
    Digits \texttt{0-9}
  \item
    Hyphen \texttt{-}
  \end{itemize}
\end{itemize}

\hypertarget{example-9}{%
\paragraph{Example}\label{example-9}}

\begin{verbatim}
Data:    6540123789-A-K-Z
Options: 
Encoder: japanpost
\end{verbatim}

\includegraphics{images/japanpost-1.eps}

\hypertarget{maxicode}{%
\subsubsection{MaxiCode}\label{maxicode}}

The \textbf{MaxiCode} barcode symbology is a 2D barcode based on a
hexagonal matrix surrounding a bulls eye pattern. It can encode a
structured carrier message and full 256 character extended-ASCII.

Also known as: UPS Code, Code 6, Dense Code.

Standards: ISO/IEC 16023, ANSI/AIM BC10 - ISS MaxiCode.

\hypertarget{data-and-options-46}{%
\paragraph{Data and Options}\label{data-and-options-46}}

\begin{itemize}
\tightlist
\item
  The \textbf{mode} option is used to specify how the data is structured
  in the symbol:

  \begin{itemize}
  \tightlist
  \item
    \texttt{mode=2} - Formatted data containing a Structured Carrier
    Message with a numeric (US domestic) postal code.
  \item
    \texttt{mode=3} - Formatted data containing a Structured Carrier
    Message with an alphanumeric (international) postal code.
  \item
    \texttt{mode=4} - Unstructured extended ASCII data using standard
    error correction.
  \item
    \texttt{mode=5} - Unstructured extended ASCII data using enhanced
    error correction.
  \item
    \texttt{mode=6} - Barcode reader programming.
  \end{itemize}
\item
  If \textbf{mode} is unspecified the encoder will default to selecting
  \texttt{mode=5} if the encoded length of the input data permits
  enhanced error correction, otherwise it will select \texttt{mode=4}
  which provides standard error correction.
\item
  The default interpretation of data by readers is in accordance with
  ISO/IEC 8859-1. When supported by the receiver characters from other
  code pages can be encoded using Extended Channel Interpretation (ECI).
\item
  When the \textbf{parse} option is specified, any instances of
  \texttt{\^{}NNN} in the data field are replaced with their equivalent
  ASCII value, useful for specifying unprintable characters.
\item
  When the \textbf{parsefnc} option is specified, non-data function
  characters can be specified by escape sequences:

  \begin{itemize}
  \tightlist
  \item
    \texttt{\^{}ECI000000} to \texttt{\^{}ECI999999}: ECI indicators
  \end{itemize}
\item
  If \texttt{mode=4}, \texttt{mode=5} or \texttt{mode=6} the data field
  may contain any extended ASCII data.
\item
  If \texttt{mode=2} or \texttt{mode=3} the data field must begin with a
  properly structured carrier message, followed by any extended ASCII
  data.
\item
  The structured carrier message contains a postal code, three-digit
  class of service and a three-digit ISO country code separated by
  \emph{GS} (ASCII 29) characters. It is formatted in the data field as
  follows:
  \texttt{{[}postal\ code{]}\^{}029{[}country\ code{]}\^{}029{[}service\ class{]}\^{}029}.
  If \texttt{mode=2} the postcode must be numeric, whilst if
  \texttt{mode=3} the postcode may contain up to six digits, upper case
  letters and spaces.
\item
  Alternatively, messages may begin with the special application field
  identifier \texttt{{[})\textgreater{}\{RS\}01\{GS\}yy} where
  \texttt{\{RS\}} represents ASCII value 30, \texttt{\{GS\}} represents
  ASCII value 29 and \texttt{yy} is a two-digit year. In parse mode this
  is represented as \texttt{{[})\textgreater{}\^{}03001\^{}0299}. If
  \texttt{mode=2} or \texttt{mode=3} this must be immediately followed
  by the structured carrier message.
\end{itemize}

\hypertarget{examples-34}{%
\paragraph{Examples}\label{examples-34}}

\begin{verbatim}
Data:    This is MaxiCode
Options: 
Encoder: maxicode
\end{verbatim}

\begin{verbatim}
Data:    This is Maxi^067ode
Options: parse
Encoder: maxicode
\end{verbatim}

\includegraphics{images/maxicode-1.eps}

\begin{verbatim}
Data:    152382802^029840^029001^0291Z00004951^029UPSN^02906X610^029159^0291234567^0291/1^029^029Y^029634 ALPHA DR^029PITTSBURGH^029PA^029^004
Options: mode=2 parse
Encoder: maxicode
\end{verbatim}

\includegraphics{images/maxicode-2.eps}

\begin{verbatim}
Data:    ABC123^029840^029001^0291Z00004951^029UPSN^02906X610^029159^0291234567^0291/1^029^029Y^029634 ALPHA DR^029PITTSBURGH^029PA^029^004
Options: mode=3 parse
Encoder: maxicode
\end{verbatim}

\includegraphics{images/maxicode-3.eps}

\begin{verbatim}
Data:    [\)>^03001^02996152382802^029840^029001^0291Z00004951^029UPSN^02906X610^029159^0291234567^0291/1^029^029Y^029634 ALPHA DR^029PITTSBURGH^029PA^029^004
Options: mode=2 parse
Encoder: maxicode
\end{verbatim}

\includegraphics{images/maxicode-4.eps}

\hypertarget{royal-mail-4-state-customer-code}{%
\subsubsection{Royal Mail 4 State Customer
Code}\label{royal-mail-4-state-customer-code}}

The \textbf{Royal Mail 4 State Customer Code} is a barcode symbology
used by the British Postal Service to encode the postcode and delivery
point identifier on letter mail.

Also known as: RM4SCC, CBC, BPO 4 State Code

\hypertarget{data-and-options-47}{%
\paragraph{Data and Options}\label{data-and-options-47}}

\begin{itemize}
\tightlist
\item
  The data may contain any of the following characters from the postcode
  and DPID:

  \begin{itemize}
  \tightlist
  \item
    Capital letters A-Z
  \item
    Digits 0-9
  \end{itemize}
\item
  The maditory checksum digit is calculated automatically and must not
  be included in the data field
\end{itemize}

\hypertarget{example-10}{%
\paragraph{Example}\label{example-10}}

\begin{verbatim}
Data:    LE28HS9Z
Options: includetext
Encoder: royalmail
\end{verbatim}

\includegraphics{images/rm4scc-1.eps}

\hypertarget{royal-mail-mailmark}{%
\subsubsection{Royal Mail Mailmark}\label{royal-mail-mailmark}}

\textbf{Royal Mail Mailmark} is an implementation of the
\protect\hyperlink{data-matrix}{Data Matrix} (ECC 200) barcode symbology
for application within the Royal Mail postal system.

Standards: Royal Mail Mailmark Barcode Definition Document

\hypertarget{data-and-options-48}{%
\paragraph{Data and Options}\label{data-and-options-48}}

\begin{itemize}
\tightlist
\item
  The data field input consists of 45 characters of Mailmark structured
  data (including required space padding) followed by variable-length,
  free-formatted customer data.
\item
  Note: Due to the fixed-length field format of the message data it is
  important to carefully preserve spaces in the data field.
\item
  The mandatory \textbf{type} option is used to specify the size of the
  symbol, either \texttt{7}, \texttt{9} or \texttt{29}.
\end{itemize}

\hypertarget{example-11}{%
\paragraph{Example}\label{example-11}}

\begin{verbatim}
Data:    JGB 012100123412345678AB19XY1A 0             www.xyz.com
Options: type=29
Encoder: mailmark
\end{verbatim}

\hypertarget{royal-tnt-post-4-state-barcode}{%
\subsubsection{Royal TNT Post 4 state
barcode}\label{royal-tnt-post-4-state-barcode}}

The \textbf{Royal TNT Post 4 state barcode} symbology is used by the
Dutch Postal Service to encode the delivery point identifier on letter
mail.

Also known as: KIX, Klant IndeX.

\hypertarget{data-and-options-49}{%
\paragraph{Data and Options}\label{data-and-options-49}}

\begin{itemize}
\tightlist
\item
  The data may contain any of the following characters from the DPID:

  \begin{itemize}
  \tightlist
  \item
    Capital letters \texttt{A-Z}
  \item
    Digits \texttt{0-9}
  \end{itemize}
\end{itemize}

\hypertarget{example-12}{%
\paragraph{Example}\label{example-12}}

\begin{verbatim}
Data:    1231FZ13XHS
Options: includetext
Encoder: kix
\end{verbatim}

\includegraphics{images/kix-1.eps}

\hypertarget{usps-intelligent-mail}{%
\subsubsection{USPS Intelligent Mail}\label{usps-intelligent-mail}}

The \textbf{USPS Intelligent Mail} barcode is used by the US Postal
service to encode the delivery and sender information on letter mail.

Also known as: USPS OneCode.

Standards: USPS-STD-11.

\hypertarget{data-and-options-50}{%
\paragraph{Data and Options}\label{data-and-options-50}}

\begin{itemize}
\tightlist
\item
  The data contains 31 digits representing the following:

  \begin{itemize}
  \tightlist
  \item
    Barcode Identifier - two digits
  \item
    Service Type Identifier - three digits
  \item
    Mailer ID, Sequence Number - either six then nine digits
    respectively or nine then six digits respectively
  \item
    Delivery Point ZIP Code - eleven digits
  \end{itemize}
\item
  The mandatory checksum digit is calculated automatically and must not
  be included in the data field.
\end{itemize}

\hypertarget{example-13}{%
\paragraph{Example}\label{example-13}}

\begin{verbatim}
Data:    0123456709498765432101234567891
Options: includetext
Encoder: onecode
\end{verbatim}

\includegraphics{images/onecode-1.eps}

\hypertarget{usps-postnet}{%
\subsubsection{USPS POSTNET}\label{usps-postnet}}

The \textbf{USPS POSTNET} barcode symbology is used by the US Postal
service to encode the ZIP code information on letter mail.

\hypertarget{data-and-options-51}{%
\paragraph{Data and Options}\label{data-and-options-51}}

\begin{itemize}
\tightlist
\item
  The data field contains the digits from the ZIP code, without dashes.
\item
  The mandatory checksum is calculated automatically and must not be
  included in the data field.
\end{itemize}

\hypertarget{example-14}{%
\paragraph{Example}\label{example-14}}

\begin{verbatim}
Data:    12345123412
Options: 
Encoder: postnet
\end{verbatim}

\includegraphics{images/postnet-1.eps}

\hypertarget{usps-planet}{%
\subsubsection{USPS PLANET}\label{usps-planet}}

The \textbf{USPS PLANET} barcode symbology is used by the US Postal
service to encode the ZIP code information on letter mail.

\hypertarget{data-and-options-52}{%
\paragraph{Data and Options}\label{data-and-options-52}}

\begin{itemize}
\tightlist
\item
  The data field contains eleven or thirteen digits, without dashes.
\item
  The mandatory checksum is calculated automatically and must not be
  included in the data field.
\end{itemize}

\hypertarget{example-15}{%
\paragraph{Example}\label{example-15}}

\begin{verbatim}
Data:   01234567890
Options: 
Encoder: planet
\end{verbatim}

\includegraphics{images/planet-1.eps}

\hypertarget{usps-fim-symbols}{%
\subsubsection{USPS FIM Symbols}\label{usps-fim-symbols}}

The \textbf{USPS FIM} encoder is used to generate static predefined
barcode symbols.

\hypertarget{data-and-options-53}{%
\paragraph{Data and Options}\label{data-and-options-53}}

\begin{itemize}
\tightlist
\item
  The data field accepts one of the following values:

  \begin{itemize}
  \tightlist
  \item
    \texttt{fima} - US Postal Service FIM-A symbol
  \item
    \texttt{fimb} - US Postal Service FIM-B symbol
  \item
    \texttt{fimc} - US Postal Service FIM-C symbol
  \item
    \texttt{fimd} - US Postal Service FIM-D symbol
  \end{itemize}
\end{itemize}

\hypertarget{examples-35}{%
\paragraph{Examples}\label{examples-35}}

A USPS FIM A symbol:

\begin{verbatim}
Data:    fima
Options: 
Encoder: symbol
\end{verbatim}

\includegraphics{images/fima-1.eps}

A USPS FIM B symbol:

\begin{verbatim}
Data:    fimb
Options: 
Encoder: symbol
\end{verbatim}

\includegraphics{images/fimb-1.eps}

A USPS FIM C symbol:

\begin{verbatim}
Data:    fimc
Options: 
Encoder: symbol
\end{verbatim}

\includegraphics{images/fimc-1.eps}

A USPS FIM D symbol:

\begin{verbatim}
Data:    fimd
Options: 
Encoder: symbol
\end{verbatim}

\includegraphics{images/fimd-1.eps}

\hypertarget{pharmaceutical-symbols}{%
\subsection{Pharmaceutical Symbols}\label{pharmaceutical-symbols}}

\hypertarget{italian-pharmacode}{%
\subsubsection{Italian Pharmacode}\label{italian-pharmacode}}

\textbf{Italian Pharmacode} is a discrete, fixed length, self-checking
barcode symbology used for pharmaceutical products in Italy.

Also known as: Code 32, IMH, Radix 32.

\hypertarget{data-and-options-54}{%
\paragraph{Data and Options}\label{data-and-options-54}}

\begin{itemize}
\tightlist
\item
  The data field must contain either eight or nine digits from the code.
  The leading \emph{A} which is provided in some applications must be
  omitted.
\item
  The mandatory check digit is calculated automatically if it is not
  provided, otherwise the provided check digit is verified.
\end{itemize}

\hypertarget{examples-36}{%
\paragraph{Examples}\label{examples-36}}

Identical symbols, input provided with and without a check digit:

\begin{verbatim}
Data:    012345676
Options: includetext
Encoder: code32
\end{verbatim}

\begin{verbatim}
Data:    01234567
Options: includetext
Encoder: code32
\end{verbatim}

\includegraphics{images/code32-1.eps}

\hypertarget{pharmacode}{%
\subsubsection{Pharmacode}\label{pharmacode}}

\textbf{Pharmacode} is a binary barcode symbology that is used by the
Pharmaceutical industry.

Also known as: Pharmaceutical Binary Code. Laetus Code.

Variants:

\begin{itemize}
\tightlist
\item
  \protect\hyperlink{two-track-pharmacode}{Two-track Pharmacode} is a
  variant of the Pharmacode barcode.
\end{itemize}

\hypertarget{data-and-options-55}{%
\paragraph{Data and Options}\label{data-and-options-55}}

\begin{itemize}
\tightlist
\item
  The data field must contain a number between \texttt{3} and
  \texttt{131070} inclusive.
\item
  The \textbf{nwidth}, \textbf{wwidth} and \textbf{swidth} options can
  be used to specify a custom width (in points) for the narrow bars,
  wide bars and inter-bar spaces respectively.
\end{itemize}

\hypertarget{example-16}{%
\paragraph{Example}\label{example-16}}

\begin{verbatim}
Data:    117480
Options: 
Encoder: pharmacode
\end{verbatim}

\includegraphics{images/pharmacode-1.eps}

\hypertarget{two-track-pharmacode}{%
\subsubsection{Two-Track Pharmacode}\label{two-track-pharmacode}}

\textbf{Two-Track Pharmacode} is a binary barcode symbology used by the
Pharmaceutical industry.

Also known as: Two-track Pharmaceutical Binary Code, Two-track Laetus
Code.

Variants:

\begin{itemize}
\tightlist
\item
  \protect\hyperlink{pharmacode}{Pharmacode} is a variant of the
  Two-track Pharmacode barcode.
\end{itemize}

\hypertarget{data-and-options-56}{%
\paragraph{Data and Options}\label{data-and-options-56}}

\begin{itemize}
\tightlist
\item
  The data field must contain a number between \texttt{4} and
  \texttt{64570080} inclusive.
\end{itemize}

\hypertarget{example-17}{%
\paragraph{Example}\label{example-17}}

\begin{verbatim}
Data:    117480
Options: 
Encoder: pharmacode2
\end{verbatim}

\includegraphics{images/pharmacode2-1.eps}

\hypertarget{pzn}{%
\subsubsection{PZN}\label{pzn}}

\textbf{PZN} is a discrete, fixed length, self-checking barcode
symbology used for pharmaceutical products in Germany.

Also known as: Pharmazentralnummer.

Variants:

\begin{itemize}
\tightlist
\item
  PZN7 is the seven-digit format.
\item
  PZN8 is the eight-digit format.
\end{itemize}

\hypertarget{data-and-options-57}{%
\paragraph{Data and Options}\label{data-and-options-57}}

\begin{itemize}
\tightlist
\item
  For the default PZN7 encoding, the data field must contain six digits
  or seven digits.
\item
  The \textbf{pzn8} option specifies that a PZN8 symbol is required, in
  which case the data field must contain seven digits or eight digits.
\item
  The mandatory check digit is calculated automatically if not provided,
  otherwise the provided check digit is verified.
\item
  Note: by definition, not all six-digit or seven-digit number sequences
  are valid inputs.
\end{itemize}

\hypertarget{examples-37}{%
\paragraph{Examples}\label{examples-37}}

Identical PZN7 symbols, input provided with and without a check digit:

\begin{verbatim}
Data:    1234562
Options: includetext
Encoder: pzn
\end{verbatim}

\begin{verbatim}
Data:    123456
Options: includetext
Encoder: pzn
\end{verbatim}

\includegraphics{images/pzn-1.eps}

Identical PZN8 symbols, input provided with and without a check digit:

\begin{verbatim}
Data:    0275808
Options: pzn8 includetext
Encoder: pzn
\end{verbatim}

\begin{verbatim}
Data:    02758089
Options: pzn8 includetext
Encoder: pzn
\end{verbatim}

\includegraphics{images/pzn-2.eps}

\hypertarget{hibc-symbols}{%
\subsection{HIBC Symbols}\label{hibc-symbols}}

\textbf{HIBC barcodes} use a number of general symbologies as carrier
symbols for data structured according to the LIC and PAS structured data
definitions.

Variants:

\begin{itemize}
\tightlist
\item
  \protect\hyperlink{hibc-code-39}{HIBC Code 39} is a variant of
  \protect\hyperlink{code-39}{Code 39}.
\item
  \protect\hyperlink{hibc-code-128}{HIBC Code 128} is a variant of
  \protect\hyperlink{code-128}{Code 128}.
\item
  \protect\hyperlink{hibc-pdf417}{HIBC PDF417} is a variant of
  \protect\hyperlink{pdf417}{PDF417}.
\item
  \protect\hyperlink{hibc-micropdf417}{HIBC MicroPDF417} is a variant of
  \protect\hyperlink{micropdf417}{MicroPDF417}.
\item
  \protect\hyperlink{hibc-qr-code}{HIBC QR Code} is a variant of
  \protect\hyperlink{qr-code}{QR Code}.
\item
  \protect\hyperlink{hibc-data-matrix}{HIBC Data Matrix} is a variant of
  \protect\hyperlink{data-matrix}{Data Matrix}.
\item
  \protect\hyperlink{hibc-codablock-f}{HIBC Codablock F} is a variant of
  \protect\hyperlink{codablock-f}{Codablock F}.
\end{itemize}

Standards: ANSI/HIBC Provider Applications Standard, ANSI/HIBC Supplier
Labelling Standard, ANSI/HIBC Positive Identification for Patient
Safety, ANSI/HIBC Syntax Standard.

\hypertarget{data-and-options-58}{%
\paragraph{Data and Options}\label{data-and-options-58}}

\begin{itemize}
\tightlist
\item
  The data should be pre-encoded to describe the intended barcode
  content.
\item
  The HIBC \emph{+} character is prefixed automatically.
\item
  The mandatory HIBC check character is automatically appended to the
  input.
\end{itemize}

\hypertarget{hibc-code-39}{%
\subsubsection{HIBC Code 39}\label{hibc-code-39}}

\begin{verbatim}
Data:    A123BJC5D6E71
Options: includetext
Encoder: hibccode39
\end{verbatim}

\includegraphics{images/hibccode39-1.eps}

\hypertarget{hibc-code-128}{%
\subsubsection{HIBC Code 128}\label{hibc-code-128}}

\begin{verbatim}
Data:    A123BJC5D6E71
Options: includetext
Encoder: hibccode128
\end{verbatim}

\includegraphics{images/hibccode128-1.eps}

\hypertarget{hibc-pdf417}{%
\subsubsection{HIBC PDF417}\label{hibc-pdf417}}

\begin{verbatim}
Data:    A123BJC5D6E71
Options: 
Encoder: hibcpdf417
\end{verbatim}

\includegraphics{images/hibcpdf417-1.eps}

\hypertarget{hibc-micropdf417}{%
\subsubsection{HIBC MicroPDF417}\label{hibc-micropdf417}}

\begin{verbatim}
Data:    A123BJC5D6E71
Options: 
Encoder: hibcmicropdf417
\end{verbatim}

\includegraphics{images/hibcmicropdf417-1.eps}

\hypertarget{hibc-qr-code}{%
\subsubsection{HIBC QR Code}\label{hibc-qr-code}}

\begin{verbatim}
Data:    A123BJC5D6E71
Options: 
Encoder: hibcqrcode
\end{verbatim}

\includegraphics{images/hibcqrcode-1.eps}

\hypertarget{hibc-data-matrix}{%
\subsubsection{HIBC Data Matrix}\label{hibc-data-matrix}}

\begin{verbatim}
Data:    A123BJC5D6E71
Options: 
Encoder: hibcdatamatrix
\end{verbatim}

\includegraphics{images/hibcdatamatrix-1.eps}

\hypertarget{hibc-codablock-f}{%
\subsubsection{HIBC Codablock F}\label{hibc-codablock-f}}

\begin{verbatim}
Data:    A123BJC5D6E71
Options: 
Encoder: hibccodablockf
\end{verbatim}

\includegraphics{images/hibccodablockf-1.eps}

\hypertarget{less-used-symbols}{%
\subsection{Less-used Symbols}\label{less-used-symbols}}

\hypertarget{bc412}{%
\subsubsection{BC412}\label{bc412}}

The \textbf{BC412} barcode symbology is single width, variable length
barcode that is used for silicon wafer identification by the
semiconductor manufacturing industry.

Also known as: BC412 SEMI, BC412 IBM.

\hypertarget{data-and-options-59}{%
\paragraph{Data and Options}\label{data-and-options-59}}

\begin{itemize}
\tightlist
\item
  The data field can hold any of the following:

  \begin{itemize}
  \tightlist
  \item
    Numbers 0-9
  \item
    Capital letters A-Z, excluding O
  \end{itemize}
\item
  The \textbf{includestartstop} option enables the display of start and
  stop bars.
\item
  The \textbf{includecheck} option calculates the check character.
\item
  The \textbf{includecheckintext} option makes the calculated checksum
  appear in the human readable text.
\item
  The \textbf{semi} option enables conformance to the SEMI standard by
  enabling start and stop bars as well as a check character.
\item
  The \textbf{inkspread} option can be used to adjust the width of the
  bars.
\end{itemize}

\hypertarget{examples-38}{%
\paragraph{Examples}\label{examples-38}}

\begin{verbatim}
Data:    BC412
Options: includecheck
Encoder: bc412
\end{verbatim}

\includegraphics{images/bc412-1.eps}

\begin{verbatim}
Data:    BC412
Options: includestartstop
Encoder: bc412
\end{verbatim}

\includegraphics{images/bc412-2.eps}

\begin{verbatim}
Data:    BC412
Options: semi
Encoder: bc412
\end{verbatim}

\includegraphics{images/bc412-3.eps}

\hypertarget{channel-code}{%
\subsubsection{Channel Code}\label{channel-code}}

\textbf{Channel Code} is a linear, continuous, self-checking,
bidirectional barcode symbology that encodes between two and seven
digits in a short space.

Standards: ANSI/AIM BC12 - USS Channel Code.

\hypertarget{data-and-options-60}{%
\paragraph{Data and Options}\label{data-and-options-60}}

\begin{itemize}
\tightlist
\item
  The data field can hold zero prefixed values from any of the following
  ranges:

  \begin{itemize}
  \tightlist
  \item
    Channel 3: \texttt{00-26}
  \item
    Channel 4: \texttt{000-292}
  \item
    Channel 5: \texttt{0000-3493}
  \item
    Channel 6: \texttt{00000-44072}
  \item
    Channel 7: \texttt{000000-576688}
  \item
    Channel 8: \texttt{0000000-7742862}
  \end{itemize}
\item
  The channel is determined to be one more than the number of digits
  given in the data field.
\item
  The \textbf{shortfinder} option generates a symbol with a shortened
  finder pattern.
\item
  The \textbf{includecheck} option appends an optional check bar suffix.
\end{itemize}

\hypertarget{examples-39}{%
\paragraph{Examples}\label{examples-39}}

A channel 3 symbol holding the value five:

\begin{verbatim}
Data:    05
Options: includetext
Encoder: channelcode
\end{verbatim}

\includegraphics{images/channelcode-1.eps}

A channel 4 symbol holding the value 123:

\begin{verbatim}
Data:    123
Options: includetext
Encoder: channelcode
\end{verbatim}

\includegraphics{images/channelcode-2.eps}

A channel 4 symbol holding the value five including optional check bars:

\begin{verbatim}
Data:    005
Options: includetext includecheck
Encoder: channelcode
\end{verbatim}

\includegraphics{images/channelcode-3.eps}

A channel 3 symbol holding the value 26 with a shorteded finder pattern:

\begin{verbatim}
Data:    26
Options: shortfinder includetext
Encoder: channelcode
\end{verbatim}

\includegraphics{images/channelcode-4.eps}

\hypertarget{codabar}{%
\subsubsection{Codabar}\label{codabar}}

\textbf{Codabar} is a linear, discrete, self-checking, bidirectional
barcode symbology that can encode digits, six symbols and four delimiter
characters. It is primarily used by libraries and blood banks, photo
labs and FedEx airbills.

Also known as: Rationalized Codabar, Ames Code, NW-7, USD-4,
USS-Codabar, ABC Codabar, Monarch, Code 2 of 7.

Standards: ANSI/AIM BC3 - USS Codabar, BS EN 798.

\hypertarget{data-and-options-61}{%
\paragraph{Data and Options}\label{data-and-options-61}}

\begin{itemize}
\tightlist
\item
  The data field must start and stop with one of the following
  delimiters

  \begin{itemize}
  \tightlist
  \item
    \texttt{ABCD}
  \item
    \texttt{TNE*} (with the \emph{altstartstop} option)
  \end{itemize}
\item
  The data field can otherwise hold any of the following

  \begin{itemize}
  \tightlist
  \item
    Digits 0-9
  \item
    Symbols \texttt{-\$:/.+}
  \end{itemize}
\item
  The \textbf{altstartstop} option specifies that the alternative set of
  delimiter characters is in use.
\item
  The \textbf{includecheck} option calculates the check digit.
\item
  The \textbf{includecheckintext} option makes the calculated check
  characters appear in the human readable text.
\end{itemize}

\hypertarget{example-18}{%
\paragraph{Example}\label{example-18}}

\begin{verbatim}
Data:    A0123456789B
Options: includecheck includetext includecheckintext
Encoder: rationalizedCodabar
\end{verbatim}

\includegraphics{images/codabar-1.eps}

\hypertarget{codablock-f}{%
\subsubsection{Codablock F}\label{codablock-f}}

The \textbf{Codablock F} barcode symbology is 2D stacked-linear barcode
that consists of a number of stacked \protect\hyperlink{code-128}{Code
128} symbols. It can encode full 256 character extended-ASCII.

Variants:

\begin{itemize}
\tightlist
\item
  \protect\hyperlink{hibc-symbols}{HIBC Codablock F} is a variant of
  Codablock F that should be used when encoding HIBC formatted data.
\end{itemize}

Standards: USS Codablock F.

\hypertarget{data-and-options-62}{%
\paragraph{Data and Options}\label{data-and-options-62}}

\begin{itemize}
\tightlist
\item
  The data field can consist of any extended-ASCII data. The default
  interpretation of data by readers is in accordance with ISO/IEC
  8859-1.
\item
  \emph{FNC4} function characters will be inserted automatically to
  allow the encoding of extended ASCII characters.
\item
  When the \textbf{parse} option is specified, any instances of
  \texttt{\^{}NNN} in the data field are replaced with their equivalent
  ASCII value, useful for specifying unprintable characters.
\item
  When the \textbf{parsefnc} option is specified, non-data function
  characters can be specified by escape sequences:

  \begin{itemize}
  \tightlist
  \item
    \texttt{\^{}FNC1}: FNC1
  \item
    \texttt{\^{}FNC2}: FNC2
  \item
    \texttt{\^{}FNC3}: FNC3
  \item
    \texttt{\^{}ECI000000} to \texttt{\^{}ECI999999}: ECI indicators
  \end{itemize}
\item
  The \textbf{columns} option specifies the number of columns in the
  symbol, default 8.
\item
  The \textbf{rows} option specifies the number of rows in the symbol,
  between 2 and 44.
\item
  If \textbf{rows} is unspecified the encoder will default to the
  smallest number of rows that can hold the given data.
\item
  The \textbf{raw} option denotes that the data field is providing the
  input as a pre-encoded codewords in \texttt{\^{}NNN} format, suitable
  for direct low-level encoding.
\item
  The \textbf{rowheight} option specifies the height of the bars in each
  row in points. The default is 10.
\item
  The \textbf{sepheight} option specifies the height of the separator
  bars enclosing the rows in points. The default is 1.
\end{itemize}

\hypertarget{examples-40}{%
\paragraph{Examples}\label{examples-40}}

\begin{verbatim}
Data:    Codablock F
Options: 
Encoder: codablockf
\end{verbatim}

\includegraphics{images/codablockf-1.eps}

\begin{verbatim}
Data:    CODABLOCK F 34567890123456789010040digit
Options: columns=8 rows=5
Encoder: codablockf
\end{verbatim}

\includegraphics{images/codablockf-2.eps}

\begin{verbatim}
Data:    Short bars, fat seperators
Options: columns=10 rows=8 rowheight=6 sepheight=2
Encoder: codablockf
\end{verbatim}

\includegraphics{images/codablockf-3.eps}

\hypertarget{code-11}{%
\subsubsection{Code 11}\label{code-11}}

\textbf{Code 11} is a linear, discrete, non-self-checking,
bidirectional, numeric barcode symbology that is primarily used for
labelling telecommunication equipment.

Also known as: USD-8.

\hypertarget{data-and-options-63}{%
\paragraph{Data and Options}\label{data-and-options-63}}

\begin{itemize}
\tightlist
\item
  The data consists of digits and the dash character \emph{-}.
\item
  The \textbf{includecheck} option calculates the check digits.
\item
  For less than 10 data digits a single check digit is used.
\item
  For 10 or more data digits two check digits are used.
\end{itemize}

\hypertarget{example-19}{%
\paragraph{Example}\label{example-19}}

\begin{verbatim}
Data:    0123456789
Options: includecheck includetext includecheckintext
Encoder: code11
\end{verbatim}

\includegraphics{images/code11-1.eps}

\hypertarget{code-16k}{%
\subsubsection{Code 16K}\label{code-16k}}

The \textbf{Code 16K} barcode symbology is 2D stacked-linear barcode
that can encode full 256 character extended-ASCII with the use of the
\emph{FNC4} shift character.

Also known as: USS-16K

Standards: ANSI/AIM BC7 - USS Code 16K, BS EN 12323.

\hypertarget{data-and-options-64}{%
\paragraph{Data and Options}\label{data-and-options-64}}

\begin{itemize}
\tightlist
\item
  The data field can consist of any extended ASCII data. The default
  interpretation of data by readers is in accordance with ISO/IEC
  8859-1.
\item
  \emph{FNC4} function characters will be inserted automatically to
  allow the encoding of extended ASCII characters.
\item
  When the \textbf{parse} option is specified, any instances of
  \texttt{\^{}NNN} in the data field are replaced with their equivalent
  ASCII value, useful for specifying unprintable characters.
\item
  When the \textbf{parsefnc} option is specified, non-data function
  characters can be specified by escape sequences:

  \begin{itemize}
  \tightlist
  \item
    \texttt{\^{}FNC1}: FNC1
  \item
    \texttt{\^{}FNC2}: FNC2
  \item
    \texttt{\^{}FNC3}: FNC3
  \end{itemize}
\item
  The \textbf{mode} option specifies the mode for the symbol. It is
  usual to leave this unspecified in which case the most appropriate
  mode that results in the shortest symbol is automatically selected
  based in the input data.

  \begin{itemize}
  \tightlist
  \item
    \texttt{mode=0} - Starting code set A
  \item
    \texttt{mode=1} - Starting code set B
  \item
    \texttt{mode=2} - Starting code set C
  \item
    \texttt{mode=3} - Starting code set B with implied \emph{FNC1}
  \item
    \texttt{mode=4} - Starting code set C with implied \emph{FNC1}
  \item
    \texttt{mode=5} - Starting code set C with implied \emph{Shift B}
  \item
    \texttt{mode=6} - Starting code set C with implied \emph{Double
    Shift B}
  \end{itemize}
\item
  The \textbf{pos} option specifies this symbol to be part of multi-part
  structured data. For example \texttt{pos=25} specifies this to be the
  second symbol in a group of five symbols.
\item
  The \textbf{rows} option specifies the number of rows in the symbol,
  between two and sixteen.
\item
  If \textbf{rows} is unspecified the encoder will default to the
  smallest number of rows that can hold the given data.
\item
  The \textbf{raw} option denotes that the data field is providing the
  input as a pre-encoded codewords in \texttt{\^{}NNN} format, suitable
  for direct low-level encoding.
\item
  The \textbf{rowheight} option specifies the height of the bars in each
  row in points. The default is 10.
\item
  The \textbf{sepheight} option specifies the height of the separator
  bars enclosing the rows in points. The default is 1.
\end{itemize}

\hypertarget{examples-41}{%
\paragraph{Examples}\label{examples-41}}

\begin{verbatim}
Data:    Abcd-1234567890-wxyZ
Options: 
Encoder: code16k
\end{verbatim}

\includegraphics{images/code16k-1.eps}

\begin{verbatim}
Data:    Code 16K
Options: rows=10
Encoder: code16k
\end{verbatim}

\includegraphics{images/code16k-2.eps}

\begin{verbatim}
Data:    Short bars, fat seperators
Options: rows=8 rowheight=5 sepheight=2
Encoder: code16k
\end{verbatim}

\includegraphics{images/code16k-3.eps}

\hypertarget{code-25}{%
\subsubsection{Code 25}\label{code-25}}

\textbf{Code 2 of 5} is a simple low density numeric barcode symbology.

Also known as: Code 25, Industrial 2 of 5, Standard 2 of 5

Variants:

\begin{itemize}
\tightlist
\item
  \protect\hyperlink{iata-2-of-5}{IATA 2 of 5}, Computer Identics 2 of
  5.
\item
  \protect\hyperlink{datalogic-2-of-5}{Datalogic 2 of 5}.
\item
  \protect\hyperlink{matrix-2-of-5}{Matrix 2 of 5}.
\item
  \protect\hyperlink{coop-2-of-5}{COOP 2 of 5}.
\end{itemize}

\hypertarget{data-and-options-65}{%
\paragraph{Data and Options}\label{data-and-options-65}}

\begin{itemize}
\tightlist
\item
  The data consists of any number of digits.
\item
  The \textbf{includecheck} option calculates the check digit.
\item
  \emph{Deprecated}: The \textbf{version} option determines which
  variant to use:

  \begin{itemize}
  \tightlist
  \item
    \texttt{version=industrial} (default) - Industrial 2 of 5.
  \item
    \texttt{version=iata} - \emph{Deprecated: Use
    \protect\hyperlink{iata-2-of-5}{IATA 2 of 5}}
  \item
    \texttt{version=datalogic} - \emph{Deprecated: Use
    \protect\hyperlink{datalogic-2-of-5}{Datalogic 2 of 5}}
  \item
    \texttt{version=matrix} - \emph{Deprecated: Use
    \protect\hyperlink{matrix-2-of-5}{Matrix 2 of 5}}
  \item
    \texttt{version=coop} - \emph{Deprecated: Use
    \protect\hyperlink{coop-2-of-5}{COOP 2 of 5}}
  \end{itemize}
\end{itemize}

\hypertarget{examples-42}{%
\paragraph{Examples}\label{examples-42}}

\begin{verbatim}
Data:    01234567
Options: 
Encoder: code2of5
\end{verbatim}

\includegraphics{images/code25-1.eps}

\hypertarget{dotcode}{%
\subsubsection{DotCode}\label{dotcode}}

The \textbf{DotCode} symbology is 2D matrix-style barcode that can
encode full 256 character extended-ASCII.

Also known as: Ribbon Code

Variants:

\begin{itemize}
\tightlist
\item
  \protect\hyperlink{gs1-dotcode}{GS1 DotCode} is a variant of DotCode
  that should be used when encoding data that is in
  \protect\hyperlink{gs1-application-identifier-standard-format}{GS1
  Application Identifier standard format}.
\end{itemize}

Standards: AIM - ISS DotCode.

\hypertarget{data-and-options-66}{%
\paragraph{Data and Options}\label{data-and-options-66}}

\begin{itemize}
\tightlist
\item
  The data field can contain any extended ASCII data. The default
  interpretation of data by readers is in accordance with ISO/IEC
  8859-1.
\item
  When the \textbf{parse} option is specified, any instances of
  \texttt{\^{}NNN} in the data field are replaced with their equivalent
  ASCII value. This is useful for specifying unprintable characters.
\item
  When the \textbf{parsefnc} option is specified, non-data function
  characters can be specified by escape sequences:

  \begin{itemize}
  \tightlist
  \item
    \texttt{\^{}FNC1}: FNC1. \emph{Recommendation: Use
    \protect\hyperlink{gs1-dotcode}{GS1 DotCode} for encoding GS1 AI
    syntax data.}
  \item
    \texttt{\^{}FNC2}: FNC2
  \item
    \texttt{\^{}FNC3}: FNC3
  \end{itemize}
\item
  The \textbf{ratio} is used to specify the aspect ratio of the symbol.
  The default is 1.5.
\item
  The \textbf{columns} and \textbf{rows} options are used to specify the
  size of the symbol. When these are not specified a symbol is generated
  that maintains the selected aspect ratio.
\item
  \emph{Deprecated: Use only if performance is a concern.} The
  \textbf{fast} option is used to enable the optional high-speed symbol
  masking algorithm.
\end{itemize}

\hypertarget{examples-43}{%
\paragraph{Examples}\label{examples-43}}

\begin{verbatim}
Data:    This is DotCode
Options: inkspread=0.16
Encoder: dotcode
\end{verbatim}

\includegraphics{images/dotcode-1.eps}

\hypertarget{ultracode}{%
\subsubsection{Ultracode}\label{ultracode}}

The \textbf{Ultracode} symbology is a colour, 2D matrix-style barcode
symbology that can encode full 256 character extended-ASCII.

Standards: AIM ISS - Ultracode.

\hypertarget{data-and-options-67}{%
\paragraph{Data and Options}\label{data-and-options-67}}

\begin{itemize}
\tightlist
\item
  The data field can contain any extended ASCII data. The default
  interpretation of data by readers is in accordance with ISO/IEC
  8859-1.
\item
  When the \textbf{parse} option is specified, any instances of
  \texttt{\^{}NNN} in the data field are replaced with their equivalent
  ASCII value, useful for specifying unprintable characters.
\item
  When the \textbf{parsefnc} option is specified, non-data function
  characters can be specified by escape sequences:

  \begin{itemize}
  \tightlist
  \item
    \texttt{\^{}FNC1}: FNC1
  \item
    \texttt{\^{}FNC3}: FNC3
  \end{itemize}
\item
  The \textbf{eclevel} option is used to specify the error correction
  level:

  \begin{itemize}
  \tightlist
  \item
    \texttt{eclevel=EC0} - Error detection only
  \item
    \texttt{eclevel=EC1} - Low
  \item
    \texttt{eclevel=EC2} - Medium (default)
  \item
    \texttt{eclevel=EC3}
  \item
    \texttt{eclevel=EC4}
  \item
    \texttt{eclevel=EC5} - Highest
  \end{itemize}
\end{itemize}

\hypertarget{examples-44}{%
\paragraph{Examples}\label{examples-44}}

\begin{verbatim}
Data:    Nice colours!
Options: eclevel=EC3
Encoder: ultracode
\end{verbatim}

\includegraphics{images/ultracode-1.eps}

\hypertarget{iata-2-of-5}{%
\subsubsection{IATA 2 of 5}\label{iata-2-of-5}}

\textbf{IATA 2 of 5} is a variant of the
\protect\hyperlink{code-25}{Code 2 of 5} barcode symbology.

Also known as: Computer Identics 2 of 5.

Variants:

\begin{itemize}
\tightlist
\item
  \protect\hyperlink{code-25}{Industrial 2 of 5}, Standard 2 of 5.
\item
  \protect\hyperlink{datalogic-2-of-5}{Datalogic 2 of 5}.
\item
  \protect\hyperlink{matrix-2-of-5}{Matrix 2 of 5}.
\item
  \protect\hyperlink{coop-2-of-5}{COOP 2 of 5}.
\end{itemize}

\hypertarget{data-and-options-68}{%
\paragraph{Data and Options}\label{data-and-options-68}}

\begin{itemize}
\tightlist
\item
  The data consists of any number of digits.
\item
  The \textbf{includecheck} option calculates the check digit.
\end{itemize}

\hypertarget{examples-45}{%
\paragraph{Examples}\label{examples-45}}

\begin{verbatim}
Data:    01234567
Options: includetext includecheck includecheckintext
Encoder: iata2of5
\end{verbatim}

\includegraphics{images/code25-2.eps}

\hypertarget{matrix-2-of-5}{%
\subsubsection{Matrix 2 of 5}\label{matrix-2-of-5}}

\textbf{Matrix 2 of 5} is a variant of the
\protect\hyperlink{code-25}{Code 2 of 5} barcode symbology.

Variants:

\begin{itemize}
\tightlist
\item
  \protect\hyperlink{code-25}{Industrial 2 of 5}, Standard 2 of 5.
\item
  \protect\hyperlink{iata-2-of-5}{IATA 2 of 5}, Computer Identics 2 of
  5.
\item
  \protect\hyperlink{datalogic-2-of-5}{Datalogic 2 of 5}.
\item
  \protect\hyperlink{coop-2-of-5}{COOP 2 of 5}.
\end{itemize}

\hypertarget{data-and-options-69}{%
\paragraph{Data and Options}\label{data-and-options-69}}

\begin{itemize}
\tightlist
\item
  The data consists of any number of digits.
\item
  The \textbf{includecheck} option calculates the check digit.
\end{itemize}

\hypertarget{datalogic-2-of-5}{%
\subsubsection{Datalogic 2 of 5}\label{datalogic-2-of-5}}

\textbf{Datalogic 2 of 5} is a variant of the
\protect\hyperlink{code-25}{Code 2 of 5} barcode symbology.

Variants:

\begin{itemize}
\tightlist
\item
  \protect\hyperlink{code-25}{Industrial 2 of 5}, Standard 2 of 5.
\item
  \protect\hyperlink{iata-2-of-5}{IATA 2 of 5}, Computer Identics 2 of
  5.
\item
  \protect\hyperlink{matrix-2-of-5}{Matrix 2 of 5}.
\item
  \protect\hyperlink{coop-2-of-5}{COOP 2 of 5}.
\end{itemize}

\hypertarget{data-and-options-70}{%
\paragraph{Data and Options}\label{data-and-options-70}}

\begin{itemize}
\tightlist
\item
  The data consists of any number of digits.
\item
  The \textbf{includecheck} option calculates the check digit.
\end{itemize}

\hypertarget{coop-2-of-5}{%
\subsubsection{COOP 2 of 5}\label{coop-2-of-5}}

\textbf{COOP 2 of 5} is a variant of the
\protect\hyperlink{code-25}{Code 2 of 5} barcode symbology.

Variants:

\begin{itemize}
\tightlist
\item
  \protect\hyperlink{code-25}{Industrial 2 of 5}, Standard 2 of 5.
\item
  \protect\hyperlink{iata-2-of-5}{IATA 2 of 5}, Computer Identics 2 of
  5.
\item
  \protect\hyperlink{datalogic-2-of-5}{Datalogic 2 of 5}.
\item
  \protect\hyperlink{matrix-2-of-5}{Matrix 2 of 5}.
\end{itemize}

\hypertarget{data-and-options-71}{%
\paragraph{Data and Options}\label{data-and-options-71}}

\begin{itemize}
\tightlist
\item
  The data consists of any number of digits.
\item
  The \textbf{includecheck} option calculates the check digit.
\end{itemize}

\hypertarget{code-49}{%
\subsubsection{Code 49}\label{code-49}}

The \textbf{Code 49} barcode symbology is 2D stacked-linear barcode that
can encode 128 character ASCII.

Also known as: USS-49.

Standards: ANSI/AIM BC6 - USS Code 49.

\hypertarget{data-and-options-72}{%
\paragraph{Data and Options}\label{data-and-options-72}}

\begin{itemize}
\tightlist
\item
  The input can consist of any ASCII data.
\item
  When the \textbf{parse} option is specified, any instances of
  \texttt{\^{}NNN} in the data field are replaced with their equivalent
  ASCII value, useful for specifying unprintable characters.
\item
  When the \textbf{parsefnc} option is specified, non-data function
  characters can be specified by escape sequences:

  \begin{itemize}
  \tightlist
  \item
    \texttt{\^{}FNC1}: FNC1
  \item
    \texttt{\^{}FNC2}: FNC2
  \item
    \texttt{\^{}FNC3}: FNC3
  \end{itemize}
\item
  The \textbf{mode} option specifies the mode for the symbol. It is
  usual to leave this unspecified in which case the most appropriate
  mode that results in the shortest symbol is automatically selected
  based in the input data.

  \begin{itemize}
  \tightlist
  \item
    \texttt{mode=0} - regular alphanumeric mode
  \item
    \texttt{mode=1} - append mode
  \item
    \texttt{mode=2} - numeric mode
  \item
    \texttt{mode=3} - group alphanumeric mode
  \item
    \texttt{mode=4} - alphanumeric mode starting shift 1
  \item
    \texttt{mode=5} - alphanumeric mode starting shift 2
  \item
    \texttt{mode=6} - reserved
  \end{itemize}
\item
  The \textbf{pos} option specifies this symbol to be part of multi-part
  structured data, i.e. selecting \texttt{mode=3}. For example
  \texttt{pos=25} specifies this to be the second symbol in a group of
  five symbols.
\item
  The \textbf{rows} option specifies the number of rows in the symbol,
  between \emph{2} and \emph{8}.
\item
  If \textbf{rows} is unspecified the encoder will default to the
  smallest number of rows that can hold the given data.
\item
  The \textbf{rowheight} option specifies the height of the bars in each
  row in points. The default is \emph{10}.
\item
  The \textbf{sepheight} option specifies the height of the separator
  bars enclosing the rows in points. The default is \emph{1}.
\end{itemize}

\hypertarget{examples-46}{%
\paragraph{Examples}\label{examples-46}}

\begin{verbatim}
Data:    MULTIPLE ROWS IN CODE 49
Options: 
Encoder: code49
\end{verbatim}

\includegraphics{images/code49-1.eps}

\begin{verbatim}
Data:    CODE 49
Options: rows=8
Encoder: code49
\end{verbatim}

\includegraphics{images/code49-2.eps}

\begin{verbatim}
Data:    Short bars, fat seperators
Options: rows=8 rowheight=6 sepheight=2
Encoder: code49
\end{verbatim}

\includegraphics{images/code49-3.eps}

\hypertarget{code-one}{%
\subsubsection{Code One}\label{code-one}}

\textbf{Code One} was the earliest public domain 2D matrix-style
barcode. It is used by the health care and recycling industry and can
encode full 256 character extended-ASCII.

Also known as: Code 1, Code 1S.

Standards: AIM USS - Code One.

\hypertarget{data-and-options-73}{%
\paragraph{Data and Options}\label{data-and-options-73}}

\begin{itemize}
\tightlist
\item
  The data field can consist of any ASCII data for \emph{standard} and
  \emph{T-type} symbols. The default interpretation of data by readers
  is in accordance with ISO/IEC 8859-1. When supported by the receiver
  characters from other code pages can be encoded using Extended Channel
  Interpretation (ECI).
\item
  Note: \emph{S-type} symbols are special in that they represent a
  numeric value so may only contain digits.
\item
  When the \textbf{parse} option is specified, any instances of
  \texttt{\^{}NNN} in the data field are replaced with their equivalent
  ASCII value, useful for specifying unprintable characters, e.g.
  \texttt{\^{}029} for \emph{GS}, etc.
\item
  When the \textbf{parsefnc} option is specified, non-data function
  characters can be specified by escape sequences:

  \begin{itemize}
  \tightlist
  \item
    \texttt{\^{}FNC1}: FNC1
  \item
    \texttt{\^{}FNC3}: FNC3
  \item
    \texttt{\^{}ECI000000} to \texttt{\^{}ECI999999}: ECI indicators
  \end{itemize}
\item
  The \textbf{version} option is used to specify the size and type of
  the symbol:

  \begin{itemize}
  \tightlist
  \item
    \texttt{A}, \texttt{B}, \texttt{C}, \texttt{D}, \texttt{E},
    \texttt{F}, \texttt{G}, \texttt{H} - for standard format symbols
    (default automatic selection)
  \item
    \texttt{version=T-16}, \texttt{version=T-32}, \texttt{version=T-48}
    - \emph{T-type} symbols
  \item
    \texttt{version=S-10}, \texttt{version=S-20}, \texttt{version=S-30}
    - \emph{S-type} symbols
  \end{itemize}
\end{itemize}

\hypertarget{examples-47}{%
\paragraph{Examples}\label{examples-47}}

\begin{verbatim}
Data:    Code One
Options: 
Encoder: codeone
\end{verbatim}

\includegraphics{images/codeone-1.eps}

\begin{verbatim}
Data:    Code One
Options: version=C
Encoder: codeone
\end{verbatim}

\includegraphics{images/codeone-2.eps}

\begin{verbatim}
Data:    Code One
Options: version=T-32
Encoder: codeone
\end{verbatim}

\includegraphics{images/codeone-3.eps}

\begin{verbatim}
Data:    406990
Options: version=S-10
Encoder: codeone
\end{verbatim}

\includegraphics{images/codeone-4.eps}

\hypertarget{msi-plessey}{%
\subsubsection{MSI Plessey}\label{msi-plessey}}

\textbf{MSI Plessey} is a continuous, non-self-checking, arbitrary
length, numeric barcode symbology.

Also known as: MSI, MSI Modified Plessey.

Variants:

\begin{itemize}
\tightlist
\item
  \protect\hyperlink{plessey}{Plessey (UK)} is the original barcode upon
  which MSI Modified Plessey was based.
\end{itemize}

\hypertarget{data-and-options-74}{%
\paragraph{Data and Options}\label{data-and-options-74}}

\begin{itemize}
\tightlist
\item
  The data can consist of any number of digits.
\item
  The \textbf{includecheck} option calculates the check digit or check
  digits.
\item
  The \textbf{includecheckintext} option makes the calculated check
  characters appear in the human readable text.
\item
  The \textbf{checktype} option is used to specify the type of checksum,
  either:

  \begin{itemize}
  \tightlist
  \item
    \texttt{checktype=mod10} (default)
  \item
    \texttt{checktype=mod1010}
  \item
    \texttt{checktype=mod11}
  \item
    \texttt{checktype=ncrmod11}
  \item
    \texttt{checktype=mod1110}
  \item
    \texttt{checktype=ncrmod1110}
  \end{itemize}
\item
  The \textbf{badmod11} option allows a \texttt{checktype=mod11}
  checksum value of 10 to be encoded with a pair of check digits
  \emph{10}. Normally in \texttt{checktype=mod11}, any input whose
  checksum evaluates to \emph{10} is considered invalid having no
  correct representation.
\end{itemize}

\hypertarget{examples-48}{%
\paragraph{Examples}\label{examples-48}}

\begin{verbatim}
Data:    0123456789
Options: includecheck includetext
Encoder: msi
\end{verbatim}

\includegraphics{images/msi-1.eps}

\begin{verbatim}
Data:    0123456789
Options: includecheck checktype=mod1110 includetext includecheckintext
Encoder: msi
\end{verbatim}

\includegraphics{images/msi-2.eps}

\begin{verbatim}
Data:    0123456785
Options: includecheck checktype=mod11 badmod11 includetext includecheckintext
Encoder: msi
\end{verbatim}

\includegraphics{images/msi-3.eps}

\hypertarget{plessey}{%
\subsubsection{Plessey}\label{plessey}}

\textbf{Plessey} is a continuous, arbitrary length barcode symbology for
encoding hexadecimal data.

Also known as: Anker Code.

Variants:

\begin{itemize}
\tightlist
\item
  \protect\hyperlink{msi-plessey}{MSI Modified Plessey} is a variant of
  the Plessey (UK) barcode developed by the MSI Data Corporation.
\end{itemize}

\hypertarget{data-and-options-75}{%
\paragraph{Data and Options}\label{data-and-options-75}}

\begin{itemize}
\tightlist
\item
  The data can contain any of the following:

  \begin{itemize}
  \tightlist
  \item
    Numbers \texttt{0-9}
  \item
    Capital letters \texttt{A-F}
  \end{itemize}
\item
  Two manditory check characters implementing a CRC check are
  automatically included.
\item
  The \textbf{includecheckintext} option makes the calculated check
  characters appear in the human readable text.
\item
  The \textbf{unidirectional} option generates a unidirectional Plessey
  symbol.
\end{itemize}

\hypertarget{examples-49}{%
\paragraph{Examples}\label{examples-49}}

Equivalent symbols, the latter displaying the two mandatory check
characters:

\begin{verbatim}
Data:    01234ABCD
Options: includetext
Encoder: plessey
\end{verbatim}

\includegraphics{images/plessey-1.eps}

\begin{verbatim}
Data:    01234ABCD
Options: includetext includecheckintext
Encoder: plessey
\end{verbatim}

\includegraphics{images/plessey-2.eps}

A unidirectional symbol:

\begin{verbatim}
Data:    01234ABCD
Options: unidirectional
Encoder: plessey
\end{verbatim}

\includegraphics{images/plessey-3.eps}

\hypertarget{posicode}{%
\subsubsection{PosiCode}\label{posicode}}

\textbf{PosiCode} is a continuous, variable length, non-self-checking,
bidirectional barcode symbology that is designed for use within printing
processes where it is difficult to precisely control the width of a bar.

Standards: ITS PosiCode.

\hypertarget{data-and-options-76}{%
\paragraph{Data and Options}\label{data-and-options-76}}

\begin{itemize}
\tightlist
\item
  The data field can hold the following:

  \begin{itemize}
  \tightlist
  \item
    For \emph{standard} symbols: Any extended ASCII data. The default
    interpretation of data by readers is in accordance with ISO/IEC
    8859-1.

    \begin{itemize}
    \tightlist
    \item
      \emph{FNC4} function characters will be inserted automatically to
      allow the encoding of extended ASCII characters.
    \end{itemize}
  \item
    For \emph{limited} symbols: letters \texttt{A-Z}, digits
    \texttt{0-9}, symbols \texttt{-} and \texttt{.}
  \end{itemize}
\item
  The \textbf{version} option is used to specify the variant of the
  symbol, either:

  \begin{itemize}
  \tightlist
  \item
    \texttt{version=a} (default)
  \item
    \texttt{version=b}
  \item
    \texttt{version=limiteda}
  \item
    \texttt{version=limitedb}
  \end{itemize}
\item
  When the \textbf{parse} option is specified, any instances of
  \texttt{\^{}NNN} in the data field are replaced with their equivalent
  ASCII value, useful for specifying unprintable characters, e.g.
  \texttt{\^{}029} for \emph{GS}, etc.
\item
  When the \textbf{parsefnc} option is specified, non-data function
  characters can be specified by escape sequences:

  \begin{itemize}
  \tightlist
  \item
    \texttt{\^{}FNC1}: FNC1
  \item
    \texttt{\^{}FNC2}: FNC2
  \item
    \texttt{\^{}FNC3}: FNC3
  \end{itemize}
\item
  The \textbf{inkspread} option can be used to adjust the width of the
  bars.
\end{itemize}

\hypertarget{example-posicode}{%
\paragraph{Example PosiCode}\label{example-posicode}}

Equivalent ways to generate a PosiCode A symbol:

\begin{verbatim}
Data:    Abc123
Options: 
Encoder: posicode
\end{verbatim}

\begin{verbatim}
Data:    Abc123
Options: version=a
Encoder: posicode
\end{verbatim}

\includegraphics{images/posicode-1.eps}

PosiCode A including a \emph{GS} (ASCII 29) character:

\begin{verbatim}
Data:    AB^029CD
Options: parse
Encoder: posicode
\end{verbatim}

\includegraphics{images/posicode-2.eps}

PosiCode A including an \emph{FNC2} special character:

\begin{verbatim}
Data:    AB^FNC2CD
Options: parsefnc
Encoder: posicode
\end{verbatim}

\includegraphics{images/posicode-3.eps}

PosiCode B symbol with widened bars:

\begin{verbatim}
Data:    Abc123
Options: version=b inkspread=-1
Encoder: posicode
\end{verbatim}

\includegraphics{images/posicode-4.eps}

\hypertarget{example-limited-posicode}{%
\paragraph{Example Limited PosiCode}\label{example-limited-posicode}}

Limited PosiCode A with narrowed bars:

\begin{verbatim}
Data:    ABC-12.3
Options: version=limiteda
Encoder: posicode
\end{verbatim}

\includegraphics{images/posicode-5.eps}

Limited PosiCode B:

\begin{verbatim}
Data:    ABC-12.3
Options: version=limitedb
Encoder: posicode
\end{verbatim}

\includegraphics{images/posicode-6.eps}

\hypertarget{telepen}{%
\subsubsection{Telepen}\label{telepen}}

\textbf{Telepen} is an arbitrary length barcode symbology for encoding
all 128 ASCII characters without the need for shift characters.

Also known as: Telepen Alpha, Telepen Full ASCII.

Variants:

\begin{itemize}
\tightlist
\item
  \protect\hyperlink{telepen-numeric}{Telepen Numeric}.
\end{itemize}

Standards: USS Telepen.

\hypertarget{data-and-options-77}{%
\paragraph{Data and Options}\label{data-and-options-77}}

\begin{itemize}
\tightlist
\item
  The data can contain any standard ASCII data, values 0-127.
\item
  When the \textbf{parse} option is specified, any instances of
  \texttt{\^{}NNN} in the data field are replaced with their equivalent
  ASCII value, useful for specifying unprintable characters.
\item
  The mandatory check digit is automatically included.
\item
  \emph{Deprecated: Use \protect\hyperlink{telepen-numeric}{Telepen
  Numeric} instead. When the \textbf{numeric} option is given, the data
  is read as either pairs of digits or \texttt{0X}, \texttt{1X}, etc.
  The singular values \texttt{\^{}000} to \texttt{\^{}016} can also be
  encoded using the \emph{parse} option.}
\end{itemize}

\hypertarget{examples-50}{%
\paragraph{Examples}\label{examples-50}}

\begin{verbatim}
Data:    Telepen
Options: includetext
Encoder: telepen
\end{verbatim}

\includegraphics{images/telepen-1.eps}

\begin{verbatim}
Data:    Telepen^013
Options: parse
Encoder: telepen
\end{verbatim}

\includegraphics{images/telepen-2.eps}

\hypertarget{telepen-numeric}{%
\subsubsection{Telepen Numeric}\label{telepen-numeric}}

\textbf{Telepen Numeric} is a variant of the Telepen symbology for
efficient encoding of numeric data.

Variants:

\begin{itemize}
\tightlist
\item
  \protect\hyperlink{telepen}{Telepen Alpha}.
\end{itemize}

Standards: USS Telepen.

\hypertarget{data-and-options-78}{%
\paragraph{Data and Options}\label{data-and-options-78}}

\begin{itemize}
\tightlist
\item
  The data is provided as either pairs of digits or \texttt{0X},
  \texttt{1X}, etc. The singular values \texttt{\^{}000} to
  \texttt{\^{}016} can also be encoded using the \emph{parse} option.
\item
  When the \textbf{parse} option is specified, any instances of
  \texttt{\^{}NNN} in the data field are replaced with their equivalent
  ASCII value, useful for specifying unprintable characters.
\item
  The mandatory check digit is automatically included.
\end{itemize}

\hypertarget{examples-51}{%
\paragraph{Examples}\label{examples-51}}

\begin{verbatim}
Data:    123456
Options: 
Encoder: telepennumeric
\end{verbatim}

\includegraphics{images/telepen-3.eps}

\begin{verbatim}
Data:    1X345X
Options: 
Encoder: telepennumeric
\end{verbatim}

\includegraphics{images/telepen-4.eps}

\begin{verbatim}
Data:    1234^005
Options: parse includetext
Encoder: telepennumeric
\end{verbatim}

\includegraphics{images/telepen-5.eps}

\hypertarget{gs1-composite-symbols}{%
\subsection{GS1 Composite Symbols}\label{gs1-composite-symbols}}

\textbf{GS1 Composite} barcode symbologies consist of a primary
component beneath a 2D component (variations of
\protect\hyperlink{micropdf417}{MicroPDF417} and
\protect\hyperlink{pdf417}{PDF417}) used to encode supplementary
\protect\hyperlink{gs1-application-identifier-standard-format}{GS1
formatted data}.

Variants:

\begin{itemize}
\tightlist
\item
  \protect\hyperlink{ean-13-composite}{EAN-13 Composite} is a variant of
  \protect\hyperlink{ean-13}{EAN-13}.
\item
  \protect\hyperlink{ean-8-composite}{EAN-8 Composite} is a variant of
  \protect\hyperlink{ean-8}{EAN-8}.
\item
  \protect\hyperlink{upc-a-composite}{UPC-A Composite} is a variant of
  \protect\hyperlink{upc-a}{UPC-A}.
\item
  \protect\hyperlink{upc-e-composite}{UPC-E Composite} is a variant of
  \protect\hyperlink{upc-e}{UPC-E}.
\item
  \protect\hyperlink{gs1-databar-omnidirectional-composite}{GS1 DataBar
  Omnidirectional Composite} is a variant of
  \protect\hyperlink{gs1-databar-omnidirectional}{GS1 DataBar
  Omnidirectional}.
\item
  \protect\hyperlink{gs1-databar-stacked-omnidirectional-composite}{GS1
  DataBar Stacked Omnidirectional Composite} is a variant of
  \protect\hyperlink{gs1-databar-stacked-omnidirectional}{GS1 DataBar
  Stacked Omnidirectional}.
\item
  \protect\hyperlink{gs1-databar-expanded-composite}{GS1 DataBar
  Expanded Composite} is a variant of
  \protect\hyperlink{gs1-databar-expanded}{GS1 DataBar Expanded}.
\item
  \protect\hyperlink{gs1-databar-expanded-stacked-composite}{GS1 DataBar
  Expanded Stacked Composite} is a variant of
  \protect\hyperlink{gs1-databar-expanded-stacked}{GS1 DataBar Expanded
  Stacked}.
\item
  \protect\hyperlink{gs1-databar-truncated-composite}{GS1 DataBar
  Truncated Composite} is a variant of
  \protect\hyperlink{gs1-databar-truncated}{GS1 DataBar Truncated}.
\item
  \protect\hyperlink{gs1-databar-stacked-composite}{GS1 DataBar Stacked
  Composite} is a variant of \protect\hyperlink{gs1-databar-stacked}{GS1
  DataBar Stacked}.
\item
  \protect\hyperlink{gs1-databar-limited-composite}{GS1 DataBar Limited
  Composite} is a variant of \protect\hyperlink{gs1-databar-limited}{GS1
  DataBar Limited}.
\item
  \protect\hyperlink{gs1-128-composite}{GS1-128 Composite} is a variant
  of \protect\hyperlink{gs1-128}{GS1-128}.
\end{itemize}

Standards: ISO/IEC 24723, ITS EAN.UCC Composite Symbology, AIM ISS -
EAN.UCC Composite Symbology, GS1 General Specifications.

\hypertarget{data-and-options-79}{%
\paragraph{Data and Options}\label{data-and-options-79}}

\begin{itemize}
\tightlist
\item
  The data field consists of a primary and secondary component separated
  by a pipe \texttt{\textbar{}} character.
\item
  The data for the primary component (preceding the pipe) is entered in
  a format identical to the corresponding non-composite barcode
  symbology.
\item
  The data for the 2D component (following the pipe) is entered in
  \protect\hyperlink{gs1-application-identifier-standard-format}{GS1
  Application Identifier standard format}.
\item
  Note that it is not necessary to specify \emph{FNC1} characters since
  the encoder will insert these automatically where necessary.
\item
  \emph{Deprecated.} The \textbf{dontlint} option disables GS1
  Application Identifier syntax validation allowing images to be
  generated for input that does not comply with GS1 standards.
\item
  For maximum efficiency, if the data for the 2D component contains a
  number of application identifiers matching any of the specifications
  below then they should be provided in this given order:

  \begin{itemize}
  \tightlist
  \item
    \texttt{(11)...(10)...}
  \item
    \texttt{(17)...(10)...}
  \item
    \texttt{(90)\{0-3\ digits\ not\ starting\ 0\}\{upper\ alpha\}...}
  \end{itemize}
\item
  \emph{Deprecated: Default is best unless your application dictates
  otherwise.} The \textbf{ccversion} option is used to select a specific
  2D component:

  \begin{itemize}
  \tightlist
  \item
    \texttt{ccversion=a} - \protect\hyperlink{cc-a}{CC-A}
  \item
    \texttt{ccversion=b} - \protect\hyperlink{cc-b}{CC-B}
  \item
    \texttt{ccversion=c} - \protect\hyperlink{cc-c}{CC-C}
    (\protect\hyperlink{gs1-128-composite}{GS1-128 Composite} only)
  \item
    If \textbf{ccversion} is not specified a CC-A component will be
    selected if the data will fit, otherwise a CC-B component will be
    used. In the case of \protect\hyperlink{gs1-128-composite}{GS1-128
    Composite} a CC-C component will be used if the data does not fit
    within either a CC-A or CC-B component.
  \end{itemize}
\end{itemize}

\hypertarget{ean-13-composite}{%
\subsubsection{EAN-13 Composite}\label{ean-13-composite}}

\begin{verbatim}
Data:    331234567890|(99)1234-abcd
Options: includetext guardwhitespace
Encoder: ean13composite
\end{verbatim}

\includegraphics{images/ean13composite-1.eps}

\hypertarget{ean-8-composite}{%
\subsubsection{EAN-8 Composite}\label{ean-8-composite}}

\begin{verbatim}
Data:    12345670|(21)A12345678
Options: includetext guardwhitespace
Encoder: ean8composite
\end{verbatim}

\includegraphics{images/ean8composite-1.eps}

\hypertarget{upc-a-composite}{%
\subsubsection{UPC-A Composite}\label{upc-a-composite}}

\begin{verbatim}
Data:    01600033610|(99)1234-abcd
Options: includetext
Encoder: upcacomposite
\end{verbatim}

\includegraphics{images/upcacomposite-1.eps}

\hypertarget{upc-e-composite}{%
\subsubsection{UPC-E Composite}\label{upc-e-composite}}

\begin{verbatim}
Data:    0121230|(15)021231
Options: includetext
Encoder: upcecomposite
\end{verbatim}

\includegraphics{images/upcecomposite-1.eps}

\hypertarget{gs1-databar-omnidirectional-composite}{%
\subsubsection{GS1 DataBar Omnidirectional
Composite}\label{gs1-databar-omnidirectional-composite}}

\begin{verbatim}
Data:    (01)03612345678904|(11)990102
Options: 
Encoder: databaromnicomposite
\end{verbatim}

\includegraphics{images/databaromnicomposite-1.eps}

\hypertarget{gs1-databar-stacked-omnidirectional-composite}{%
\subsubsection{GS1 DataBar Stacked Omnidirectional
Composite}\label{gs1-databar-stacked-omnidirectional-composite}}

\begin{verbatim}
Data:    (01)03612345678904|(11)990102
Options: 
Encoder: databarstackedomnicomposite
\end{verbatim}

\includegraphics{images/databarstackedomnicomposite-1.eps}

\hypertarget{gs1-databar-expanded-composite}{%
\subsubsection{GS1 DataBar Expanded
Composite}\label{gs1-databar-expanded-composite}}

\begin{verbatim}
Data:    (01)93712345678904(3103)001234|(91)1A2B3C4D5E
Options: 
Encoder: databarexpandedcomposite
\end{verbatim}

\includegraphics{images/databarexpandedcomposite-1.eps}

\hypertarget{gs1-databar-expanded-stacked-composite}{%
\subsubsection{GS1 DataBar Expanded Stacked
Composite}\label{gs1-databar-expanded-stacked-composite}}

\begin{verbatim}
Data:    (01)00012345678905(10)ABCDEF|(21)12345678
Options: segments=4
Encoder: databarexpandedstackedcomposite
\end{verbatim}

\includegraphics{images/databarexpandedstackedcomposite-1.eps}

\hypertarget{gs1-databar-truncated-composite}{%
\subsubsection{GS1 DataBar Truncated
Composite}\label{gs1-databar-truncated-composite}}

\begin{verbatim}
Data:    (01)03612345678904|(11)990102
Options: 
Encoder: databartruncatedcomposite
\end{verbatim}

\includegraphics{images/databartruncatedcomposite-1.eps}

\hypertarget{gs1-databar-stacked-composite}{%
\subsubsection{GS1 DataBar Stacked
Composite}\label{gs1-databar-stacked-composite}}

\begin{verbatim}
Data:    (01)03412345678900|(17)010200
Options: 
Encoder: databarstackedcomposite
\end{verbatim}

\includegraphics{images/databarstackedcomposite-1.eps}

\hypertarget{gs1-databar-limited-composite}{%
\subsubsection{GS1 DataBar Limited
Composite}\label{gs1-databar-limited-composite}}

\begin{verbatim}
Data:    (01)03512345678907|(21)abcdefghijklmnopqrstuv
Options: 
Encoder: databarlimitedcomposite
\end{verbatim}

\includegraphics{images/databarlimitedcomposite-1.eps}

\hypertarget{gs1-128-composite}{%
\subsubsection{GS1-128 Composite}\label{gs1-128-composite}}

GS1-128 Composite with a CC-A 2D component:

\begin{verbatim}
Data:    (01)03212345678906|(21)A1B2C3D4E5F6G7H8
Options: 
Encoder: gs1-128composite
\end{verbatim}

\includegraphics{images/ean128composite-1.eps}

GS1-128 Composite with a CC-C 2D component:

\begin{verbatim}
Data:    (00)030123456789012340|(02)13012345678909(37)24(10)1234567ABCDEFG
Options: ccversion=c
Encoder: gs1-128composite
\end{verbatim}

\includegraphics{images/ean128composite-2.eps}

\hypertarget{cc-a}{%
\subsubsection{CC-A}\label{cc-a}}

Isolated CC-A 2D component:

\begin{verbatim}
Data:    (01)95012345678903
Options: ccversion=a cccolumns=3
Encoder: gs1-cc
\end{verbatim}

\includegraphics{images/gs1composite-1.eps}

\hypertarget{cc-b}{%
\subsubsection{CC-B}\label{cc-b}}

Isolated CC-B 2D component:

\begin{verbatim}
Data:    (01)95012345678903(3103)000123
Options: ccversion=b cccolumns=4
Encoder: gs1-cc
\end{verbatim}

\includegraphics{images/gs1composite-2.eps}

\hypertarget{cc-c}{%
\subsubsection{CC-C}\label{cc-c}}

Isolated CC-C 2D component:

\begin{verbatim}
Data:    (02)13012345678909(37)24(10)1234567ABCDEFG
Options: ccversion=c cccolumns=5
Encoder: gs1-cc
\end{verbatim}

\includegraphics{images/gs1composite-3.eps}

\hypertarget{raw-symbols}{%
\subsection{Raw Symbols}\label{raw-symbols}}

\hypertarget{daft}{%
\subsubsection{DAFT}\label{daft}}

\textbf{DAFT} is an encoder for directly specifying the descender,
ascender, full-height, tracker-bar succession for a custom 4 state
barcode symbol.

\hypertarget{data-and-options-80}{%
\paragraph{Data and Options}\label{data-and-options-80}}

\begin{itemize}
\tightlist
\item
  The data field contains a sequence of the characters \texttt{D},
  \texttt{A}, \texttt{F} or \texttt{T} to denote the descender,
  ascender, full-height and tracker bars of a custom 4 state symbol.
\end{itemize}

\hypertarget{example-20}{%
\paragraph{Example}\label{example-20}}

\begin{verbatim}
Data:    FATDAFTDAD
Options: 
Encoder: daft
\end{verbatim}

\includegraphics{images/daft-1.eps}

\hypertarget{flattermarken}{%
\subsubsection{Flattermarken}\label{flattermarken}}

\textbf{Flattermarken} are identification marks used in book production
that facilitate the proper arrangement of bound sections by a book
binder.

\hypertarget{data-and-options-81}{%
\paragraph{Data and Options}\label{data-and-options-81}}

\begin{itemize}
\tightlist
\item
  The data field can holding any sequence of digits corresponds to a 9
  module width with the following meaning:

  \begin{itemize}
  \tightlist
  \item
    \texttt{1-9}: a single mark exists in the corresponding module
    position
  \item
    \texttt{0}: unmarked sequence of modules
  \end{itemize}
\item
  The \textbf{inkspread} option can be used to adjust the width of the
  bars.
\item
  If greater fidelity is required then the \protect\hyperlink{raw}{raw}
  encoder should be used instead.
\end{itemize}

\hypertarget{example-21}{%
\paragraph{Example}\label{example-21}}

\begin{verbatim}
Data:    1304
Options: inkspread=-1
Encoder: flattermarken
\end{verbatim}

\includegraphics{images/flattermarken-1.eps}

\hypertarget{raw}{%
\subsubsection{Raw}\label{raw}}

The \textbf{raw} encoder is used for directly specifying the space/bar
succession of a custom barcode symbol.

\hypertarget{data-and-options-82}{%
\paragraph{Data and Options}\label{data-and-options-82}}

\begin{itemize}
\tightlist
\item
  The data field contains an alternating sequence of widths (1 to 9) for
  the bars and spaces of a custom symbol.
\end{itemize}

\hypertarget{example-22}{%
\paragraph{Example}\label{example-22}}

\begin{verbatim}
Data:    331132131313411122131311333213114131131221323
Options: height=0.5
Encoder: raw
\end{verbatim}

\includegraphics{images/raw-1.eps}

\hypertarget{partial-symbols}{%
\subsection{Partial Symbols}\label{partial-symbols}}

\hypertarget{ean-2}{%
\subsubsection{EAN-2}\label{ean-2}}

\textbf{EAN-2} is the two-digit add-on code that accompanies a EAN or
UPC type barcode symbol such as an \protect\hyperlink{isbn}{ISBN} or
\protect\hyperlink{issn}{ISSN}.

Also known as: Two-Digit Add-On, Two-Digit Supplement, UPC-2

\hypertarget{data-and-options-83}{%
\paragraph{Data and Options}\label{data-and-options-83}}

\begin{itemize}
\tightlist
\item
  The data field must contain two digits.
\item
  The \textbf{includetext} option should normally be supplied.
\end{itemize}

\hypertarget{example-23}{%
\paragraph{Example}\label{example-23}}

\begin{verbatim}
Data:    05
Options: includetext guardwhitespace
Encoder: ean2
\end{verbatim}

\includegraphics{images/ean2-1.eps}

\hypertarget{ean-5}{%
\subsubsection{EAN-5}\label{ean-5}}

\textbf{EAN-5} is the five-digit add-on code that accompanies an EAN or
UPC type barcode symbol such as an \protect\hyperlink{isbn}{ISBN} or
\protect\hyperlink{issn}{ISSN}.

Also known as: Five-Digit Add-On, Five-Digit Supplement, UPC-5

\hypertarget{data-and-options-84}{%
\paragraph{Data and Options}\label{data-and-options-84}}

\begin{itemize}
\tightlist
\item
  The data field must contain five digits.
\item
  The \textbf{includetext} option should normally be supplied.
\end{itemize}

\hypertarget{example-24}{%
\paragraph{Example}\label{example-24}}

\begin{verbatim}
Data:    90200
Options: includetext guardwhitespace
Encoder: ean5
\end{verbatim}

\includegraphics{images/ean5-1.eps}

\hypertarget{gs1-application-identifier-standard-format}{%
\subsection{GS1 Application Identifier Standard
Format}\label{gs1-application-identifier-standard-format}}

Certain barcode symbologies (including
\protect\hyperlink{gs1-128}{GS1-128},
\protect\hyperlink{gs1-databar-expanded}{GS1 DataBar Expanded},
\protect\hyperlink{gs1-datamatrix}{GS1 DataMatrix},
\protect\hyperlink{gs1-qr-code}{GS1 QR Code} and
\protect\hyperlink{gs1-composite-symbols}{GS1 Composite Symbols})
represent standardized GS1 data and require that their data field is
provided in GS1 Application Identifier standard format, consisting of a
concatenated string of \emph{AIs} along with their corresponding values.

The AIs are a set of approximately one hundred two-, three- or
four-digit prefixes written within parentheses that represent physical
attributes and business information, e.g.

\begin{itemize}
\tightlist
\item
  \emph{(00)} is an eighteen-digit SSCC.
\item
  \emph{(01)} is a fourteen-digit GTIN.
\item
  \emph{(403)} is a variable-length routing code.
\end{itemize}

The following input represents GTIN \emph{0061414199996}; Expiration
Date \emph{1 January 2010}; Batch \emph{123ABC}; Serial
\emph{1234567890}:

\begin{verbatim}
(01)0061414199996(17)100101(10)123ABC(21)1234567890
\end{verbatim}

Encoders for barcode symbologies that expect data in GS1 Application
Identifier standard format will take care of parsing the input and
inserting any necessary \emph{FNC1} characters to delimit variable
length fields.

Note: Instances of parentheses, \texttt{(} and \texttt{)}, within
Application Identifier values in the input must be escaped using the
\protect\hyperlink{input-processingux5cux23parse}{\texttt{parse} option}
as \texttt{\^{}040} and \texttt{\^{}041}, respectively.

\hypertarget{gs1-application-identifier-linting}{%
\paragraph{GS1 Application Identifier
Linting}\label{gs1-application-identifier-linting}}

The input is checked against the structural rules for the GS1
Application Identifiers, as described in
\href{https://www.linkedin.com/pulse/gs1-application-identifier-syntax-dictionary-terry-burton/}{this
article}, unless the \textbf{dontlint} option is given. The Application
Identifier definitions are provided in the
\href{http://www.gs1.org/barcodes-epcrfid-id-keys/gs1-general-specifications}{GS1
General Specifications}.

\hypertarget{options-reference}{%
\section{Options Reference}\label{options-reference}}

\hypertarget{check-digits}{%
\subsection{Check Digits}\label{check-digits}}

\hypertarget{includecheck}{%
\subsubsection{includecheck}\label{includecheck}}

Generate check digit(s) for symbologies where the use of check digits is
optional.

\hypertarget{example-25}{%
\paragraph{Example}\label{example-25}}

Calculate the optional check characters of this Code 93 symbol:

\begin{verbatim}
Data:    CHECK ME OUT
Options: includecheck
Encoder: code93
\end{verbatim}

\includegraphics{images/optincludecheck.eps}

\hypertarget{includecheckintext}{%
\subsubsection{includecheckintext}\label{includecheckintext}}

Show the calculated check digit in the human readable text.

\hypertarget{notes}{%
\paragraph{Notes}\label{notes}}

\begin{itemize}
\tightlist
\item
  For barcode symbologies where the check digit is not mandatory, this
  option must be used in combination with \textbf{includecheck}.
\item
  If any part of the checksum does not have a printable representation
  then that part is not displayed.
\end{itemize}

\hypertarget{example-26}{%
\paragraph{Example}\label{example-26}}

Display the check digit of this Royal Mail barcode:

\begin{verbatim}
Data:    LE28HS9Z
Options: includetext includecheckintext
Encoder: royalmail
\end{verbatim}

\includegraphics{images/optincludecheckintext.eps}

\hypertarget{input-processing}{%
\subsection{Input Processing}\label{input-processing}}

\hypertarget{parse}{%
\subsubsection{parse}\label{parse}}

In supporting barcode symbologies, when the \emph{parse} option is
specified, any instances of \texttt{\^{}NNN} in the data field are
replaced with their equivalent ASCII value, useful for specifying
unprintable characters.

Additionally, control character names can be used to specify equivalent
ASCII values, as follows:

\begin{longtable}[]{@{}lr@{}}
\toprule()
Sequence & ASCII value \\
\midrule()
\endhead
\texttt{\^{}NUL} & 0 \\
\texttt{\^{}SOH} & 1 \\
\texttt{\^{}STX} & 2 \\
\texttt{\^{}ETX} & 3 \\
\texttt{\^{}EOT} & 4 \\
\texttt{\^{}ENQ} & 5 \\
\texttt{\^{}ACK} & 6 \\
\texttt{\^{}BEL} & 7 \\
\texttt{\^{}BS} & 8 \\
\texttt{\^{}TAB} & 9 \\
\texttt{\^{}LF} & 10 \\
\texttt{\^{}VT} & 11 \\
\texttt{\^{}FF} & 12 \\
\texttt{\^{}CR} & 13 \\
\texttt{\^{}DLE} & 16 \\
\texttt{\^{}DC1} & 17 \\
\texttt{\^{}DC2} & 18 \\
\texttt{\^{}DC3} & 19 \\
\texttt{\^{}DC4} & 20 \\
\texttt{\^{}NAK} & 21 \\
\texttt{\^{}SYN} & 22 \\
\texttt{\^{}ETB} & 23 \\
\texttt{\^{}CAN} & 24 \\
\texttt{\^{}EM} & 25 \\
\texttt{\^{}SUB} & 26 \\
\texttt{\^{}ESC} & 27 \\
\texttt{\^{}FS} & 28 \\
\texttt{\^{}GS} & 29 \\
\texttt{\^{}RS} & 30 \\
\texttt{\^{}US} & 31 \\
\bottomrule()
\end{longtable}

ASCII control characters \texttt{SO} and \texttt{SI} cannot be encoded
by name (since \texttt{SO} would prefix clash with \texttt{SOH}).

Note: When this option is enabled, literal instances of "\^{}" in the
data should be escaped as \texttt{\^{}094} to avoid replacement if the
subsequent data characters represent a valid substitution. For example,
literal "\^{}123" in the input data can be escaped as
\texttt{\^{}094123}, literal "\^{}RS" in the data can be escaped as
\texttt{\^{}094RS}, and literal "\^{}ABC" in the data cam be escaped as
\texttt{\^{}094ABC} (in case \texttt{\^{}ABC} becomes a valid
substitution in the future).

\hypertarget{example-27}{%
\paragraph{Example}\label{example-27}}

Equivalent symbols:

\begin{verbatim}
Data:    This is Data Matrix
Options: 
Encoder: datamatrix
\end{verbatim}

\begin{verbatim}
Data:    This is ^068ata Matrix
Options: parse
Encoder: datamatrix
\end{verbatim}

\includegraphics{images/optparse.eps}

\hypertarget{parsefnc}{%
\subsubsection{parsefnc}\label{parsefnc}}

In supporting barcode symbologies, when the \textbf{parsefnc} option is
specified, non-data function characters can be specified by escaped
combinations such as \texttt{\^{}FNC1}, \texttt{\^{}FNC4},
\texttt{\^{}SFT/} and \texttt{\^{}ECI000003}.

Note: When this option is enabled, literal instances of \texttt{\^{}} in
the input data can be escaped as \texttt{\^{}\^{}}.

\hypertarget{example-28}{%
\paragraph{Example}\label{example-28}}

Manually-composed Code 93 Extended demonstrating use of the special
shift combination \texttt{(/)A} to represent \emph{!}:

\begin{verbatim}
Data:    TERRY^SFT/A
Options: parsefnc includecheck
Encoder: code93
\end{verbatim}

\includegraphics{images/optparsefnc.eps}

\hypertarget{symbol-dimensions}{%
\subsection{Symbol Dimensions}\label{symbol-dimensions}}

\hypertarget{height}{%
\subsubsection{height}\label{height}}

Height of longest bar, in inches.

\hypertarget{example-29}{%
\paragraph{Example}\label{example-29}}

A 1/2 inch tall EAN-13:

\begin{verbatim}
Data:    977147396801
Options: includetext height=0.5
Encoder: ean13
\end{verbatim}

\includegraphics{images/optheight.eps}

\hypertarget{width}{%
\subsubsection{width}\label{width}}

Stretch the symbol to precisely this width, in inches.

\hypertarget{notes-1}{%
\paragraph{Notes}\label{notes-1}}

\begin{itemize}
\tightlist
\item
  This parameter literally stretches the symbol and text to the desired
  width which will may distort the human readable text.
\item
  For information about resizing symbols read the article on
  \protect\hyperlink{resizing-symbols}{resizing symbols}.
\end{itemize}

\hypertarget{example-30}{%
\paragraph{Example}\label{example-30}}

A 2 inch wide Code 93 symbol:

\begin{verbatim}
Data:    TERRY
Options: width=2
Encoder: code93
\end{verbatim}

\includegraphics{images/optwidth.eps}

\hypertarget{element-properties}{%
\subsection{Element Properties}\label{element-properties}}

These are options that allow you to compensate for print gain or print
loss resulting from the properties of the printing process or substrate.

\hypertarget{inkspread}{%
\subsubsection{inkspread}\label{inkspread}}

For linear barcodes, the amount by which to reduce the bar widths to
compensate for inkspread, in points.

For matrix barcodes, the amount by which the reduce the width and height
of dark modules to compensate for inkspread, in points.

\hypertarget{notes-2}{%
\paragraph{Notes}\label{notes-2}}

\begin{itemize}
\tightlist
\item
  Negative values will increase the bar width.
\end{itemize}

\hypertarget{example-31}{%
\paragraph{Example}\label{example-31}}

Code 39 for a printer with very runny ink:

\begin{verbatim}
Data:    TEZ
Options: inkspread=0.6
Encoder: code39
\end{verbatim}

\includegraphics{images/optinkspread.eps}

\hypertarget{inkspreadh}{%
\subsubsection{inkspreadh}\label{inkspreadh}}

For matrix barcodes, the amount by which the reduce the width of dark
modules to compensate for inkspread, in points.

Note: inkspreadh is most useful for stacked-linear type barcodes such as
PDF417 and Codablock F.

\hypertarget{inkspreadv}{%
\subsubsection{inkspreadv}\label{inkspreadv}}

For matrix barcodes, the amount by which the reduce the height of dark
modules to compensate for inkspread, in points.

\hypertarget{dotty}{%
\subsubsection{dotty}\label{dotty}}

For matrix barcodes, render the modules as dots rather than squares. The
dot radius can be adjusted using the \textbf{inkspread} option.

\hypertarget{text-properties}{%
\subsection{Text Properties}\label{text-properties}}

\hypertarget{includetext}{%
\subsubsection{includetext}\label{includetext}}

Show human readable text for data in symbol.

\hypertarget{notes-3}{%
\paragraph{Notes}\label{notes-3}}

\begin{itemize}
\tightlist
\item
  If a character in the data does not have a printable representation
  then it is not displayed
\end{itemize}

\hypertarget{example-32}{%
\paragraph{Example}\label{example-32}}

Display the text encoded in this Code 39 symbol:

\begin{verbatim}
Data:    SEE ME
Options: includetext
Encoder: code39
\end{verbatim}

\includegraphics{images/optincludetext.eps}

\hypertarget{textfont}{%
\subsubsection{textfont}\label{textfont}}

The PostScript font name for text.

\hypertarget{notes-4}{%
\paragraph{Notes}\label{notes-4}}

\begin{itemize}
\tightlist
\item
  The font name must be the literal name of a PostScript font that is
  available to the system.
\item
  This option should be used in combination with the
  \textbf{includetext} option.
\item
  The PostScript font name for a font file can be determined with the
  following command from the fontconfig library:

  \begin{itemize}
  \tightlist
  \item
    \texttt{fc-scan\ -\/-format\ "\%\{postscriptname\}\textbackslash{}n"\ /usr/share/fonts/truetype/somefont.ttf}
  \end{itemize}
\end{itemize}

\hypertarget{example-33}{%
\paragraph{Example}\label{example-33}}

Customise the human readable text of this USPS POSTNET symbol

\begin{verbatim}
Data:    64501
Options: includetext textfont=Times-Roman textsize=9
Encoder: postnet
\end{verbatim}

\includegraphics{images/opttextfont.eps}

\hypertarget{textsize}{%
\subsubsection{textsize}\label{textsize}}

The font size of the text in points.

\hypertarget{note}{%
\paragraph{Note}\label{note}}

\begin{itemize}
\tightlist
\item
  This option should be used in combination with the
  \textbf{includetext} option.
\end{itemize}

\hypertarget{example-34}{%
\paragraph{Example}\label{example-34}}

Customise the human readable text of this USPS POSTNET symbol

\begin{verbatim}
Data:    64501
Options: includetext textfont=Times-Roman textsize=9
Encoder: postnet
\end{verbatim}

\includegraphics{images/opttextfont.eps}

\hypertarget{textgaps}{%
\subsubsection{textgaps}\label{textgaps}}

The inter-character spacing of the text.

\hypertarget{note-1}{%
\paragraph{Note}\label{note-1}}

\begin{itemize}
\tightlist
\item
  This option should be used in combination with the \textbf{textxalign}
  option.
\end{itemize}

\hypertarget{alttext}{%
\subsubsection{alttext}\label{alttext}}

Specify text to display other than what is provided in the data field.

\hypertarget{note-2}{%
\paragraph{Note}\label{note-2}}

\begin{itemize}
\tightlist
\item
  This option should be used in combination with the includetext option.
\end{itemize}

\hypertarget{text-positioning}{%
\subsection{Text Positioning}\label{text-positioning}}

\hypertarget{textxalign}{%
\subsubsection{textxalign}\label{textxalign}}

The \textbf{textxalign} option is used to specify where to horizontally
position the text.

\begin{itemize}
\tightlist
\item
  \texttt{textxalign=offleft}
\item
  \texttt{textxalign=left}
\item
  \texttt{textxalign=center}
\item
  \texttt{textxalign=right}
\item
  \texttt{textxalign=offright}
\item
  \texttt{textxalign=justify}
\end{itemize}

\hypertarget{textyalign}{%
\subsubsection{textyalign}\label{textyalign}}

The \textbf{textyalign} option is used to specify where to vertically
position the text.

\begin{itemize}
\tightlist
\item
  \texttt{textyalign=below}
\item
  \texttt{textyalign=center}
\item
  \texttt{textyalign=above}
\end{itemize}

\hypertarget{textxoffset}{%
\subsubsection{textxoffset}\label{textxoffset}}

The horizontal position of the text in points relative to the default
position.

\hypertarget{textyoffset}{%
\subsubsection{textyoffset}\label{textyoffset}}

The vertical position of the text in points relative to the default
position.

\hypertarget{notes-5}{%
\paragraph{Notes}\label{notes-5}}

\begin{itemize}
\tightlist
\item
  By default (in the absence of \textbf{textxalign} or
  \textbf{textyalign}), each character of text is placed immediately
  below the corresponding modules where this is possible.
\item
  Where there isn\textquotesingle t such a direct relationship then the
  default is to position the text centrally beneath the symbol.
\end{itemize}

\hypertarget{border-properties}{%
\subsection{Border Properties}\label{border-properties}}

\hypertarget{showborder}{%
\subsubsection{showborder}\label{showborder}}

Display a border around the symbol.

\hypertarget{example-35}{%
\paragraph{Example}\label{example-35}}

Display a customised border around this Codabar symbol:

\begin{verbatim}
Data:    A3+14-63B
Options: includetext showborder borderwidth=2 borderbottom=8
Encoder: rationalizedCodabar
\end{verbatim}

\includegraphics{images/optborder.eps}

\hypertarget{borderwidth}{%
\subsubsection{borderwidth}\label{borderwidth}}

Width of the border, in points.

\hypertarget{borderleft}{%
\subsubsection{borderleft}\label{borderleft}}

Gap between the left edge of the border and the symbol, in points.

\hypertarget{borderright}{%
\subsubsection{borderright}\label{borderright}}

Gap between the right edge of the border and the symbol, in points.

\hypertarget{bordertop}{%
\subsubsection{bordertop}\label{bordertop}}

Gap between the top edge of the border and the symbol, in points.

\hypertarget{borderbottom}{%
\subsubsection{borderbottom}\label{borderbottom}}

Gap between the bottom edge of the border and the symbol, in points.

\hypertarget{symbol-colors}{%
\subsection{Symbol Colors}\label{symbol-colors}}

\hypertarget{barcolor}{%
\subsubsection{barcolor}\label{barcolor}}

Color of the bars or dark modules, either as a hex \texttt{RRGGBB} value
or a hex \texttt{CCMMYYKK} value.

\hypertarget{backgroundcolor}{%
\subsubsection{backgroundcolor}\label{backgroundcolor}}

Color of the light background or light modules, either as a hex
\texttt{RRGGBB} value or a hex \texttt{CCMMYYKK} value.

\hypertarget{bordercolor}{%
\subsubsection{bordercolor}\label{bordercolor}}

Color of the border, either as a hex \texttt{RRGGBB} value or a hex
\texttt{CCMMYYKK} value.

\hypertarget{textcolor}{%
\subsubsection{textcolor}\label{textcolor}}

Color of the respective component, either as a hex \texttt{RRGGBB} value
or a hex \texttt{CCMMYYKK} value.

\hypertarget{examples-52}{%
\paragraph{Examples}\label{examples-52}}

Colorized Code 11 symbol using the RGB colour space:

\begin{verbatim}
Data:    0123456789
Options: includetext barcolor=AA0000 textcolor=008888 backgroundcolor=FFFF60 showborder bordercolor=0000FF textyoffset=-8
Encoder: code11
\end{verbatim}

\includegraphics{images/optcolorrgb.eps}

Colorized Code 11 symbol using the CMYK colour space:

\begin{verbatim}
Data:    0123456789
Options: includetext barcolor=AA000000 textcolor=00888844 backgroundcolor=CCCC6000 showborder bordercolor=00FF0088
Encoder: code11
\end{verbatim}

\includegraphics{images/optcolorcmyk.eps}

\hypertarget{ean-upc-add-ons}{%
\subsection{EAN-UPC Add Ons}\label{ean-upc-add-ons}}

\hypertarget{addontextfont}{%
\subsubsection{addontextfont}\label{addontextfont}}

The font name of the add on text.

\hypertarget{notes-6}{%
\paragraph{Notes}\label{notes-6}}

\begin{itemize}
\tightlist
\item
  The font name must be the literal name of a PostScript available font.
\item
  These options must be used in combination with the
  \textbf{includetext} option.
\end{itemize}

\hypertarget{addontextsize}{%
\subsubsection{addontextsize}\label{addontextsize}}

The size of the add on text, in points

\hypertarget{notes-7}{%
\paragraph{Notes}\label{notes-7}}

\begin{itemize}
\tightlist
\item
  These options must be used in combination with the
  \textbf{includetext} option.
\end{itemize}

\hypertarget{example-36}{%
\paragraph{Example}\label{example-36}}

Customise the human readable text of this USPS POSTNET symbol

\begin{verbatim}
Data:    64501
Options: includetext textfont=Times-Roman textsize=9
Encoder: postnet
\end{verbatim}

\includegraphics{images/opttextfont.eps}

\hypertarget{addontextxoffset}{%
\subsubsection{addontextxoffset}\label{addontextxoffset}}

Overrides the default horizontal positioning of the add on text.

\hypertarget{addontextyoffset}{%
\subsubsection{addontextyoffset}\label{addontextyoffset}}

Overrides the default vertical positioning of the add on text.

\hypertarget{ean-upc-guards}{%
\subsection{EAN-UPC Guards}\label{ean-upc-guards}}

\hypertarget{guardwhitespace}{%
\subsubsection{guardwhitespace}\label{guardwhitespace}}

Display whitespace guards.

\hypertarget{example-37}{%
\paragraph{Example}\label{example-37}}

ISBN-13 with tiny whitespace guard:

\begin{verbatim}
Data:    978-1-873671-00
Options: includetext guardwhitespace guardwidth=3 guardheight=3
Encoder: isbn
\end{verbatim}

\includegraphics{images/optguardsize.eps}

\hypertarget{guardwidth}{%
\subsubsection{guardwidth}\label{guardwidth}}

Width of the whitespace guards, in points.

\hypertarget{guardheight}{%
\subsubsection{guardheight}\label{guardheight}}

Height of the whitespace guards, in points.

\hypertarget{guardleftpos}{%
\subsubsection{guardleftpos}\label{guardleftpos}}

Override the default horizontal position of the left whitespace guard.

\hypertarget{guardrightpos}{%
\subsubsection{guardrightpos}\label{guardrightpos}}

Override the default horizontal position of the right whitespace guard.

\hypertarget{guardleftypos}{%
\subsubsection{guardleftypos}\label{guardleftypos}}

Override the default vertical position of the left whitespace guard.

\hypertarget{guardrightypos}{%
\subsubsection{guardrightypos}\label{guardrightypos}}

Override the default vertical position of the right whitespace guard.

\hypertarget{example-38}{%
\paragraph{Example}\label{example-38}}

EAN-8 with very mangled whitespace guards:

\begin{verbatim}
Data:    01335583
Options: includetext guardwhitespace guardleftpos=14 guardrightpos=7 guardleftypos=15 guardrightypos=4
Encoder: ean8
\end{verbatim}

\includegraphics{images/optguardpos.eps}
